% La philosophie antique — Philosophie 1re Tech — Cours signé OnBuch+
% Programme officiel MINESEC. Compiler avec : tectonic N16-philosophie-antique.tex
\newcommand{\DOCMATIERE}{Philosophie}
\newcommand{\DOCNIVEAU}{1re Tech}
\newcommand{\DOCMODULE}{Philosophie – classe de Première technique}
\newcommand{\DOCLECON}{N16}
\newcommand{\DOCTITRE}{La philosophie antique}
% OnBuch+ — préambule « cours signés » ENRICHI (compilé avec tectonic / XeLaTeX)
% Système visuel « Pop » d'OnBuch+ : fond crème (jamais blanc), bordures encre
% épaisses, ombres « dures » décalées, titres Archivo Black en capitales.
% Version MATHÉMATIQUES : graphes (pgfplots/tikz), boîtes pédagogiques
% (définition, propriété, méthode, attention, le savais-tu, exercice résolu,
% exercice, TP, bilan), page de garde et signature OnBuch+ ancrée en bas.
%
% Macros attendues dans main.tex AVANT \input{preamble} :
%   \DOCMATIERE  \DOCNIVEAU  \DOCMODULE  \DOCLECON (numéro)  \DOCTITRE
\documentclass[11pt]{article}

\usepackage{fontspec}
\usepackage[french]{babel}
\usepackage[a4paper,margin=2cm,top=3.4cm,bottom=2.8cm,headheight=1.4cm,footskip=1.3cm]{geometry}
\usepackage{amsmath,amssymb,amsfonts}
\usepackage{mathtools} % \xmapsto, \coloneqq... souvent utilisés par les modèles
\usepackage{cancel} % simplifications barrées (bilans de forces)
% \abs{z} (module, valeur absolue) et \norm{u} (norme), utilisés par les modèles
\providecommand{\abs}{}\let\abs\relax\DeclarePairedDelimiter{\abs}{\lvert}{\rvert}
\providecommand{\norm}{}\let\norm\relax\DeclarePairedDelimiter{\norm}{\lVert}{\rVert}
\usepackage[table]{xcolor} % [table] : fournit aussi \rowcolors (alternance de lignes)
\usepackage{pagecolor}
\usepackage{tikz}
\usetikzlibrary{arrows.meta,positioning,calc,shapes.geometric,shapes.misc,shapes.arrows,decorations.pathmorphing,decorations.pathreplacing,decorations.markings,patterns,shadows,mindmap,trees,fit,backgrounds,matrix,intersections,angles,quotes}
\usepackage{pgfplots}
\usepackage[version=4]{mhchem} % équations nucléaires \ce{^{235}_{92}U}
\usepackage[european,siunitx]{circuitikz} % schémas électriques
\pgfplotsset{compat=1.18}
\usepgfplotslibrary{fillbetween,groupplots} % aires entre courbes (calcul intégral)
\usepackage{enumitem}
\usepackage{fancyhdr}
% [most] charge toutes les bibliothèques tcolorbox (listings, minted, raster,
% documentation...) et gonfle inutilement la mémoire TeX sur un document long
% (~300 boîtes) : on ne charge que ce qui est réellement utilisé.
\usepackage[skins,breakable]{tcolorbox}
\usepackage{graphicx}
\usepackage{colortbl}
\usepackage{booktabs}
\usepackage{tabularx}
\usepackage{multirow}
\usepackage{diagbox} % case d'en-tête diagonale des tableaux à double entrée
% Colonnes extensibles centrée / gauche / droite (C, L, R), souvent utilisées par les modèles
\newcolumntype{C}{>{\centering\arraybackslash}X}
\newcolumntype{L}{>{\raggedright\arraybackslash}X}
\newcolumntype{R}{>{\raggedleft\arraybackslash}X}
\usepackage{array}
\usepackage{multicol}
\usepackage{siunitx}
% parse-numbers=false : accepte aussi \SI{2,5 \times 10^{-3}}{...}
\sisetup{locale=FR,output-decimal-marker={,},group-minimum-digits=4,parse-numbers=false}
\DeclareSIUnit\ohm{\ensuremath{\symup{\Omega}}} % U+2126 absent de la police
\DeclareSIUnit\division{div} % réglages d'oscilloscope (ms/div, V/div)
\DeclareSIUnit\tr{tr} % tours (vitesse de rotation, ex. \SI{1500}{\tr\per\minute})
\DeclareSIUnit\molar{M} % concentration molaire (ex. \SI{3}{\molar}, \SI{1}{\micro\molar})
\DeclareSIUnit\an{an} % année (le modèle confond parfois avec \year, un compteur TeX) — ex. \SI{5}{\kilo\meter\per\an}
\DeclareSIUnit\FCFA{FCFA} % franc CFA (ex. \SI{500}{\FCFA\per\kilo\gram})
\DeclareSIUnit\degreeBrix{\textdegree Brix} % teneur en sucre d'un sirop/jus (réfractométrie)
\DeclareSIUnit\month{mois} % le modèle utilise parfois \month, un compteur TeX, comme unité de durée
\DeclareSIUnit\microliter{\micro\liter} % le modèle écrit parfois \microliter en un seul token
\DeclareSIUnit\million{millions} % dénombrement (ex. \SI{15}{\million\per\mL} spermatozoïdes)
\usepackage{qrcode}
% \degree : commande fréquemment utilisée par les modèles pour un angle ou une
% température en dehors de \SI{}{} (non fournie par les paquets chargés).
\providecommand{\degree}{^{\circ}}
% \qed (fin de démonstration), utilisé par les modèles sans amsthm
\providecommand{\qed}{\hfill$\square$}
% \barre : double barre des valeurs interdites dans un tableau de variations (tkz-tab)
\providecommand{\barre}{\ensuremath{\|}}
% environnement proof (amsthm non chargé) utilisé par les modèles
\ifdefined\proof\else\newenvironment{proof}[1][Démonstration]{\par\noindent\textit{#1.}\ }{\qed\par}\fi
\usepackage{lastpage}
\usepackage{hyperref}
\hypersetup{hidelinks}

% ============================================================
% Palette « Pop » OnBuch+ — mêmes hex que la classe P (plus_ui.dart)
% ============================================================
\definecolor{poporange}{HTML}{F59321}
\definecolor{popdark}{HTML}{E07A0C}
\definecolor{popcream}{HTML}{FFF6E8}
\definecolor{popink}{HTML}{1C1714}
\definecolor{poppurple}{HTML}{7A5AE0}
\definecolor{popgreen}{HTML}{2E9E6A}
\definecolor{poppink}{HTML}{E0457B}
\definecolor{popblue}{HTML}{2F7DE1}
\definecolor{popgold}{HTML}{C98A12}
\definecolor{popmuted}{HTML}{8A7F76}
\definecolor{popline}{HTML}{EADFD2}
\definecolor{popgray}{HTML}{8A7F76}
\colorlet{popgrey}{popgray}
% Alias tolérants (le modèle nomme parfois les couleurs différemment)
\colorlet{popwhite}{white}
\colorlet{popblack}{popink}
\colorlet{popred}{poppink}
\colorlet{popyellow}{popgold}
% Teintes claires (fonds de tableaux/encadrés) — le modèle nomme parfois
% les couleurs avec un suffixe L (ex. popblueL) sans les avoir définies.
\colorlet{poporangeL}{poporange!20}
\colorlet{popdarkL}{popdark!20}
\colorlet{poppurpleL}{poppurple!20}
\colorlet{popgreenL}{popgreen!20}
\colorlet{poppinkL}{poppink!20}
\colorlet{popblueL}{popblue!20}
\colorlet{popgoldL}{popgold!20}
\colorlet{popredL}{popred!20}
\colorlet{popgrayL}{popgray!20}
% Teintes claires pour fonds de tableaux / zones
\colorlet{popblueL}{popblue!10!white}
\colorlet{poporangeL}{poporange!14!white}
\colorlet{popgreenL}{popgreen!12!white}
\colorlet{poppurpleL}{poppurple!12!white}
\colorlet{poppinkL}{poppink!12!white}
\colorlet{popgoldL}{popgold!14!white}

\pagecolor{popcream}

% ============================================================
% Typographie
% ============================================================
\defaultfontfeatures{Ligatures=TeX}
\setmainfont{PlusJakartaSans-Medium.ttf}[Path=../../../fonts/,BoldFont=PlusJakartaSans-Bold.ttf]
% Maths en sans-serif (Fira Math) avec chiffres et lettres droites de Plus
% Jakarta, pour que les formules se fondent dans le texte.
\usepackage[math-style=ISO,bold-style=ISO]{unicode-math}
\setmathfont{FiraMath-Regular.otf}[Path=../../../fonts/]
\setmathfont{PlusJakartaSans-Medium.ttf}[Path=../../../fonts/,range={up/num,up/{Latin,latin},\mathup}]
\usepackage{icomma} % virgule décimale sans espace en mode math
% Caractères mathématiques Unicode absents des polices de texte (Plus Jakarta,
% Archivo Black) : rendus par leur équivalent mathématique, en texte comme en
% formule — sinon ils s'affichent en carré vide (ex. « Arithmétique dans ℤ »).
\usepackage{newunicodechar}
\newunicodechar{ℝ}{\ensuremath{\mathbb{R}}}
\newunicodechar{ℤ}{\ensuremath{\mathbb{Z}}}
\newunicodechar{ℂ}{\ensuremath{\mathbb{C}}}
\newunicodechar{ℕ}{\ensuremath{\mathbb{N}}}
\newunicodechar{ℚ}{\ensuremath{\mathbb{Q}}}
\newunicodechar{𝔻}{\ensuremath{\mathbb{D}}}
\newunicodechar{α}{\ensuremath{\alpha}}
\newunicodechar{β}{\ensuremath{\beta}}
\newunicodechar{γ}{\ensuremath{\gamma}}
\newunicodechar{δ}{\ensuremath{\delta}}
\newunicodechar{ε}{\ensuremath{\varepsilon}}
\newunicodechar{θ}{\ensuremath{\theta}}
\newunicodechar{λ}{\ensuremath{\lambda}}
\newunicodechar{μ}{\ensuremath{\mu}}
\newunicodechar{σ}{\ensuremath{\sigma}}
\newunicodechar{τ}{\ensuremath{\tau}}
\newunicodechar{φ}{\ensuremath{\varphi}}
\newunicodechar{ψ}{\ensuremath{\psi}}
\newunicodechar{ω}{\ensuremath{\omega}}
\newunicodechar{Σ}{\ensuremath{\Sigma}}
% majuscules produites par \MakeUppercase dans les titres (α-aminés -> Α)
\newunicodechar{Α}{\ensuremath{\alpha}}
\newunicodechar{Β}{\ensuremath{\beta}}
\newunicodechar{Γ}{\ensuremath{\Gamma}}
\newunicodechar{Δ}{\ensuremath{\Delta}}
\newunicodechar{Ε}{\ensuremath{\varepsilon}}
\newunicodechar{Θ}{\ensuremath{\Theta}}
\newunicodechar{Λ}{\ensuremath{\Lambda}}
\newunicodechar{Μ}{\ensuremath{\mu}}
\newunicodechar{Τ}{\ensuremath{\tau}}
\newunicodechar{Φ}{\ensuremath{\Phi}}
\newunicodechar{Ψ}{\ensuremath{\Psi}}
\newunicodechar{Ω}{\ensuremath{\Omega}}
\newunicodechar{↦}{\ensuremath{\mapsto}}
\newunicodechar{∈}{\ensuremath{\in}}
\newunicodechar{∉}{\ensuremath{\notin}}
\newunicodechar{∧}{\ensuremath{\wedge}}
\newunicodechar{⇔}{\ensuremath{\Leftrightarrow}}
\newunicodechar{⟺}{\ensuremath{\Longleftrightarrow}}
\newunicodechar{⇒}{\ensuremath{\Rightarrow}}
\newunicodechar{⟹}{\ensuremath{\Longrightarrow}}
\newunicodechar{∘}{\ensuremath{\circ}}
\newunicodechar{∪}{\ensuremath{\cup}}
\newunicodechar{∩}{\ensuremath{\cap}}
\newunicodechar{≡}{\ensuremath{\equiv}}
\newunicodechar{⊂}{\ensuremath{\subset}}
\newunicodechar{⊥}{\ensuremath{\perp}}
\newunicodechar{∃}{\ensuremath{\exists}}
\newunicodechar{∀}{\ensuremath{\forall}}
\newunicodechar{∛}{\ensuremath{\sqrt[3]{\,}}}
\newunicodechar{∜}{\ensuremath{\sqrt[4]{\,}}}
\newunicodechar{⃗}{\ensuremath{{}^{\to}}}
\newunicodechar{ⁿ}{\ifmmode^{n}\else\textsuperscript{\ensuremath{n}}\fi}
\newunicodechar{ˣ}{\ifmmode^{x}\else\textsuperscript{\ensuremath{x}}\fi}
\newunicodechar{ᵇ}{\ifmmode^{b}\else\textsuperscript{\ensuremath{b}}\fi}
\newunicodechar{ʳ}{\ifmmode^{r}\else\textsuperscript{\ensuremath{r}}\fi}
\newunicodechar{ᵘ}{\ifmmode^{u}\else\textsuperscript{\ensuremath{u}}\fi}
\newunicodechar{ʸ}{\ifmmode^{y}\else\textsuperscript{\ensuremath{y}}\fi}
\newunicodechar{ᵃ}{\ifmmode^{a}\else\textsuperscript{\ensuremath{a}}\fi}
\newunicodechar{ᵉ}{\ifmmode^{e}\else\textsuperscript{\ensuremath{e}}\fi}
\newunicodechar{ᵏ}{\ifmmode^{k}\else\textsuperscript{\ensuremath{k}}\fi}
\newunicodechar{ᵖ}{\ifmmode^{p}\else\textsuperscript{\ensuremath{p}}\fi}
\newunicodechar{ᴺ}{\ifmmode^{N}\else\textsuperscript{\ensuremath{N}}\fi}
\newunicodechar{ᶜ}{\ifmmode^{c}\else\textsuperscript{\ensuremath{c}}\fi}
\newunicodechar{ʰ}{\ifmmode^{h}\else\textsuperscript{\ensuremath{h}}\fi}
\newunicodechar{ᵠ}{\ifmmode^{q}\else\textsuperscript{\ensuremath{q}}\fi}
\newunicodechar{ₙ}{\ifmmode_{n}\else\textsubscript{\ensuremath{n}}\fi}
\newunicodechar{ₐ}{\ifmmode_{a}\else\textsubscript{\ensuremath{a}}\fi}
\newunicodechar{ᵢ}{\ifmmode_{i}\else\textsubscript{\ensuremath{i}}\fi}
\newunicodechar{ᵣ}{\ifmmode_{r}\else\textsubscript{\ensuremath{r}}\fi}
\newunicodechar{ₖ}{\ifmmode_{k}\else\textsubscript{\ensuremath{k}}\fi}
\newunicodechar{ᵦ}{\ifmmode_{\beta}\else\textsubscript{\ensuremath{\beta}}\fi}
\newunicodechar{ₓ}{\ifmmode_{x}\else\textsubscript{\ensuremath{x}}\fi}
\newfontfamily\popdisplay{ArchivoBlack-Regular.ttf}[Path=../../../fonts/]
\newfontfamily\poplogo{SpaceGrotesk-Bold.ttf}[Path=../../../fonts/]
\newfontfamily\popsemi{PlusJakartaSans-SemiBold.ttf}[Path=../../../fonts/]
\newfontfamily\popxbold{PlusJakartaSans-ExtraBold.ttf}[Path=../../../fonts/]
\newfontfamily\popxboldls{PlusJakartaSans-ExtraBold.ttf}[Path=../../../fonts/,LetterSpace=3.0]

\setlist[enumerate]{leftmargin=1.7em,itemsep=5pt,topsep=4pt}
\setlist[itemize]{leftmargin=1.7em,itemsep=4pt,topsep=4pt,label={\color{poporange}\textbullet}}
\setlength{\parindent}{0pt}
\setlength{\parskip}{5pt}
\renewcommand{\arraystretch}{1.25}

% pgfplots : style Pop par défaut
\pgfplotsset{
  every axis/.append style={axis line style={popink,line width=1pt},tick style={popink},
    grid style={popline,line width=0.5pt},label style={font=\small},tick label style={font=\footnotesize}},
  popcurve/.style={poporange,line width=1.8pt,no markers,smooth},
}
% Même style disponible dans un \draw TikZ simple (hors pgfplots)
\tikzset{popcurve/.style={poporange,line width=1.8pt}}
\tikzset{popgrid/.style={popline,very thin}}
\colorlet{popcurve}{poporange} % le nom du style est parfois utilisé comme couleur

% ============================================================
% Pill / squiggle
% ============================================================
\newcommand{\poppill}[3]{%
  \tikz[baseline=-0.55ex]\node[draw=popink,line width=1.2pt,rounded corners=2.6pt,%
    fill=#2,inner xsep=6.5pt,inner ysep=3pt]{%
    \popxboldls\fontsize{7}{7}\selectfont\color{#3}\MakeUppercase{#1}};%
}
\newcommand{\popsquiggle}[1]{%
  \tikz[baseline=-0.4ex]{
    \draw[#1,line width=2.6pt,line cap=round]
      plot [smooth] coordinates {(0,0.09) (0.34,0.01) (0.68,0.21) (1.02,0.01) (1.36,0.10)};
  }%
}
% Mot-clé surligné (marqueur orange sous le texte)
\newcommand{\cle}[1]{{\popxbold\color{popdark}#1}}

% ============================================================
% En-tête / pied
% ============================================================
\pagestyle{fancy}
\fancyhf{}
\renewcommand{\headrulewidth}{0pt}
\renewcommand{\footrulewidth}{0pt}
\lhead{\raisebox{-0.15ex}{\poplogo\fontsize{13}{13}\selectfont\color{popink}OnBuch\color{poporange}+}}
\rhead{\poppill{\DOCMATIERE\ \textbullet\ \DOCNIVEAU\ \textbullet\ Leçon \DOCLECON}{white}{popink}}
\lfoot{\popxbold\fontsize{7.6}{7.6}\selectfont\color{popmuted}\MakeUppercase{Cours signé OnBuch+}\ \textbullet\ {\popsemi\DOCTITRE}}
\rfoot{\tikz[baseline=-0.6ex]\node[rounded corners=8.5pt,fill=popink,minimum height=17pt,minimum width=17pt,inner xsep=5pt,inner sep=0]{\popxbold\fontsize{8}{8}\selectfont\color{popcream}\thepage\,/\,\pageref*{LastPage}};}

% ============================================================
% Titres
% ============================================================
\newcommand{\chaptitre}[2]{%
  \begin{tcolorbox}[enhanced,colback=poporange,colframe=popink,boxrule=2.6pt,arc=11pt,
      drop shadow=popink,left=15pt,right=15pt,top=12pt,bottom=11pt,before skip=3pt,after skip=18pt]
    \poppill{#2}{popink}{popcream}\par\vspace{9pt}
    {\raggedright\hyphenpenalty=10000\exhyphenpenalty=10000\popdisplay\fontsize{22}{24}\selectfont\color{popcream}\MakeUppercase{#1}\par}
  \end{tcolorbox}
}
% Section numérotée : pastille orange avec le numéro + titre Archivo
\newcounter{popsec}
\newcommand{\coursec}[1]{%
  \stepcounter{popsec}\setcounter{popsub}{0}%
  \par\vspace{16pt}\noindent
  \begin{minipage}{\linewidth}
  \tikz[baseline=-0.7ex]\node[draw=popink,line width=1.6pt,fill=poporange,rounded corners=4pt,minimum size=22pt,inner sep=0]{\popdisplay\fontsize{12}{12}\selectfont\color{popcream}\thepopsec};%
  \hspace{8pt}{\popdisplay\fontsize{15}{17}\selectfont\color{popink}\MakeUppercase{#1}}\par
  \vspace{4pt}\noindent{\color{poporange}\rule{36pt}{3.6pt}}
  \end{minipage}\par\nopagebreak\vspace{8pt}
}
\newcounter{popsub}
\newcommand{\courssub}[1]{%
  \stepcounter{popsub}%
  \par\vspace{9pt}\noindent{\popxbold\fontsize{11.5}{13}\selectfont\color{popink}{\color{poporange}\thepopsec.\thepopsub}\hspace{6pt}#1}\par\nopagebreak\vspace{4pt}
}

% ============================================================
% Boîtes pédagogiques — même recette : fond blanc, bordure épaisse colorée,
% ombre dure encre, badge plein. Titre optionnel [#1].
% ============================================================
\tcbset{popbase/.style={breakable,enhanced,colback=white,boxrule=2.3pt,arc=9pt,
  shadow={3pt}{-3pt}{0pt}{popink},left=14pt,right=14pt,top=11pt,bottom=11pt,before skip=11pt,after skip=15pt,
  fontupper=\popsemi\fontsize{10.6}{14.5}\selectfont\color{popink},
  fontlower=\popsemi\fontsize{10.6}{14.5}\selectfont\color{popink}}}
% Coche dessinée (le glyphe ✓ n'existe pas dans les polices du document)
\newcommand{\popcheck}[1]{\tikz[baseline=-0.5ex]\draw[#1,line width=1.6pt,line cap=round,line join=round](-0.13,0.0)--(-0.04,-0.09)--(0.13,0.1);}
\providecommand{\coche}{\popcheck{popgreen}}
\providecommand{\resultat}[1]{\par\smallskip\noindent\fcolorbox{popgreen}{popgreenL}{\parbox{\dimexpr\linewidth-2\fboxsep-2\fboxrule\relax}{\textbf{Résultat.} #1}}\par\smallskip}
\newcommand{\popboxhead}[3]{\poppill{#1}{#2}{white}\ifx\relax#3\relax\else\hspace{6pt}{\popxbold\fontsize{10.6}{12}\selectfont\color{popink}#3}\fi\par\vspace{7pt}}

\newtcolorbox{aretenir}[1][]{popbase,colframe=popgreen,before upper={\popboxhead{À retenir}{popgreen}{#1}}}
\newtcolorbox{exemplebox}[1][]{popbase,colframe=poporange,before upper={\popboxhead{Exemple}{poporange}{#1}}}
\newtcolorbox{definition}[1][]{popbase,colframe=popblue,before upper={\popboxhead{Définition}{popblue}{#1}}}
\newtcolorbox{propriete}[1][]{popbase,colframe=poppurple,before upper={\popboxhead{Propriété}{poppurple}{#1}}}
\newtcolorbox{methode}[1][]{popbase,colframe=popgold,before upper={\popboxhead{Méthode}{popgold}{#1}}}
\newtcolorbox{attention}[1][]{popbase,colframe=poppink,before upper={\popboxhead{Attention}{poppink}{#1}}}
\newtcolorbox{experience}[1][]{popbase,colframe=popblue,colback=popblueL,before upper={\popboxhead{Expérience}{popblue}{#1}}}
% « Le savais-tu ? » : carte violette pleine (contexte, culture, Cameroun)
\newtcolorbox{savaistu}[1][]{popbase,colback=poppurple,colframe=popink,
  fontupper=\popsemi\fontsize{10.4}{14.2}\selectfont\color{white},
  before upper={\poppill{Le savais-tu\,?}{white}{poppurple}\ifx\relax#1\relax\else\hspace{6pt}{\popxbold\color{white}#1}\fi\par\vspace{7pt}}}
% Exercice résolu : énoncé en haut, \tcblower puis la solution sur fond crème
\newcounter{popexo}\newcounter{popexr}
\newtcolorbox{exoresolu}[1][]{popbase,colframe=poporange,colbacklower=popcream,
  segmentation style={popink,line width=1.2pt,dashed},
  before upper={\stepcounter{popexr}\popboxhead{Exercice résolu \thepopexr}{poporange}{#1}},
  before lower={\poppill{Solution}{popink}{popcream}\par\vspace{6pt}}}
% Exercice (non résolu en place) — la correction est dans la partie Corrigés
\newtcolorbox{exercice}[1][]{popbase,colframe=popink,
  before upper={\stepcounter{popexo}\popboxhead{Exercice \thepopexo}{popink}{#1}}}
% Corrigé
\newtcolorbox{corrige}[1][]{popbase,colframe=popgreen,colback=popgreenL,
  before upper={\popboxhead{Corrigé}{popgreen}{#1}}}
% Objectifs / prérequis / bilan
\newtcolorbox{objectifs}{popbase,colframe=popink,colback=poporangeL,
  before upper={\popboxhead{Objectifs de la leçon}{popink}{}}}
\newtcolorbox{prerequis}{popbase,colframe=popink,colback=popblueL,
  before upper={\popboxhead{Prérequis}{popblue}{}}}
\newtcolorbox{bilan}{popbase,colframe=popink,colback=white,boxrule=2.8pt,
  before upper={\popboxhead{Fiche bilan}{poporange}{L'essentiel à connaître pour l'examen}}}
% Figure encadrée avec légende
\newtcolorbox{popfigure}[1][]{enhanced,breakable=false,colback=white,colframe=popline,boxrule=1.4pt,arc=8pt,
  left=10pt,right=10pt,top=10pt,bottom=8pt,before skip=10pt,after skip=12pt,halign=center,
  lower separated=false,halign lower=center,
  fontlower=\popsemi\fontsize{9}{11.5}\selectfont\color{popmuted},
  #1}
\newcommand{\legende}[1]{\par\vspace{6pt}{\popsemi\fontsize{9}{11.5}\selectfont\color{popmuted}#1}}

% ============================================================
% Signature OnBuch+
% ============================================================
\newcommand{\poplogoinline}[1]{{\poplogo\fontsize{#1}{#1}\selectfont\color{popink}OnBuch\color{poporange}+}}
\newcommand{\signatureonbuch}{%
  \begin{tcolorbox}[enhanced,colback=popink,colframe=popink,boxrule=2.2pt,arc=10pt,
      drop shadow=poporange,left=14pt,right=14pt,top=10pt,bottom=10pt,before skip=0pt,after skip=0pt]
    \tikz[baseline=-0.7ex]\node[circle,fill=popgreen,draw=popcream,line width=1.3pt,minimum size=22pt,inner sep=0]{\popcheck{white}};%
    \hspace{9pt}\begin{minipage}[c]{0.62\linewidth}
      {\popxbold\fontsize{12}{13}\selectfont\color{popcream}Signé\ \textbullet\ {\poplogo OnBuch\color{poporange}+}}\par\vspace{2pt}
      {\raggedright\popsemi\fontsize{8}{10.5}\selectfont\color{popline}\MakeUppercase{Équipe pédagogique OnBuch+}\\\MakeUppercase{Conforme au programme officiel MINESEC}\par}
    \end{minipage}\hfill
    {\poplogo\fontsize{22}{22}\selectfont\color{popcream}OnBuch\color{poporange}+}
  \end{tcolorbox}%
}

% ============================================================
% Page de garde
% #1 titre · #2 matière · #3 niveau · #4 module · #5 n° leçon · #6 durée ·
% #7 description / objectifs (texte ou itemize)
% ============================================================
\newcommand{\pagegarde}[7]{%
  \thispagestyle{empty}
  \newgeometry{a4paper,margin=1.7cm,top=1.6cm,bottom=1.4cm}%
  \begin{tikzpicture}[remember picture,overlay]
    \fill[poporange!18] ([xshift=-3.2cm,yshift=-3.6cm]current page.north east) circle (4.6cm);
    \fill[poppurple!14] ([xshift=2.4cm,yshift=7.6cm]current page.south west) circle (3.8cm);
  \end{tikzpicture}%
  {\poplogo\fontsize{30}{30}\selectfont\color{popink}OnBuch\color{poporange}+}\hfill
  \raisebox{0.6ex}{\poppill{Cours signé \textbullet\ Premium}{popink}{popcream}}\par
  \vspace{1pt}\popsquiggle{poppurple}\par
  \vspace{8pt}
  \begin{tcolorbox}[enhanced,colback=poporange,colframe=popink,boxrule=2.8pt,arc=13pt,
      drop shadow=popink,left=18pt,right=18pt,top=12pt,bottom=13pt,after skip=10pt]
    \poppill{#2 \textbullet\ #3}{popink}{popcream}\hspace{5pt}\poppill{#4}{white}{popink}\par\vspace{8pt}
    {\popdisplay\fontsize{14}{14}\selectfont\color{popink}LEÇON #5}\par\vspace{4pt}
    {\raggedright\hyphenpenalty=10000\exhyphenpenalty=10000\popdisplay\fontsize{25}{27}\selectfont\color{popcream}\MakeUppercase{#1}\par}
    \vspace{8pt}
    {\popxbold\fontsize{9.5}{9.5}\selectfont\color{popink}\MakeUppercase{Durée conseillée : #6}}
  \end{tcolorbox}
  \begin{tcolorbox}[enhanced,colback=white,colframe=popink,boxrule=2.2pt,arc=10pt,
      drop shadow=popink,left=16pt,right=16pt,top=10pt,bottom=10pt,after skip=8pt]
    \poppill{Dans cette leçon}{poppurple}{white}\par\vspace{5pt}
    \popsemi\fontsize{9.3}{12.6}\selectfont\color{popink}#7
  \end{tcolorbox}
  \noindent\begin{minipage}[t]{0.487\linewidth}
    \begin{tcolorbox}[enhanced,colback=popgreen,colframe=popink,boxrule=2.2pt,arc=10pt,
        drop shadow=popink,left=11pt,right=11pt,top=10pt,bottom=10pt,height=3.1cm,valign=center]
      \tcbox[colback=white,colframe=popink,boxrule=1.3pt,arc=5pt,boxsep=0pt,nobeforeafter,
        left=3pt,right=3pt,top=3pt,bottom=3pt,valign=center]{\qrcode[height=1.9cm]{https://play.google.com/store/apps/details?id=com.luvvix.onbuch&hl=fr}}%
      \hspace{8pt}\begin{minipage}[c]{0.5\linewidth}
        {\popdisplay\fontsize{11}{12.5}\selectfont\color{white}\MakeUppercase{Télécharge OnBuch}}\par\vspace{4pt}
        {\raggedright\popsemi\fontsize{8.4}{11}\selectfont\color{white}Gratuit \textbullet\ Google Play\\Tuteur IA, annales, concours, résultats\par}
      \end{minipage}
    \end{tcolorbox}
  \end{minipage}\hfill
  \begin{minipage}[t]{0.487\linewidth}
    \begin{tcolorbox}[enhanced,colback=poppurple,colframe=popink,boxrule=2.2pt,arc=10pt,
        drop shadow=popink,left=11pt,right=11pt,top=10pt,bottom=10pt,height=3.1cm,valign=center]
      \tikz[baseline=-0.7ex]\node[circle,fill=white,draw=popink,line width=1.3pt,minimum size=1.6cm,inner sep=0]{\popdisplay\fontsize{24}{24}\selectfont\color{poppurple}+};%
      \hspace{8pt}\begin{minipage}[c]{0.6\linewidth}
        {\popdisplay\fontsize{11}{12.5}\selectfont\color{white}\MakeUppercase{Passe à OnBuch+}}\par\vspace{4pt}
        {\raggedright\popsemi\fontsize{8.4}{11}\selectfont\color{white}Tous les cours signés, Tuteur IA illimité, fiches PDF — onglet OnBuch+ de l'app\par}
      \end{minipage}
    \end{tcolorbox}
  \end{minipage}
  \vspace{8pt}
  \popcontenu
  \vfill
  \signatureonbuch
  \restoregeometry
  \newpage
}

% Rangée « Ce cours contient » (page de garde)
\newcommand{\poptile}[3]{% #1 couleur #2 gros texte #3 légende
  \begin{tcolorbox}[enhanced,colback=#1,colframe=popink,boxrule=2pt,arc=9pt,shadow={2.5pt}{-2.5pt}{0pt}{popink},
      left=6pt,right=6pt,top=9pt,bottom=9pt,halign=center,valign=center,height=1.75cm,width=\linewidth,nobeforeafter]
    {\popdisplay\fontsize{12.5}{13}\selectfont\color{white}\MakeUppercase{#2}}\par\vspace{3pt}
    {\centering\popsemi\fontsize{7.8}{9.5}\selectfont\color{white}#3\par}
  \end{tcolorbox}}
\newcommand{\popcontenu}{%
  \par\noindent{\popxboldls\fontsize{7.5}{7.5}\selectfont\color{popmuted}CE COURS SIGNÉ CONTIENT}\par\vspace{7pt}
  \noindent\begin{minipage}[t]{0.235\linewidth}\poptile{popblue}{Cours}{complet, illustré, pas à pas}\end{minipage}\hfill
  \begin{minipage}[t]{0.235\linewidth}\poptile{poporange}{Méthodes}{et exercices résolus}\end{minipage}\hfill
  \begin{minipage}[t]{0.235\linewidth}\poptile{poppink}{Exercices}{type examen + corrigés}\end{minipage}\hfill
  \begin{minipage}[t]{0.235\linewidth}\poptile{popgreen}{Fiche bilan}{l'essentiel à retenir}\end{minipage}\par}

% Sommaire
\newtcolorbox{sommaire}{enhanced,colback=white,colframe=popink,boxrule=2.2pt,arc=10pt,
  drop shadow=popink,left=18pt,right=18pt,top=13pt,bottom=13pt,before skip=6pt,after skip=20pt,
  before upper={\poppill{Sommaire}{poppurple}{white}\par\vspace{9pt}\popsemi\fontsize{10.6}{14.5}\selectfont\color{popink}}}

% Signature de fin de cours — poussée tout en bas de la dernière page
\newcommand{\findecours}{%
  \par\vfill\nopagebreak
  \begin{center}\popsquiggle{poporange}\end{center}\vspace{4pt}
  \signatureonbuch
}
\AtBeginDocument{\def\llbracket{\mathopen{[\mkern-3.5mu[}}\def\rrbracket{\mathclose{]\mkern-3.5mu]}}} % intervalles d'entiers (glyphe ⟦ absent de la police)
% Glyphes absents de Fira Math : redéfinis à partir de glyphes disponibles
\AtBeginDocument{%
  \def\setminus{\mathbin{\backslash}}\let\smallsetminus\setminus
  \def\vdots{\mathord{\vbox{\baselineskip3.2pt\lineskiplimit0pt\kern4pt\hbox{.}\hbox{.}\hbox{.}}}}%
  \def\ddots{\mathinner{\mkern1mu\raise6.5pt\hbox{.}\mkern2mu\raise3.6pt\hbox{.}\mkern2mu\raise0.7pt\hbox{.}\mkern1mu}}%
  \def\triangle{\mathord{\scalebox{0.9}{$\symup{\Delta}$}}}%
  \def\nabla{\mathord{\raisebox{\depth}{\scalebox{1}[-1]{$\symup{\Delta}$}}}}%
  \def\checkmark{\popcheck{popdark}}%
  \def\star{\mathbin{\ast}}\def\bigstar{\mathord{\ast}}%
  \def\diamond{\mathbin{\scalebox{0.6}{$\Diamond$}}}%
  \def\bigcup{\mathop{\mathchoice{\raisebox{-0.3em}{\scalebox{1.7}{$\cup$}}}{\scalebox{1.25}{$\cup$}}{\cup}{\cup}}\displaylimits}%
  \def\bigcap{\mathop{\mathchoice{\raisebox{-0.3em}{\scalebox{1.7}{$\cap$}}}{\scalebox{1.25}{$\cap$}}{\cap}{\cap}}\displaylimits}%
  \def\lesseqqgtr{\mathrel{\lessgtr}}\def\bigoplus{\mathop{\oplus}\displaylimits}%
}
\providecommand{\ssi}{si et seulement si} % abréviation fréquente
\AtBeginDocument{\providecommand{\wideparen}{\overparen}} % arc de cercle

%% FIGFIX-BEGIN v5 — correctifs globaux des figures (docs/onbuch-generation/tools/figfix)
\usepackage{adjustbox}\usepackage{etoolbox}\tcbuselibrary{raster}
% toute figure TikZ/pgfplots plus large que la ligne est réduite à la largeur disponible
\BeforeBeginEnvironment{tikzpicture}{\begin{adjustbox}{max width=\linewidth}}
\AfterEndEnvironment{tikzpicture}{\end{adjustbox}}
% légendes pgfplots lisibles par-dessus les courbes
\pgfplotsset{every axis/.append style={legend style={fill=white,fill opacity=0.92,text opacity=1,draw=black!15,inner sep=3pt}}}
% étiquettes d'arêtes (syntaxe edge["..."]) avec fond blanc
\tikzset{every edge quotes/.append style={fill=white,inner sep=1.2pt}}
%% FIGFIX-END
\begin{document}

\pagegarde{La philosophie antique}{Philosophie}{1re Tech}{Philosophie – classe de Première technique}{N16}{3 séances (fiche de progression harmonisée)}{Cette leçon explore les fondements de la pensée occidentale à travers la philosophie pharaonique, l'enquête des Présocratiques sur l'Un et le Multiple, et l'anthropologie philosophique de Socrate et Platon. Elle développe l'esprit critique en articulant concepts rigoureux et situations concrètes camerounaises pour préparer aux épreuves du Baccalauréat technique.
\vspace{4pt}
\begin{multicols}{2}\small
\begin{itemize}
\item À la fin de cette leçon, je sais expliquer le concept de Maât et son rôle structurant dans la philosophie pharaonique en le comparant aux valeurs traditionnelles camerounaises (ex: le Nguon chez les Bamoun).
\item Je sais distinguer la rupture épistémologique opérée par les Présocratiques (passage du mythe au logos) et caractériser l'opposition fondamentale entre le Devenir (Héraclite) et l'Être (Parménide).
\item Je sais analyser la méthode socratique (ironie, maïeutique, définition) comme fondement de l'éthique rationnelle et de l'exigence de cohérence.
\item Je sais exposer la théorie des Idées (Formes) de Platon, l'allégorie de la Caverne et la cité juste comme réponses aux problèmes posés par ses prédécesseurs.
\item Je sais appliquer les concepts d'Unité, d'Identité et de Devenir à des situations concrètes de la vie courante (ex: gestion d'une coopérative, identité numérique, justice sociale).
\item Je sais rédiger une dissertation philosophique structurée (introduction, thèse, antithèse, synthèse, conclusion) et justifier chaque étape de mon raisonnement.
\end{itemize}\end{multicols}}

\begin{sommaire}
\begin{multicols}{2}
\begin{enumerate}
\item 1. Les sources africaines : La philosophie pharaonique et l'ordre du monde (Maât)
\item 2. La naissance de la philosophie grecque : Les Présocratiques, Héraclite et Parménide
\item 3. Le tournant anthropologique : Socrate, l'éthique par le dialogue
\item 4. La systématisation platonicienne : Être, Idées, Politique et Éducation
\item 5. Synthèse : Problématique de l'Unité et du Multiple – Application aux situations camerounaises
\item 6. Atelier méthodologique : Rédiger et justifier une démarche philosophique (Dissertation \& Explication)
\item Activité d'intégration
\item Méthodes et exercices résolus
\item Exercices
\item Corrigés des exercices
\item Fiche bilan
\end{enumerate}
\end{multicols}
\end{sommaire}

\begin{objectifs}
\begin{itemize}
\item À la fin de cette leçon, je sais expliquer le concept de Maât et son rôle structurant dans la philosophie pharaonique en le comparant aux valeurs traditionnelles camerounaises (ex: le Nguon chez les Bamoun).
\item Je sais distinguer la rupture épistémologique opérée par les Présocratiques (passage du mythe au logos) et caractériser l'opposition fondamentale entre le Devenir (Héraclite) et l'Être (Parménide).
\item Je sais analyser la méthode socratique (ironie, maïeutique, définition) comme fondement de l'éthique rationnelle et de l'exigence de cohérence.
\item Je sais exposer la théorie des Idées (Formes) de Platon, l'allégorie de la Caverne et la cité juste comme réponses aux problèmes posés par ses prédécesseurs.
\item Je sais appliquer les concepts d'Unité, d'Identité et de Devenir à des situations concrètes de la vie courante (ex: gestion d'une coopérative, identité numérique, justice sociale).
\item Je sais rédiger une dissertation philosophique structurée (introduction, thèse, antithèse, synthèse, conclusion) et justifier chaque étape de mon raisonnement.
\item Je sais expliquer un texte philosophique antique en identifiant la thèse, les arguments, les concepts clés et en en évaluant la portée actuelle.
\item Je sais construire un schéma conceptuel (carte heuristique, tableau comparatif) synthétisant l'évolution de la problématique de l'Être de l'Antiquité à nos jours.
\end{itemize}
\end{objectifs}

\begin{prerequis}
\begin{itemize}
\item Notion de mythe et de logos (programme de 4ème / Seconde).
\item Maîtrise de la structure d'une dissertation et de l'explication de texte (début d'année 1re).
\item Repères chronologiques et géographiques de base : Égypte antique, Grèce archaïque et classique (Méditerranée).
\item Vocabulaire philosophique élémentaire : concept, définition, argument, thèse, antithèse, synthèse.
\item Connaissance des grandes chefferies traditionnelles camerounaises (organisation politique, justice coutumière).
\end{itemize}
\end{prerequis}

\newpage

\coursec{Les sources africaines : La philosophie pharaonique et l'ordre du monde (Maât)}

\courssub{Situation d'accroche et problématique}

Au palais du Sultan des Bamoun à Foumban, la transmission du pouvoir et la résolution des conflits reposent sur une sagesse ancestrale codifiée, tout comme à l'époque pharaonique où \cle{Maât} garantissait l'ordre cosmique et social. Aujourd'hui, face aux défis de la modernité (réseaux sociaux, intelligence artificielle, crises écologiques), les jeunes Camerounais cherchent des repères stables : tout change-t-il sans cesse comme le fleuve Wouri, ou existe-t-il des vérités immuables comme le mont Cameroun ? Cette tension entre le flux et la permanence, entre l'opinion (\cle{doxa}) et la science (\cle{epistémè}), est au cœur de la naissance de la philosophie en Grèce antique, un héritage qui structure encore notre pensée critique. Mais avant Athènes, il y eut \cle{Kemet} (l'Égypte noire). La philosophie ne naît pas qu'en Grèce ; l'Afrique antique a produit une pensée rationnelle de l'ordre, de la justice et de la vérité. Comment le concept de Maât structure-t-il une philosophie de l'action, du politique et de l'éthique ? En quoi cette sagesse institutionnalisée éclaire-t-elle nos défis contemporains ?

\courssub{Kemet : Géographie, histoire et émergence de la pensée structurée}

L'Égypte antique, que ses habitants nommaient \textbf{Kemet} (« la terre noire », en référence au limon fertile du Nil), n'est pas une simple périphérie du monde méditerranéen. C'est le berceau africain d'une pensée \cle{archétypale} qui a structuré la civilisation pendant plus de trois millénaires (vers -3150 à -30).

\begin{popfigure}
\begin{tikzpicture}[scale=0.9, every node/.style={font=\small}]
% Timeline axis
\draw[popdark, thick, ->] (0,0) -- (15,0) node[right] {Chronologie (av. J.-C.)};
\foreach \x/\label in {0/-3150, 2.5/-2700, 5/-2100, 7.5/-1500, 10/-1000, 12.5/-500, 15/0}
    \draw (\x,0.1) -- (\x,-0.1) node[below=2pt] {\label};
% Egypt periods
\draw[popblue, line width=8pt, opacity=0.6] (0,1) -- (15,1) node[right, opacity=1, popblue] {Égypte (Kemet)};
\node[popblue, align=center] at (1.25,1.6) {Période thinite\\ (-3150)};
\node[popblue, align=center] at (3.75,1.6) {Ancien Empire\\ Pyramides};
\node[popblue, align=center] at (6.25,1.6) {Moyen Empire\\ Textes Sarcophages};
\node[popblue, align=center] at (8.75,1.6) {Nouvel Empire\\ Livre des Morts};
\node[popblue, align=center] at (11.25,1.6) {Basse Époque};
% Greece
\draw[poporange, line width=8pt, opacity=0.6] (10,-1) -- (15,-1) node[right, opacity=1, poporange] {Grèce};
\node[poporange, align=center] at (12.5,-0.4) {Présocratiques\\ (-600)};
\node[poporange, align=center] at (14,-0.4) {Socrate/Platon\\ (-400)};
% China
\draw[popgreen, line width=8pt, opacity=0.6] (5,-2) -- (15,-2) node[right, opacity=1, popgreen] {Chine};
\node[popgreen, align=center] at (10,-1.4) {Confucius/Lao Tseu\\ (-500)};
% India
\draw[poppurple, line width=8pt, opacity=0.6] (5,-3) -- (15,-3) node[right, opacity=1, poppurple] {Inde};
\node[poppurple, align=center] at (10,-2.4) {Upanishads/Bouddha\\ (-500)};
% Cameroon Kingdoms (schematic)
\draw[poppink, line width=8pt, opacity=0.6] (10,-4) -- (15,-4) node[right, opacity=1, poppink] {Cameroun (Royaumes)};
\node[poppink, align=center] at (12.5,-3.4) {Bamoun, Bamiléké,\\ Kotoko (XIIe-XIXe s.)};
% Connection arrow
\draw[popdark, dashed, ->, thick] (7.5,1) -- (11,-1) node[midway, above, sloped, font=\tiny] {Influence possible (Platon en Égypte)};
\end{tikzpicture}
\legende{Frise chronologique comparative : L'antériorité de la pensée structurée en Kemet (Égypte) par rapport aux autres foyers civilisationnels. Les royaumes camerounais (Bamoun, Bamiléké) structurent leur pensée politique et éthique bien plus tard, mais sur des schémas comparables d'ordre cosmique.}
\end{popfigure}

La pensée égyptienne n'est pas une « pré-philosophie » mythique. Elle répond aux critères de la rationalité : elle \textbf{conceptualise} (Maât, Ka, Ba, Akh), elle \textbf{argumente} (textes de sagesse, débats juridiques, hymnes théologiques) et elle \textbf{systématise} (cosmogonie d'Héliopolis, théologie memphite). Le Nil, par sa crue prévisible, offre le modèle d'un cosmos ordonné, cyclique mais stable : l'intelligence humaine doit s'accorder à cet ordre.

\begin{definition}[Kemet et la rationalité africaine antique]
\cle{Kemet} (« Terre Noire ») désigne l'Égypte pharaonique comme espace civilisationnel africain. Sa pensée se caractérise par :
\begin{itemize}
    \item L'\textbf{immanence} : le divin n'est pas hors du monde, il \emph{est} l'ordre du monde (Maât).
    \item La \textbf{fonctionnalité} : la vérité est ce qui « fait tenir » le monde, la société, l'individu.
    \item L'\textbf{institutionnalisation} : la philosophie n'est pas l'affaire de sages isolés (comme en Grèce archaïque), mais une charge d'État confiée aux scribes, vizirs, prêtres et au Pharaon lui-même.
\end{itemize}
\end{definition}

\courssub{Maât : Concept central, fondement ontologique, éthique et politique}

\cle{Maât} (souvent traduit par Vérité, Justice, Ordre, Harmonie, Droiture) est le concept \textbf{arché} (premier principe) de la philosophie pharaonique. Ce n'est pas une déesse parmi d'autres (bien qu'elle soit personnifiée), c'est la \textbf{structure même du réel}.

\begin{popfigure}
\begin{tikzpicture}[scale=0.85, every node/.style={font=\small, align=center}]
% Central concept
\node[circle, draw=popblue, thick, fill=popblueL, minimum width=2.5cm, align=center] (maat) at (0,0){\textbf{MAÂT}\\\emph{Ordre / Justice / Vérité}};
% Cosmic level
\node[ellipse, draw=popgold, thick, fill=popgoldL, align=center] (cosmic) at (-4,3){Niveau Cosmique\\ (Régularité des astres,\\ Crue du Nil, Saisons)};
\draw[popgold, thick, ->] (cosmic) -- (maat) node[midway, above, sloped] {manifestation};
% Divine level
\node[ellipse, draw=poppurple, thick, fill=poppurpleL, align=center] (divine) at (4,3){Niveau Divin\\ (Atum/Rê vit \emph{en} Maât\\ Les dieux sont « Seigneurs de Maât »)};
\draw[poppurple, thick, ->] (divine) -- (maat) node[midway, above, sloped] {fondement};
% Royal level
\node[rectangle, draw=poporange, thick, fill=poporangeL, rounded corners, align=center] (royal) at (0,1.8){PHARAON\\ « Seigneur de Maât »\\ Garant de l'ordre};
\draw[poporange, thick, ->] (royal) -- (maat) node[midway, left] {incarnation};
% Social level
\node[ellipse, draw=popgreen, thick, fill=popgreenL, align=center] (social) at (-4,-2.5){Niveau Social\\ (Justice, Droit, Éthique,\\ « Faire Maât » = bien agir)};
\draw[popgreen, thick, ->] (social) -- (maat) node[midway, below, sloped] {application};
% Individual level
\node[ellipse, draw=poppink, thick, fill=poppinkL, align=center] (indiv) at (4,-2.5){Niveau Individuel\\ (Cœur léger, 42 confessions,\\ Accès à l'éternité)};
\draw[poppink, thick, ->] (indiv) -- (maat) node[midway, below, sloped] {réalisation};
% Isfet (Chaos)
\node[ellipse, draw=popdark, thick, fill=popmuted, align=center] (isfet) at (0,-4.5){\textbf{ISFET} (Chaos, Mensonge, Désordre)\\ \emph{Antagoniste radical de Maât}};
\draw[popdark, thick, <->, dashed] (maat) -- (isfet) node[midway, right] {Lutte permanente};
\end{tikzpicture}
\legende{Carte mentale de la cosmologie égyptienne : Maât comme principe immanent unifiant le cosmique, le divin, le royal, le social et l'individuel. Isfet (le chaos) est l'ennemi constant à repousser par l'action juste.}
\end{popfigure}

\begin{propriete}[Structure ontologique de Maât]
Maât possède trois dimensions inséparables :
\begin{enumerate}
    \item \textbf{Ontologique (Être) :} Maât \emph{est} l'intelligibilité du monde. Sans Maât, retour au \cle{Nun} (eau primordiale, chaos indifférencié). « Je suis Maât » dit le dieu créateur Atum.
    \item \textbf{Éthique (Agir) :} \emph{Faire Maât} (\textit{iri Maât}), c'est agir conformément à la vérité et la justice. C'est une éthique de la \textbf{conduite} (orthopraxie) avant d'être une éthique de la croyance (orthodoxie).
    \item \textbf{Politique (Gouverner) :} Le Pharaon est le \textbf{seul} garant de Maât sur terre. Sa légitimité n'est pas dynastique seulement, mais \emph{fonctionnelle} : s'il ne fait pas Maât, il perd le mandat du ciel (perte du \cle{Ka} royal).
\end{enumerate}
\end{propriete}

\begin{exemplebox}[Maât dans l'administration : Le Vizir « Premier Juge »]
Le vizir (chef du gouvernement) porte le titre de « \emph{Propriétaire de Maât} ». Dans le \textbf{Papyrus de la Tombe de Rekhmire} (Nouvel Empire, vers -1450), ses instructions sont explicites :
\begin{quote}
\textit{« Regarde, le vizir est le support de Maât... Ne fais pas de différence entre le fils d'un grand et le fils d'un petit... Car la force de l'un est Maât, la force de l'autre est Maât. »}
\end{quote}
C'est une \textbf{théorie de l'égalité devant la loi} et de l'impartialité judiciaire formulée 1000 ans avant Solon à Athènes. Au Cameroun aujourd'hui, l'article 1er de la Constitution (« L'État assure l'égalité devant la loi... ») en est l'héritier structural, bien que la pratique en soit souvent écartée.
\end{exemplebox}

\courssub{Les sources documentaires : Pyramides, Sarcophages, Livre des Morts}

La philosophie pharaonique ne nous est pas parvenue par des traités abstraits (comme Aristote), mais par des \textbf{textes à visée opératoire} : ils visent à \emph{faire advenir} la vérité dans le réel (magie rituelle, justice, royauté).

\begin{center}
\begin{tabularx}{\linewidth}{|l|X|X|X|}
\hline
\rowcolor{popblueL}
\textbf{Corpus} & \textbf{Période} & \textbf{Support / Auteurs} & \textbf{Contenu philosophique majeur} \\
\hline
\textbf{Textes des Pyramides} & Ancien Empire (v. -2400 à -2200) & Murs des pyramides royales (Ounas, Téti, Pépi II) & Cosmogonie héliopolitaine (Atum, Ennéade). Royauté divine. Ascension du roi vers l'éternité. Première formulation écrite de Maât comme ordre cosmique. \\
\hline
\textbf{Textes des Sarcophages} & Moyen Empire (v. -2050 à -1750) & Sarcophages en bois (nobles, fonctionnaires) & \textbf{Démocratisation de l'au-delà} : l'éternité n'est plus réservée au roi. Développement de l'éthique individuelle (Jugement post-mortem). Apparition du concept de \cle{Ba} (personnalité) distinct du \cle{Ka} (force vitale). \\
\hline
\textbf{Livre des Morts} (\textit{Per em herou}) & Nouvel Empire à Basse Époque (v. -1550 à -30) & Papyrus funéraires illustrés (ex: Papyrus d'Ani) & \textbf{Systématisation éthique} : Chapitre 125 (Pesée du cœur, 42 confessions négatives). Guide pratique pour « sortir au jour » (vivre éternellement). Maât comme code moral explicite. \\
\hline
\end{tabularx}
\end{center}

\begin{aretenir}[Évolution de la pensée : Du royal à l'universel]
Le passage des \textbf{Textes des Pyramides} (royaux) aux \textbf{Textes des Sarcophages} puis au \textbf{Livre des Morts} (accessibles à tout lettré) marque une \textbf{démocratisation de la philosophie} : la responsabilité éthique devient individuelle. Tout homme devient « pharaon de son propre destin » s'il réalise Maât. C'est une anticipation de l'autonomie morale kantienne.
\end{aretenir}

\courssub{Anthropologie pharaonique : Corps, Ka, Ba, Akh et le Jugement}

L'homme égyptien n'est pas une substance unique. Il est une \textbf{unité complexe} dont la cohérence dépend de Maât.

\begin{definition}[Les constituants de l'être humain]
\begin{itemize}
    \item \textbf{\cle{Khat} (Corps) :} Support matériel, périssable. Doit être momifié pour servir d'ancre aux autres composantes.
    \item \textbf{\cle{Ka} (Double / Force vitale) :} Énergie vitale reçue du créateur à la naissance. Identique pour tous. Se nourrit d'offrandes. C'est le \emph{droit à la vie} et à la subsistance.
    \item \textbf{\cle{Ba} (Personnalité / « Âme » mobile) :} Ce qui fait l'unicité (caractère, mémoire, réputation). Représenté oiseau à tête humaine. Il quitte le corps à la mort mais doit y revenir. C'est la \emph{responsabilité morale}.
    \item \textbf{\cle{Akh} (Esprit glorifié / Transfiguré) :} Résultat de la fusion réussie Ka + Ba \emph{après} le jugement. L'être accompli, éternel, « efficace » dans l'au-delà.
    \item \textbf{\cle{Ren} (Nom) \& \cle{Shout} (Ombre) :} Garants de l'identité sociale et de la présence au monde.
\end{itemize}
\end{definition}

\begin{popfigure}
\begin{tikzpicture}[scale=0.9, every node/.style={font=\small, align=center}]
% Scene: Weighing of the Heart
% Balance
\draw[popdark, very thick] (0,2.5) -- (0,0.5);
\draw[popdark, thick] (-1.5,2.5) -- (1.5,2.5); % Beam
\draw[popdark, thick] (-1.5,2.5) -- (-1.3,2.2) -- (-1.5,1.9); % Left pivot decor
\draw[popdark, thick] (1.5,2.5) -- (1.3,2.2) -- (1.5,1.9); % Right pivot decor
% Left Pan: Heart
\draw[popdark, thick] (-1.5,2.5) -- (-1.5,1.2);
\node[poppink, fill=poppink!10, draw=poppink, thick, rounded corners, minimum width=1.2cm, minimum height=0.8cm, align=center] (heart) at (-1.5,0.8){\textbf{CŒUR}\\\emph{(Ib)}\\\emph{\textit{Siège de la conscience,}}\\\emph{\textit{mémoire, intention}}};
\node[below=2pt of heart, font=\tiny, popdark] {ACTIONS / PENSÉES};
% Right Pan: Feather of Maât
\draw[popdark, thick] (1.5,2.5) -- (1.5,1.2);
\node[popblue, fill=popblueL, draw=popblue, thick, rounded corners, minimum width=1.2cm, minimum height=0.8cm, align=center] (feather) at (1.5,0.8){\textbf{PLUME}\\\emph{DE MAÂT}\\\textit{Vérité / Justice}\\\textit{\textit{Étalon immuable}}};
\node[below=2pt of feather, font=\tiny, popdark] {ORDRE COSMIQUE};
% Anubis (Jackal) adjusting balance
\node[poporange, draw=poporange, thick, fill=poporangeL, rounded corners, anchor=north, align=center] (anubis) at (-2.5,1.5){\textbf{ANUBIS}\\\emph{Guide des morts}\\\emph{Vérifie l'équilibre}};
\draw[poporange, thick] (anubis.south east) -- (-1.5,1.5);
% Thoth (Ibis) recording
\node[poppurple, draw=poppurple, thick, fill=poppurpleL, rounded corners, anchor=north, align=center] (thoth) at (2.5,1.5){\textbf{THOT}\\\emph{Scribe divin}\\\emph{Enregistre le verdict}};
\draw[poppurple, thick] (thoth.south west) -- (1.5,1.5);
% Ammit (Devourer) waiting
\node[popmuted, draw=popdark, thick, fill=popmuted!30, rounded corners, anchor=north, align=center] (ammit) at (0,-1.5){\textbf{AMMIT}\\\emph{\textit{Dévoratrice}}\\\emph{\textit{ (Crocodile/Lion/Hippopotame)}}\\\emph{\textit{Si cœur > Plume : Néant}}};
% Arrow result
\draw[popgreen, very thick, ->] (0,0.2) -- (0,-1.2) node[midway, right] {\textbf{ÉQUILIBRE} = \cle{Akh} (Éternité)};
\draw[poppink, very thick, ->] (0,0.2) -- (-2,-1.2) node[midway, left] {\textbf{CŒUR LOURD} = \cle{Isfet}};
% Deceased observing
\node[popgreen, draw=popgreen, thick, fill=popgreenL, rounded corners, anchor=north, align=center] (deceased) at (-4.5,0){\textbf{LE DÉFUNT}\\\emph{(Ani)}\\ \textit{Récite les 42 confessions}\\\textit{ négatives :}\\\textit{ « Je n'ai pas volé...}\\\textit{ Je n'ai pas menti... »}};
% Maât goddess overseeing
\node[popgold, draw=popgold, thick, fill=popgoldL, rounded corners, anchor=south, align=center] (maatg) at (0,3.5){\textbf{MAÂT}\\\emph{Préside la scène}\\\emph{Garante de l'impartialité}};
\end{tikzpicture}
\legende{Schéma de la Pesée du Cœur (Livre des Morts, Chap. 125). Le cœur (siège de l'intelligence morale) est pesé face à la Plume de Maât (étalon de vérité). Anubis vérifie, Thot note. Si le cœur est « léger » (pur), accès à l'Akh ; s'il est « lourd » (coupable), dévoré par Ammit (néant). C'est une \textbf{juridiction de la conscience} sans avocat ni juge partial.}
\end{popfigure}

\begin{exoresolu}[Analyse éthique des 42 Confessions Négatives (Chapitre 125)]
\textbf{Énoncé :} Le défunt doit déclarer : « Je n'ai pas fait le mal », « Je n'ai pas volé », « Je n'ai pas tué », « Je n'ai pas menti », « Je n'ai pas détourné les mesures du grain », « Je n'ai pas calomnié », « Je n'ai pas fait souffrir », « Je n'ai pas commis d'adultère », « Je n'ai pas pollué l'eau », etc. (Papyrus d'Ani).
\textbf{Analyse philosophique :}
\begin{enumerate}
    \item \textbf{Négativité formelle :} L'éthique se définit par \emph{l'abstention} du mal (Isfet) plutôt que par la prescription du bien. Similaire à la \cle{Règle d'argent} (Confucius : « Ne fais pas à autrui... »).
    \item \textbf{Universalisme concret :} Les fautes couvrent le droit pénal (vol, meurtre), le droit civil (mesures, limites), l'éthique sociale (calomnie, souffrance), l'écologie (eau, air), la vie privée (adultère).
    \item \textbf{Responsabilité sans excuse :} Nul ne peut dire « je ne savais pas ». La loi (Maât) est inscrite dans l'ordre du monde. Le cœur \emph{est} le témoin interne.
    \item \textbf{Actualité camerounaise :} « Je n'ai pas détourné les mesures du grain » $\rightarrow$ Lutte contre la corruption (CONAC), pesée des denrées aux marchés (Mfoundi, Mokolo). « Je n'ai pas pollué l'eau » $\rightarrow$ Gestion du Wouri, pollution plastique à Douala.
\end{enumerate}
\textbf{Conclusion :} Ce code éthique explicite, antérieur de 1500 ans au Décalogue ou aux Tables de la Loi, fonde une \textbf{citoyenneté responsable} : l'individu répond de ses actes devant l'Ordre universel.
\tcblower
\textbf{Solution détaillée :} Voir analyse ci-dessus. L'exercice type au Bac : « Expliquez en quoi la Pesée du Cœur fonde une éthique de la responsabilité personnelle. » $\rightarrow$ Structure : 1. Scène juridictionnelle impartiale. 2. Cœur = conscience morale interne (pas de mensonge possible). 3. Plume = norme universelle (Maât). 4. Verdict binaire (Éternité / Néant) = enjeu radical de l'existence.
\end{exoresolu}

\courssub{Comparaison structurée : Maât (Égypte) vs Nguon / Mvet (Cameroun)}

La philosophie comparée n'est pas du folklore : elle révèle les \textbf{invariants anthropologiques} de la pensée africaine. Utilisons la méthode du tableau à trois colonnes.

\begin{methode}[Comparer deux pensées philosophiques : Tableau à 5 colonnes]
Pour comparer Maât (Égypte antique) et une tradition camerounaise (ex: \cle{Nguon} chez les Bamoun, \cle{Mvet} chez les Fang/Beti) :
\begin{enumerate}
    \item \textbf{Concept clé} : Identifier la notion centrale (Ordre, Justice, Parole, Initiation).
    \item \textbf{Formulation Égypte (Maât)} : Terme, définition, support textuel.
    \item \textbf{Formulation Cameroun (Nguon/Mvet)} : Terme vernaculaire, définition, support oral/rituel.
    \item \textbf{Point commun (Invariant)} : Structure logique identique (ex: immanence, responsabilité, oralité sacralisée).
    \item \textbf{Différence (Contexte)} : Écriture vs Oralité, État centralisé vs Société de chefferies, Universalisme cosmique vs Ancestralisme lignager.
\end{enumerate}
\end{methode}

\begin{center}
\begin{tabularx}{\linewidth}{|l|X|X|X|X|}
\hline
\rowcolor{popblueL}
\textbf{Concept} & \textbf{Égypte : Maât} & \textbf{Cameroun : Nguon (Bamoun) / Mvet (Fang)} & \textbf{Point commun (Universel)} & \textbf{Différence (Spécifique)} \\
\hline
\textbf{Ordre cosmique \& social} & Maât : Loi unique liant ciel, roi, homme, nature. & \textbf{Nguon} : Fête/institution rénovant l'alliance Roi-Peuple-Ancêtres. \textbf{Mvet} : Épopée cosmogonique structurant le temps et l'éthique. & L'ordre n'est pas naturel, il est \textbf{construit, entretenu, célébré} (rite). Le désordre (Isfet / \textit{Nkù}) menace toujours. & Égypte : État bureaucratique, textes fixes. Cameroun : Sociétés à oralité sacrée, adaptation contextuelle. \\
\hline
\textbf{Justice / Vérité} & Pesée du cœur (intériorité). 42 confessions. Loi écrite (Papyrus). & \textbf{Nguon} : Tribunal public (Mbombam), serment sur le \textit{Nguon}, vérité par l'épreuve. \textbf{Mvet} : Parole initiée (\textit{Mvet Oyeng}) comme vérité absolue. & La vérité est \textbf{performative} : la dire juste, c'est la faire advenir. Justice = restauration du lien social rompu. & Égypte : Jugement post-mortem individuel. Cameroun : Jugement social immédiat, communautaire, réconciliateur. \\
\hline
\textbf{Légitimité du pouvoir} & Pharaon = « Seigneur de Maât ». Mandat cosmique révocable (théoriquement). & Sultan / Fon = Garant de l'ordre ancestral. Investiture par les sociétés secrètes (Ngiri, Ngouon) / Notables. & Le pouvoir n'est pas \emph{propriete} mais \emph{fonction} : « Servir l'Ordre ». Redevabilité envers le peuple/ancêtres. & Égypte : Théocratie centralisée, roi dieu. Cameroun : Diarchie (Roi + Sociétés secrètes/Notables), roi premier serviteur. \\
\hline
\textbf{Éducation / Transmission} & Écoles des scribes (Per-Ankh). Copie de textes de sagesse (Ptahhotep, Amenemope). & \textbf{Nguon} : Initiation par étapes (Mbombam, Nguon). \textbf{Mvet} : École des initiés (Bebey), mémorisation épique. & \textbf{Paideia} africaine : Former l'homme complet (corps, cœur, parole, devoir). Maîtrise de la parole (\textit{Maa} / \textit{Zom}). & Égypte : Lettrés (élite scribale). Cameroun : Initiés (souvent lignagers, mais ouverture au mérite). \\
\hline
\end{tabularx}
\end{center}

\begin{savaistu}[Platon en Égypte : L'héritage africain de la République]
La tradition antique (Diogène Laërce, Clément d'Alexandrie) rapporte que \textbf{Platon} a séjourné 13 ans en Égypte (vers -390) auprès des prêtres d'Héliopolis et de Sais.
\begin{itemize}
    \item Il y a découvert la \textbf{division tripartite de la société} (Prêtres/Guerriers/Artisans) $\rightarrow$ Réplique exacte dans \emph{La République} (Gardiens/Auxiliaires/Producteurs).
    \item La \textbf{communauté des biens et des femmes} pour les gardiens $\rightarrow$ Inspiré du mode de vie des prêtres égyptiens (célibat, biens communs).
    \item L'\textbf{éducation gymnastique/musicale} $\rightarrow$ Modèle de l'éducation des scribes (corps/esprit).
    \item Le \textbf{mythe d'Er} (fin de la République, jugement des âmes, choix de vie) $\rightarrow$ Transposition quasi littérale de la \textbf{Pesée du Cœur} et du Livre des Morts.
\end{itemize}
\emph{Conclusion :} La philosophie grecque, souvent présentée comme « miracle grec », a des racines africaines documentées. Le \cle{Probatoire} attend que vous sachiez situer cet apport : « La philosophie ne naît pas qu'en Grèce ; l'Afrique antique a produit une pensée rationnelle de l'ordre. »
\end{savaistu}

\courssub{Application aux situations concrètes camerounaises (Savoir-faire)}

Le programme exige d'\textbf{appliquer les notions à des situations concrètes}. Voici trois cas types pour l'oral ou l'écrit.

\begin{exercice}[Cas pratique 1 : La corruption au guichet administratif]
\textbf{Situation :} Un fonctionnaire à la Mairie de Yaoundé VI exige 5000 FCFA « pour accélérer » un acte de naissance gratuit.
\textbf{Question :} Analysez cette situation à la lumière du concept de Maât et des 42 confessions. En quoi le fonctionnaire rompt-il Maât ? Quel est l'impact sur l'\cle{Akh} social (la confiance) ?
\end{exercice}

\begin{corrige}[Exercice 1]
\begin{enumerate}
    \item \textbf{Rupture de Maât :} Le fonctionnaire commet \textit{Isfet} (désordre). Il viole « Je n'ai pas volé », « Je n'ai pas détourné la mesure », « Je n'ai pas fait souffrir ». Il substitue l'intérêt privé (argent) à l'intérêt général (service public = Maât).
    \item \textbf{Poids du cœur :} Son « cœur » (conscience professionnelle) est alourdi par la cupidité. Il ne peut plus « sortir au jour » (exercer sa fonction avec légitimité).
    \item \textbf{Conséquence collective :} La société (le \textit{Ka} collectif) se délite. La confiance (Maât social) s'effondre. L'État (Pharaon moderne) perd sa légitimité s'il ne sanctionne pas (Ammit).
    \item \textbf{Réparation :} Restitution, sanction disciplinaire, réaffirmation du serment professionnel (nouvelle « confession négative »).
\end{enumerate}
\end{corrige}

\begin{exercice}[Cas pratique 2 : Conflit foncier et justice coutumière (Chefferie Bamiléké)]
\textbf{Situation :} Deux familles se disputent une limite de champ. Le Chef de groupement convoque les parties, fait prêter serment sur le \textit{Têté} (masque sacré / ancêtres), écoute les anciens, rend un jugement « sur la tête » (définitif, sans appel).
\textbf{Question :} Montrez en quoi cette procédure réalise l'équivalent fonctionnel de la Pesée du Cœur (Maât) dans un cadre oral.
\end{exercice}

\begin{corrige}[Exercice 2]
\begin{itemize}
    \item \textbf{Équivalent de la Plume :} Le serment sur les ancêtres (Têté) = Étalon de vérité immuable (Maât). On ne ment pas aux morts.
    \item \textbf{Équivalent du Cœur :} La parole publique des parties. La « confession » est orale, face à la communauté.
    \item \textbf{Équivalent d'Anubis/Thot :} Le Chef + Conseil des notables. Ils « pèsent » la parole (cohérence, témoins, ancienneté).
    \item \textbf{Équivalent de l'Akh :} La réconciliation (\textit{Nkwa}), la paix revenue. Le jugement restaure Maât (lien social), il ne se contente pas de punir.
    \item \textbf{Différence clé :} Pas de jugement post-mortem, mais efficacité \emph{ici et maintenant} (urgence agraire).
\end{itemize}
\end{corrige}

\begin{exercice}[Cas pratique 3 : Intelligence Artificielle et « Vérité » algorithmique]
\textbf{Situation :} Un algorithme de notation sociale (score citoyen) ou de recrutement automatique prend des décisions opaques. Les citoyens ne comprennent pas les critères (boîte noire).
\textbf{Question :} En quoi l'opacité algorithmique est-elle une forme moderne d'\cle{Isfet} (chaos/mensonge) par rapport à l'exigence de Maât (transparence, mesure, justice) ?
\end{exercice}

\begin{corrige}[Exercice 3]
\begin{itemize}
    \item \textbf{Maât = Mesure et Transparence :} « Je n'ai pas faussé la balance » (Confession 125). L'algorithme opaque est une balance truquée.
    \item \textbf{Responsabilité :} En Égypte, le vizir répondait personnellement. Aujourd'hui, qui répond de l'erreur de l'IA ? (Problème de l'\cle{agency} moral).
    \item \textbf{Ordre vs Chaos :} Un système dont on ne connaît pas les règles génère de l'arbitraire (Isfet), pas de l'ordre (Maât).
    \item \textbf{Piste de régulation :} Exiger une « Maât algorithmique » : auditabilité, explicabilité (\textit{explainable AI}), droit à la contestation (équivalent de la défense devant Osiris).
\end{itemize}
\end{corrige}

\courssub{Synthèse de la section : Maât, première philosophie de l'ordre}

\begin{aretenir}[Ce qu'il faut retenir pour le Probatoire / Bac Tech]
\begin{enumerate}
    \item \textbf{Antériorité :} Kemet (Égypte) produit une philosophie structurée (ontologie, éthique, politique) dès -2500, soit 2000 ans avant Socrate.
    \item \textbf{Concept clé :} \cle{Maât} = Ordre cosmique immanent = Justice sociale = Vérité épistémique = Devoir éthique. Un concept \textbf{totalisant}.
    \item \textbf{Anthropologie :} L'homme est un être de responsabilité (Ka/Ba/Cœur). La mort n'est pas une fin mais un \textbf{jugement de vérité} (Pesée du cœur).
    \item \textbf{Éthique concrète :} 42 confessions négatives = Code moral universel (droit, social, écologie, personnel). Préfigure la loi morale kantienne et les Droits de l'Homme.
    \item \textbf{Politique :} Le pouvoir est \emph{fonction} (garant de Maât), non propriété. Redevabilité, impartialité, service.
    \item \textbf{Continuité africaine :} Nguon, Mvet, Ngondo, Duala, Bassa... reproduisent la structure : Ordre sacré, Justice restaurative, Parole performative, Roi/Sultan serviteur.
    \item \textbf{Apport au Bac :} \textbf{La philosophie ne naît pas qu'en Grèce.} L'Afrique antique a produit une pensée rationnelle de l'ordre, documentée par des textes (Pyramides, Sarcophages, Livre des Morts) et vivante dans nos traditions.
\end{enumerate}
\end{aretenir}

\begin{attention}[Pièges à éviter à l'examen]
\begin{itemize}
    \item \textbf{Confondre Religion et Philosophie :} Dire « Les Égyptiens croyaient en des dieux, donc c'est de la religion, pas de la philosophie. » \textbf{Faux.} Maât est un \textbf{concept philosophique} (principe d'intelligibilité et d'action) avant d'être une déesse. La théologie memphite (Pierre de Shabaka, -700) est une \textbf{cosmologie rationnelle} (création par la Parole/Logos = \textit{logos spermatikos}).
    \item \textbf{Essentialiser l'Afrique :} Dire « La pensée africaine est uniquement orale/émotionnelle/communautaire. » \textbf{Faux.} L'Égypte a l'\textbf{écriture} (hiéroglyphes, hiératique, démotique), l'\textbf{abstraction} (Maât, Nun, Atum), l'\textbf{universalisme} (valable pour tout homme), l'\textbf{institution} (écoles, administration, droit).
    \item \textbf{Oublier la comparaison structurée :} Au Bac, si le sujet porte sur « Ordre et Justice », ne citez pas seulement Platon/Aristote. Intégrez \textbf{Maât} et \textbf{Nguon/Mvet} comme \textbf{références majeures} du programme de Première Tech. C'est la marque de l'ancrage camerounais exigé.
\end{itemize}
\end{attention}

\coursec{La naissance de la philosophie grecque : Les Présocratiques, Héraclite et Parménide}

Dans la section précédente, nous avons vu que la pensée pharaonique fondait l'ordre du monde sur \cle{Maât}, principe d'harmonie cosmique, sociale et morale. Cette sagesse africaine cherchait déjà à répondre à une question fondamentale : qu'est-ce qui maintient le monde en ordre ? Au VI\textsuperscript{e} siècle avant J.-C., sur la côte de l'Ionie (Asie Mineure actuelle), des penseurs grecs vont poser cette même question, mais avec une exigence nouvelle : répondre non plus par les récits mythiques, mais par la \cle{raison} seule. C'est cette rupture que l'on appelle la \emph{naissance de la philosophie}. Mais aussitôt née, la philosophie se divise : le monde est-il changement perpétuel (Héraclite) ou immobilité éternelle (Parménide) ? Telle est l'aporie que cette section va explorer.

\courssub{Le contexte : Milet et la naissance de la « physique »}

\textbf{Un creuset de civilisations}\par

Milet est une cité commerçante prospère d'Ionie, au contact des Lydiens, des Égyptiens et des Babyloniens. Les échanges marchands y brassent les savoirs : astronomie babylonienne, géométrie égyptienne, techniques de navigation. Or, dans ce carrefour culturel, aucune autorité religieuse unique n'impose de dogme. Les citoyens débattent dans l'\emph{agora} (place publique). C'est dans ce climat de liberté et de discussion que des hommes comme Thalès, Anaximandre et Anaximène — les \cle{Milesiens} — osent expliquer le monde sans invoquer les dieux.

\begin{definition}[L'archè]
L'\cle{archè} (en grec : \emph{commencement}, \emph{principe}) désigne le principe premier, l'origine constitutive de toutes choses, ce dont tout provient et à quoi tout retourne. Exemples : l'\textbf{Eau} pour Thalès, l'\textbf{Apeiron} (l'illimité) pour Anaximandre, l'\textbf{Air} pour Anaximène, le \textbf{Feu/Logos} pour Héraclite, l'\textbf{Être} pour Parménide.
\end{definition}

\textbf{Du mythe au logos}\par

Ce qui est révolutionnaire à Milet, ce n'est pas seulement la réponse donnée (l'eau, l'air...), c'est la \emph{manière} de chercher. Les Présocratiques pratiquent la \cle{physis} (la « physique », c'est-à-dire l'étude de la nature) : ils veulent une explication \emph{naturelle}, \emph{universelle} et \emph{rationnelle} des phénomènes. Le tonnerre n'est plus la colère de Zeus, mais un phénomène à expliquer. On passe ainsi du \emph{muthos} (le récit mythique) au \cle{logos} (le discours raisonné, l'argument). Cette exigence de justification rationnelle est l'acte de naissance de la philosophie et de la science.

\begin{aretenir}[L'essentiel du contexte]
La philosophie naît à Milet (Ionie) au VI\textsuperscript{e} s. av. J.-C., lorsque des penseurs cherchent l'\cle{archè} (principe premier) de la nature (\emph{physis}) par la raison (\emph{logos}) et non plus par le mythe. La question centrale : \emph{qu'est-ce qui demeure derrière la diversité changeante des choses ?}
\end{aretenir}

\begin{popfigure}
\begin{center}
\begin{tikzpicture}[
  node distance=0.55cm and 1.1cm,
  every node/.style={font=\small},
  mil/.style={rectangle, draw=popblue, fill=popblueL, rounded corners=2pt, minimum height=0.8cm, align=center},
  py/.style={rectangle, draw=popgreen, fill=popgreenL, rounded corners=2pt, minimum height=0.8cm, align=center},
  hp/.style={rectangle, draw=poporange, fill=poporangeL, rounded corners=2pt, minimum height=0.8cm, align=center},
  plu/.style={rectangle, draw=poppurple, fill=poppurpleL, rounded corners=2pt, minimum height=0.8cm, align=center},
  arr/.style={-{Stealth[length=2.5mm]}, thick, popmuted}
]
\node[mil, align=center] (thales){Thalès\\ Eau};
\node[mil, right=of thales, align=center] (anaximandre){Anaximandre\\ Apeiron};
\node[mil, right=of anaximandre, align=center] (anaximene){Anaximène\\ Air};
\node[py, below=1.1cm of anaximandre, align=center] (pyth){Pythagore\\ Nombre};
\node[hp, below left=1.2cm and 0.4cm of pyth, align=center] (hera){Héraclite\\ Feu / Devenir};
\node[hp, below right=1.2cm and 0.4cm of pyth, align=center] (parm){Parménide\\ Être};
\node[plu, below left=1.2cm and -0.2cm of $(hera)!0.5!(parm)$, align=center] (emp){Empédocle\\ 4 éléments};
\node[plu, right=0.6cm of emp, align=center] (anax){Anaxagore\\ Noûs};
\node[plu, right=0.6cm of anax, align=center] (atom){Atomistes\\ Atomes};
\draw[arr] (thales) -- (anaximandre);
\draw[arr] (anaximandre) -- (anaximene);
\draw[arr] (anaximandre) -- (pyth);
\draw[arr] (pyth) -- (hera);
\draw[arr] (pyth) -- (parm);
\draw[arr] (hera.south) |- (emp.north);
\draw[arr] (parm.south) |- (anax.north);
\draw[arr] (hera.south) -- ++(0,-0.35) -| (atom.north);
\draw[arr] (parm.south) -- ++(0,-0.35) -| (atom.north);
\end{tikzpicture}
\end{center}
\legende{Arbre conceptuel des Présocratiques : des Milesiens (Eau, Air, Apeiron) à Pythagore (Nombre), puis au grand partage Héraclite (Devenir) / Parménide (Être), que les pluralistes (Empédocle, Anaxagore, Atomistes) tenteront de dépasser.}
\end{popfigure}

\courssub{Héraclite d'Éphèse : le philosophe du Devenir}

\textbf{« Tout coule » : l'intuition du fleuve}\par

Observe un fleuve, par exemple la Sanaga : l'eau que tu vois passer à cet instant n'est déjà plus la même à l'instant suivant. Pourtant, tu dis toujours « la Sanaga », comme s'il s'agissait d'une seule et même chose. C'est cette intuition qu'Héraclite d'Éphèse (vers 544--480 av. J.-C.) radicalise : \emph{« On ne se baigne jamais deux fois dans le même fleuve »}. Tout est en mouvement, tout devient : c'est la formule célèbre \cle{Panta rhei} (« tout coule »), rapportée par Simplicius. La réalité n'est pas faite de choses fixes, mais de \emph{processus} : le jour devient nuit, la vie devient mort, le jeune devient vieux.

\textbf{L'unité des contraires}\par

Mais Héraclite ne dit pas seulement que tout change. Sa thèse profonde est que les contraires \emph{s'unissent} dans le changement : \emph{« Le chemin en montée et en descente est un et le même »}. La même route est « montante » ou « descendante » selon le sens où on la parcourt : les contraires sont donc deux faces d'une même réalité. De même, la mort des uns est la vie des autres ; le froid se réchauffe, le chaud se refroidit. Le monde est une \emph{tension harmonieuse} d'opposés, comparable à celle de l'arc et de la lyre : c'est la tension des contraires qui produit l'harmonie.

\begin{propriete}[Le Logos héraclitéen]
Le \cle{Logos} chez Héraclite n'est pas un simple discours humain : c'est la \textbf{loi rationnelle universelle} qui gouverne le Devenir ordonné des contraires. Le changement n'est pas un chaos : il obéit à une mesure, à un rythme, à une « harmonie invisible » (plus belle que l'harmonie visible). Les hommes vivent dans le Logos mais la plupart ne le comprennent pas : ils « sont comme endormis ».
\end{propriete}

\textbf{Le feu, symbole du Devenir}\par

Héraclite choisit le \cle{feu} comme archè : non pas une matière au sens des Milesiens, mais le symbole même de la transformation. Le feu détruit et consume, mais il est aussi mesure et lumière : le monde est « un feu éternellement vivant, qui s'allume avec mesure et s'éteint avec mesure ». Le feu figure à la fois le changement permanent et l'ordre (Logos) qui le règle.

\begin{attention}[Piège fréquent]
Ne pas réduire Héraclite au « tout change ». S'il n'y avait que du changement pur, rien ne serait identifiable ni pensable. Héraclite affirme l'\textbf{unité dans le changement} : c'est le Logos, l'harmonie invisible des contraires, qui fait que le devenir est un \emph{ordre} et non un chaos. Oublier le Logos, c'est trahir Héraclite.
\end{attention}

\courssub{Parménide d'Élée : le philosophe de l'Être}

\textbf{Le poème \emph{Sur la Nature}}\par

Parménide d'Élée (vers 515--450 av. J.-C.) expose sa doctrine dans un poème philosophique, \emph{Sur la Nature}, composé de trois parties : un \textbf{proème} (le voyage initiatique du jeune philosophe vers la déesse Vérité), la \textbf{Voie de la Vérité} (ce que la raison établit) et la \textbf{Voie de l'Opinion} (ce que les sens et les mortels croient). C'est la première fois dans l'histoire de la pensée qu'un philosophe procède par \emph{démonstration logique} rigoureuse.

\textbf{« L'Être est, le Non-être n'est pas »}\par

Le raisonnement de Parménide est d'une simplicité redoutable. Penser, c'est toujours penser \emph{quelque chose} : on ne peut penser le néant. Or ce qui peut être pensé \emph{est}. Donc : \emph{l'Être est, et le Non-être n'est pas}. Conséquences déduites par la raison :
\begin{itemize}
\item L'Être est \textbf{ingénéré et impérissable} : s'il naissait, il viendrait du non-être, ce qui est impossible (le non-être n'est pas) ;
\item L'Être est \textbf{Un} : s'il y avait plusieurs êtres, ils seraient séparés par du non-être, ce qui est impossible ;
\item L'Être est \textbf{immuable et immobile} : changer, c'est passer de ce qu'on est à ce qu'on n'est pas, donc supposer le non-être ;
\item L'Être est \textbf{éternel} et \textbf{sphérique} : parfait, homogène, achevé de toutes parts, comme une sphère.
\end{itemize}

\begin{definition}[L'ontologie]
Parménide fonde l'\cle{ontologie}, la science de l'Être en tant qu'Être. Sa méthode : seule la raison (\cle{noûs}) atteint la vérité ; les \textbf{sens trompent}, car ils nous montrent un monde de mouvement, de naissance et de mort qui n'est qu'\emph{opinion} (\emph{doxa}), apparence illusoire.
\end{definition}

\textbf{Le conflit raison / sens}\par

Voici le scandale parménidien : mes sens me montrent le mouvement, mais ma raison démontre que le mouvement est impossible. Parménide tranche : c'est la raison qui dit le vrai, les sens mentent. Le monde du changement n'est qu'apparence. C'est la première grande affirmation de la supériorité de la raison sur l'expérience sensible — geste fondateur de toute la métaphysique occidentale.

\begin{popfigure}
\begin{center}
\begin{tikzpicture}[scale=1]
% Héraclite : flèches circulaires
\begin{scope}[xshift=-3.4cm]
  \draw[thick, poporange, -{Stealth[length=3mm]}] (0:1.15) arc (0:120:1.15);
  \draw[thick, poporange, -{Stealth[length=3mm]}] (120:1.15) arc (120:240:1.15);
  \draw[thick, poporange, -{Stealth[length=3mm]}] (240:1.15) arc (240:360:1.15);
  \node[font=\small\bfseries, poporange] at (0,0) {Devenir};
  \node[font=\footnotesize, align=center, popdark] at (0,-1.75) {Héraclite : tout coule,\\ les contraires s'unissent};
\end{scope}
% VS
\node[font=\large\bfseries, popmuted] at (0,0.1) {vs};
% Parménide : sphère statique
\begin{scope}[xshift=3.4cm]
  \draw[thick, popblue, fill=popblueL] (0,0) circle (1.15);
  \draw[popblue] (0,0) ellipse (1.15 and 0.42);
  \draw[popblue, dashed] (0,0) ellipse (0.42 and 1.15);
  \node[font=\small\bfseries, popblue] at (0,0) {Être};
  \node[font=\footnotesize, align=center, popdark] at (0,-1.75) {Parménide : Un, immuable,\\ éternel, sphérique};
\end{scope}
\end{tikzpicture}
\end{center}
\legende{Schéma dialectique : à gauche, le Devenir héraclitéen (flèches circulaires du changement ordonné par le Logos) ; à droite, l'Être parménidien (sphère parfaite, statique et homogène).}
\end{popfigure}

\courssub{L'aporie fondamentale : Être ou Devenir ?}

\begin{definition}[Aporie]
Une \cle{aporie} (en grec : \emph{impasse}) est une difficulté rationnelle apparemment insoluble, où deux thèses également argumentées semblent s'opposer sans que l'on puisse trancher.
\end{definition}

Le face-à-face Héraclite / Parménide constitue la première grande aporie de la philosophie :
\begin{itemize}
\item \textbf{Problème 1} : Comment penser le \emph{changement} si l'Être ne change pas ? Si Parménide a raison, naître, vieillir, bouger sont des illusions — mais alors pourquoi mes sens me trompent-ils sans cesse ?
\item \textbf{Problème 2} : Comment penser la \emph{stabilité} si tout change ? Si Héraclite a raison, rien ne demeure identique — mais alors comment connaître quoi que ce soit, puisque connaître suppose saisir ce qu'une chose \emph{est} durablement ?
\end{itemize}

Cette aporie n'est pas un échec : elle est \emph{productive}. C'est pour la résoudre que les pluralistes multiplieront les principes (Empédocle : quatre éléments ; Anaxagore : le Noûs ordonnateur ; les Atomistes : atomes éternels se combinant dans le vide) et que Platon, comme nous le verrons dans la section 4, distinguera le monde sensible (devenir) et le monde intelligible (être).

\begin{exemplebox}[Illustration camerounaise : la Sanaga et le Mont Cameroun]
\begin{itemize}
\item \textbf{Héraclite et le fleuve Sanaga} : on ne s'y baigne jamais deux fois dans les mêmes eaux. Ses crues et décrues, son cours qui charrie et transforme, figurent le Devenir. Pourtant, « la Sanaga » demeure une réalité une : c'est le Logos, l'unité dans le changement.
\item \textbf{Parménide et le Mont Cameroun} : le volcan semble immuable, éternel, identique à lui-même à travers les générations — image de l'Être. Mais le géologue sait que cette permanence est \emph{apparente} : éruptions (1999, 2000), érosion, soulèvements lents. Question aporétique : le Mont Cameroun « est-il » vraiment, ou « devient-il » lui aussi ?
\end{itemize}
\end{exemplebox}

\begin{popfigure}
\begin{center}
\begin{tikzpicture}
\begin{axis}[
  width=0.92\linewidth, height=7cm,
  xlabel={Changement $\rightarrow$},
  ylabel={Stabilité $\rightarrow$},
  xmin=0, xmax=10, ymin=0, ymax=10,
  axis lines=left,
  xtick=\empty, ytick=\empty,
  clip=false
]
\addplot[only marks, mark=*, mark size=2.5pt, popblue] coordinates {(8.5,1.5)};
\node[popblue, font=\small, anchor=west] at (axis cs:8.5,1.5) {Héraclite};
\addplot[only marks, mark=*, mark size=2.5pt, poporange] coordinates {(1.2,9)};
\node[poporange, font=\small, anchor=west] at (axis cs:1.2,9) {Parménide};
\addplot[only marks, mark=*, mark size=2.5pt, popgreen] coordinates {(4.5,4)};
\node[popgreen, font=\small, anchor=west] at (axis cs:4.5,4) {Thalès};
\addplot[only marks, mark=*, mark size=2.5pt, popgreen] coordinates {(5.5,5)};
\node[popgreen, font=\small, anchor=south] at (axis cs:5.5,5) {Pythagore};
\addplot[only marks, mark=*, mark size=2.5pt, poppurple] coordinates {(6,6.5)};
\node[poppurple, font=\small, anchor=south west] at (axis cs:6,6.5) {Atomistes};
\addplot[only marks, mark=*, mark size=2.5pt, poppurple] coordinates {(3.5,6)};
\node[poppurple, font=\small, anchor=east] at (axis cs:3.5,6) {Empédocle};
\end{axis}
\end{tikzpicture}
\end{center}
\legende{Positionnement conceptuel des Présocratiques sur deux axes (stabilité / changement). Les pluralistes occupent des positions intermédiaires : ils tentent de concilier l'exigence de Parménide (des principes permanents) et le constat d'Héraclite (le monde change).}
\end{popfigure}

\courssub{Travaux pratiques et méthode d'analyse d'un fragment}

\begin{methode}[Analyser un fragment présocratique en trois étapes]
\begin{enumerate}
\item \textbf{Identifier la thèse} : le fragment défend-il le primat de l'Être (permanence) ou du Devenir (changement) ?
\item \textbf{Repérer l'argument} : le philosophe s'appuie-t-il sur la raison (contre les sens) ou sur l'observation du monde sensible ?
\item \textbf{Formuler le concept clé} : dégager la notion centrale (Logos, Noûs, Unité, archè) et montrer comment elle organise la pensée de l'auteur.
\end{enumerate}
\end{methode}

\begin{experience}[TP : construire un diagramme de Venn comparant Héraclite et Parménide]
\textbf{Matériel} : feuille, crayons de couleur, ou tableau. \textbf{Protocole} :
\begin{enumerate}
\item Tracer deux cercles qui se recoupent : l'un étiqueté « Héraclite », l'autre « Parménide ».
\item Dans la partie propre à Héraclite, placer : \emph{Devenir, Panta rhei, Feu, unité des contraires}.
\item Dans la partie propre à Parménide, placer : \emph{Être, immobilité, éternité, sphère, critique des sens}.
\item Dans l'intersection, placer ce qui leur est commun : \emph{recherche de l'archè, primat de la raison (Logos / Noûs), unité du réel, rejet de l'opinion commune}.
\end{enumerate}
\textbf{Observation et interprétation} : l'intersection, loin d'être vide, montre que les deux philosophes partagent la même exigence rationaliste ; ils divergent seulement sur la \emph{nature} de ce que la raison découvre : un ordre changeant (Logos) ou une réalité immuable (Être).
\end{experience}

\begin{exoresolu}[Analyse d'un fragment d'Héraclite]
Énoncé : Analysez le fragment suivant selon la méthode en trois étapes : \emph{« Le chemin en montée et en descente est un et le même »} (Héraclite).
\tcblower
\textbf{Solution détaillée.}
\textbf{Étape 1 — Thèse} : le fragment relève de la philosophie du Devenir, mais précisément de sa dimension la plus profonde : l'unité des contraires. Il ne dit pas seulement que tout change ; il dit que les opposés (monter / descendre) sont une seule et même réalité.
\textbf{Étape 2 — Argument} : Héraclite s'appuie sur une observation sensible (une même route est montante ou descendante selon le sens du parcours), mais il en tire une conclusion rationnelle et universelle : les contraires ne s'excluent pas, ils se fondent dans l'unité. C'est le Logos, loi d'harmonie des opposés, qui se manifeste ici.
\textbf{Étape 3 — Concept clé} : le concept central est l'\emph{unité des contraires}, expression du Logos. Le fragment illustre que le devenir n'est pas un chaos mais un ordre : la tension des opposés constitue l'harmonie invisible du monde.
\textbf{Conclusion} : ce fragment résume la pensée héraclitéenne : le changement est réel, mais il est gouverné par une loi d'unité que seule la raison éveillée peut saisir.
\end{exoresolu}

\begin{savaistu}[Deux figures fondatrices]
Héraclite était surnommé \emph{« l'Obscur »} (ou « le Ténébreux ») à cause de son style énigmatique et de ses aphorismes volontairement difficiles ; la tradition le dit aussi « pleureur », affligé par la folie des hommes endormis. Parménide, lui, est le fondateur de l'\emph{ontologie}, la science de l'Être, et son disciple Zénon d'Élée inventa les paradoxes célèbres (Achille et la tortue) pour défendre l'immobilité de l'Être. Leur débat — Être contre Devenir — structure toute la métaphysique occidentale : Platon, Aristote, puis Hegel (qui fera d'Héraclite un précurseur de la dialectique) ne cesseront de le reprendre. On ne conserve d'Héraclite qu'environ \textbf{130 fragments}, tandis que le poème de Parménide, \emph{Sur la Nature}, nous est parvenu en partie grâce aux citations de Simplicius.
\end{savaistu}

\begin{exercice}[Pour s'entraîner]
\begin{enumerate}
\item Définis l'archè et donne un exemple chez deux Présocratiques différents.
\item Explique en quoi le Logos d'Héraclite empêche de réduire sa pensée au « tout change ».
\item Reconstitue le raisonnement de Parménide : pourquoi le changement est-il, selon lui, impossible ?
\item Sujet de dissertation : \emph{Faut-il croire les sens ou la raison ?}
\end{enumerate}
\end{exercice}

\begin{corrige}[Exercice — éléments de réponse]
\begin{enumerate}
\item L'archè est le principe premier, l'origine constitutive de toutes choses. Exemples : l'Eau chez Thalès, l'Être chez Parménide.
\item Le Logos est la loi rationnelle universelle qui ordonne le devenir : le changement obéit à une mesure et à une harmonie des contraires. Dire seulement « tout change » oublie cette unité et fait du monde un chaos, ce que Héraclite récuse.
\item Penser c'est penser quelque chose ; le non-être est impensable et n'est pas. Or changer, c'est passer de l'être au non-être (ou naître à partir du non-être). Comme le non-être n'est pas, le changement est impossible : c'est une illusion des sens.
\item Plan possible : I. Les sens trompent et la raison seule atteint la vérité (Parménide) ; II. Mais l'expérience sensible est le point de départ de toute connaissance du monde changeant (Héraclite, observation du fleuve) ; III. Dépassement : la raison doit partir du sensible pour le dépasser (annonce de Platon).
\end{enumerate}
\end{corrige}

\begin{aretenir}[À retenir de la section]
La philosophie naît à Milet au VI\textsuperscript{e} s. av. J.-C. comme recherche rationnelle de l'\cle{archè}. Héraclite affirme le \cle{Devenir} (« tout coule ») ordonné par le \cle{Logos}, loi d'unité des contraires. Parménide affirme l'\cle{Être} Un, immuable et éternel, accessible au seul \cle{noûs} contre les sens trompeurs. Leur opposition constitue l'aporie fondatrice — comment penser ensemble stabilité et changement ? — que Platon tentera de résoudre, comme nous le verrons dans la section 4.
\end{aretenir}

\coursec{Le tournant anthropologique : Socrate, l'éthique par le dialogue}

Après avoir exploré la recherche présocratique de l'\emph{archè} (principe premier) chez Héraclite et Parménide — l'un voyant le monde comme un devenir perpétuel, l'autre affirmant l'immuabilité de l'Être —, nous assistons à un basculement décisif. La philosophie quitte la \emph{physiologie} (étude de la nature) pour se tourner vers l'\emph{anthropologie} (étude de l'homme). Ce « tournant anthropologique » s'opère à Athènes au V\textsuperscript{e} siècle av. J.-C., dans une cité en pleine effervescence démocratique mais traversée par une crise profonde des valeurs traditionnelles.

\courssub{Contexte : Athènes, la crise des valeurs et le défi sophistique}

Au V\textsuperscript{e} siècle, Athènes est au sommet de sa puissance (siècle de Périclès). La démocratie directe permet à tout citoyen de parler à l'\emph{Ecclésia} (assemblée) et de siéger au tribunal (\emph{Héliée}). Pour convaincre, il faut maîtriser l'art du discours : la rhétorique. C'est le marché des \cle{Sophistes} (Protagoras, Gorgias, Hippias...), « maîtres de sagesse » itinérants qui enseignent la \emph{technè} (technique) du succès politique contre rémunération.

Leur thèse centrale, formulée par \cle{Protagoras}, ébranle les fondements de la cité : \og L'homme est la mesure de toutes choses, de celles qui sont en tant qu'elles sont, de celles qui ne sont pas en tant qu'elles ne sont pas \fg{} (\emph{Mesure}, DK 80A1). C'est le \cle{relativisme} : il n'y a pas de vérité objective, pas de justice absolue, seulement des opinions (\emph{doxa}) variables selon les individus ou les cités. La loi (\emph{nomos}) n'est plus d'origine divine mais convention humaine ; le plus fort impose sa loi.

\begin{exemplebox}[La leçon de Protagoras à un élève camerounais]
Imagine un élève de Douala qui affirme : \og La corruption, c'est mal \fg{}. Un sophiste lui répondrait : \og À Douala peut-être, mais dans une autre cité ou une autre époque, c'est la norme pour accélérer les dossiers. Le bien et le mal dépendent du point de vue de celui qui juge. \fg{} Face à ce nihilisme pratique, Socrate refuse que la philosophie devienne un outil de manipulation. Il cherche l'\cle{universel} : ce qui est valable pour tous les hommes, en tout lieu, en tout temps.
\end{exemplebox}

\courssub{La méthode socratique : Ironie, Maïeutique et Soin de l'âme}

Socrate (470–399 av. J.-C.), fils d'un sculpteur et d'une sage-femme, n'a rien écrit. Il pratique la philosophie \emph{dans la rue}, sur l'Agora, au gymnase, auprès des jeunes. Sa méthode est une réponse radicale à la prétention sophistique de \emph{vendre} un savoir tout fait.

\begin{definition}[Maïeutique (Art obstétrique des esprits)]
Terme grec \emph{maïeutikè technè} (art d'accoucher). Socrate, se comparant à sa mère Phénarète, affirme n'enseigner rien : il aide l'interlocuteur à \textbf{enfanter sa propre vérité} par un questionnement rigoureux. Le philosophe n'est pas un dépositaire de savoirs (\emph{sophos}), mais un \emph{philosopho}s (ami du savoir) qui stimule la naissance de la connaissance (\emph{epistémè}) à partir de l'opinion (\emph{doxa}).
\end{definition}

La méthode se déploie en trois temps indissociables :

\begin{enumerate}
    \item \textbf{L'Ironie (\emph{eirôneia}) :} Socrate feint l'ignorance (\og Je sais que je ne sais rien \fg{}, \emph{Apologie} 21d). Il interroge l'expert (le général sur le courage, le prêtre sur la piété) pour le mettre face à sa propre insuffisance. Ce n'est pas de la moquerie, mais une \textbf{purification} : vider l'esprit de ses fausses certitudes.
    \item \textbf{L'Élenchos (Réfutation) et la Maïeutique :} Par une série de questions (\emph{ti estin ?} : « Qu'est-ce que c'est ? »), Socrate traque les contradictions. L'interlocuteur propose une définition (ex. : \og Le courage, c'est tenir bon au combat \fg{}), Socrate oppose un contre-exemple (ex. : \og Et le cavalier qui fuit pour mieux revenir charger ? \fg{}). L'interlocuteur tombe dans l'\cle{apore} (impasse, perplexité). C'est le point de départ nécessaire : on ne cherche vraiment que quand on sait qu'on ne sait pas.
    \item \textbf{Le Soin de l'âme (\emph{epiméleia tès psychès}) :} La finalité n'est pas intellectuelle mais \textbf{existentielle}. « Le plus important, ce n'est pas de vivre, c'est de vivre bien » (\emph{Criton} 48b). L'âme est le siège de la raison et de la morale. Philosopher, c'est soigner son âme de l'ignorance, source de tous les maux.
\end{enumerate}

\begin{popfigure}
\begin{tikzpicture}[
    node distance=1.2cm and 2.5cm,
    every node/.style={font=\small, align=center},
    box/.style={draw=popblue, thick, rounded corners, fill=popblue!5, minimum width=3.2cm, minimum height=1.4cm},
    arrow/.style={-Stealth, thick, popdark},
    start/.style={draw=popgreen, thick, rounded corners, fill=popgreen!10, minimum width=3.5cm},
    end/.style={draw=poporange, thick, rounded corners, fill=poporange!10, minimum width=3.5cm},
    apore/.style={draw=poppink, thick, rounded corners, fill=poppink!10, minimum width=3.5cm}
]
% Nœuds
\node[start, align=center] (start){Question initiale\\ \og Qu'est-ce que X ? \fg{}\\(Courage, Justice, Piété...)};
\node[box, right=of start, align=center] (doxa){Réponse spontanée\\(Doxa / Opinion)\\\textit{Ex: « Tenir bon »}};
\node[box, right=of doxa, align=center] (elenchos){Contre-exemple\\(Élenchos / Réfutation)\\\textit{Ex: « Fuir pour mieux charger »}};
\node[apore, below=of elenchos, align=center] (apore){APORIE (Blocage)\\\og Je ne sais plus \fg{}\\\textbf{Prise de conscience de l'ignorance}};
\node[box, right=of elenchos, align=center] (hypo){Nouvelle hypothèse\\plus rigoureuse\\\textit{Ex: « Endurance de l'âme »}};
\node[end, right=of hypo, align=center] (episteme){Définition approchée\\(Epistémè / Science)\\\textbf{Essence commune (Eidos)}};
% Flèches
\draw[arrow] (start) -- (doxa) node[midway, above] {Interrogation};
\draw[arrow] (doxa) -- (elenchos) node[midway, above] {Test};
\draw[arrow] (elenchos) -- (apore) node[midway, right] {Contradiction};
\draw[arrow] (apore) -- (hypo) node[midway, below] {Recherche};
\draw[arrow] (hypo) -- (episteme) node[midway, above] {Validation};
% Boucle possible
\draw[arrow, popmuted, dashed] (hypo) to[out=180, in=90, looseness=1.5] (elenchos) node[midway, left, align=center] {Nouveau\\test};
\end{tikzpicture}
\legende{Schéma de la dialectique socratique (Maïeutique) : de l'opinion (Doxa) à la science (Epistémè) par le passage obligé de l'aporie. La flèche pointillée indique la nature itérative du processus.}
\end{popfigure}

\courssub{Le problème de la définition : À la recherche de l'Essence (Eidos)}

Pour Socrate, savoir agir (vertu) suppose savoir \emph{ce qu'est} la chose. On ne peut être courageux si l'on ne définit pas le Courage. La question \og \emph{Ti estin ?} \fg{} (Qu'est-ce que c'est ?) vise l'\cle{Eidos} (Forme, Essence, Idée) : ce qui est commun à tous les actes courageux, au-delà de leurs différences particulières.

\begin{exoresolu}[Analyse d'un dialogue : \emph{Lachès} — La définition du courage]
\textbf{Énoncé :} Lachès, général athénien, propose trois définitions successives du courage. Socrate les réfute par l'élenchos.
\tcblower
\textbf{Solution détaillée :}
\begin{enumerate}
    \item \textbf{Définition 1 :} \og Le courage, c'est rester à son rang au combat sans fuir. \fg{} (Définition \emph{particulière}, militaire, statique).
    \item \textbf{Réfutation :} Socrate cite les Scythes (cavaliers qui fuient pour charger) ou les hoplites de Sparte (qui reculent stratégiquement). La définition est trop étroite : elle exclut des formes légitimes de courage.
    \item \textbf{Définition 2 :} \og Le courage, c'est une endurance de l'âme. \fg{} (Définition \emph{générale}, psychologique).
    \item \textbf{Réfutation :} L'endurance peut être folie (tenir bon stupidement). Il faut y ajouter la \emph{sagesse} (savoir ce qu'on fait). Courage = Endurance \textbf{+} Sagesse.
    \item \textbf{Définition 3 :} \og Le courage, c'est la science de ce qui est à craindre et de ce qui ne l'est pas. \fg{} (Identification Vertu = Science).
    \item \textbf{Apore finale :} Si le courage est une science, quelle est son objet ? Les biens et les maux futurs. Mais alors le courage devient une science générale du bien et du mal, c'est-à-dire \emph{toute} la vertu. Le dialogue s'achève sans définition close, mais l'exigence de l'universel est posée.
\end{enumerate}
\textbf{Enseignement :} La définition socratique exige \textbf{l'universalité} (valable pour tous les cas), \textbf{l'essentialité} (ce qui fait \emph{être} la chose) et \textbf{l'opérationnalité} (permettre de juger les cas nouveaux).
\end{exoresolu}

\courssub{Le paradoxe moral : « Nul n'est méchant volontairement »}

C'est le cœur de l'éthique socratique, exposé notamment dans le \emph{Protagoras} et le \emph{Ménon}.

\begin{propriete}[Théorème socratique : La Vertu est Science (\emph{Aretè = Epistémè})]
\textbf{Énoncé :} Le mal est ignorance (\emph{amathia}). Nul ne fait le mal volontairement (\emph{oudeis hekôn hamartanei}). Si l'on connaît véritablement le Bien, on le fait nécessairement. La vertu (courage, tempérance, justice, piété) n'est pas une habitude ou un don naturel, mais une \textbf{connaissance} (science du Bien et du Mal).
\textbf{Preuve (schématique) :}
\begin{enumerate}
    \item Tout être humain désire le bonheur (eudaimonia).
    \item Le bonheur s'obtient par les biens (santé, richesse, honneur, vertu).
    \item Les biens « extérieurs » (richesse, pouvoir) peuvent nuire s'ils sont mal utilisés (ex. : un tyran riche et malade).
    \item Seul l'usage \textbf{sage} (dirigé par la connaissance) des biens rend heureux.
    \item Donc, la seule chose bonne en soi est la \textbf{sagesse / science} (connaissance de l'usage des choses).
    \item Celui qui fait le mal (injuste, lâche) croit à tort que cela lui profitera. Il se trompe sur le vrai Bien. Il est \emph{ignorant}, non \emph{méchant} par nature.
\end{enumerate}
\end{propriete}

\begin{attention}[Piège fréquent : Confondre « Science » au sens moderne et « Science » socratique]
La « science » socratique n'est pas une connaissance technique (savoir faire une chaussure, calculer une trajectoire). C'est une \textbf{sagesse pratique} (\emph{phronésis}) : la connaissance de ce qui est bon pour l'homme, de la valeur des choses. C'est une \emph{connaissance éthique} qui transforme l'être. Au Bac, ne réduisez pas « Vertu = Science » à « Il faut aller à l'école pour être gentil ». C'est une thèse ontologique : l'ignorance du Bien est la cause unique du mal.
\end{attention}

\courssub{Socrate et la Cité : Cohérence de la vie et de la pensée}

Socrate refuse la carrière politique (\emph{Apologie} 31c-32a) : son \emph{daimonion} (voix intérieure) l'en empêche. Il sert la cité autrement : en exerçant sa mission divine d'\emph{éperon} (taon) qui pique le cheval athénien endormi pour le réveiller. Son procès (-399) et sa mort illustrent l'accord parfait entre sa parole et son acte.

\begin{exemplebox}[Le procès de -399 : Chiffres et symbole]
\begin{itemize}
    \item \textbf{Tribunal :} Héliée, 501 juges (héliastes) tirés au sort parmi les citoyens de plus de 30 ans.
    \item \textbf{Accusations :} 1. Corrompre la jeunesse. 2. Ne pas croire aux dieux de la cité, introduire des divinités nouvelles (le daimonion).
    \item \textbf{Vote culpabilité :} 280 voix \textbf{contre} 221 (majorité relative de 59 voix). Socrate manque de peu l'acquittement.
    \item \textbf{Peine proposée :} Socrate propose... d'être nourri au Prytanée (honneur réservé aux vainqueurs olympiques) ou une amende dérisoire (1 mine, puis 30 mines grâce à ses amis). Provocation suprême.
    \item \textbf{Vote peine :} 360 voix pour la mort (ciguë) contre 141. Le durcissement est dû à l'arrogance perçue de sa défense.
    \item \textbf{Refus de l'évasion (Criton) :} Ses amis ont tout prévu pour le faire fuir. Il refuse : \og On ne doit pas rendre le mal pour le mal \fg{} ; désobéir aux lois, c'est détruire la cité qui l'a fait naître, éduqué, protégé. \textbf{Obéir aux lois, c'est obéir à la Justice.}
\end{itemize}
\end{exemplebox}

\begin{savaistu}[Le Daimonion : une conscience morale anticipée]
Le \emph{daimonion} n'est pas un « démon » au sens chrétien (maléfique), ni un dieu olympien. C'est un \textbf{signe intérieur}, une \emph{phônè} (voix) qui \textbf{n'intervient qu'en négatif} : « Il me détourne toujours de ce que je vais faire, jamais ne m'y pousse » (\emph{Apologie} 31d). C'est une \textbf{intuition morale immédiate}, antérieure au raisonnement, qui protège l'intégrité de l'âme. On peut le lire comme la métaphore de la \emph{conscience morale} : ce « quelque chose en moi » qui dit « Non » avant même que je calcule les conséquences.
\end{savaistu}

\courssub{Ancrage local : La Palabre africaine, espace de maïeutique collective ?}

Le dialogue socratique trouve un écho saisissant dans la tradition africaine de la \textbf{Palabre} (du portugais \emph{palavra}, « parole »).

\begin{center}
\begin{tabularx}{\linewidth}{|l|X|X|}
\hline
\rowcolor{popblueL} \textbf{Critère} & \textbf{Dialogue Socratique (Athènes)} & \textbf{Palabre Africaine (ex. Bamiléké, Bassa, Douala)} \\ \hline
\textbf{But} & Recherche de la \textbf{Vérité objective} (Eidos), définition universelle. & Recherche du \textbf{Consensus social}, restauration de l'harmonie du groupe. \\ \hline
\textbf{Structure} & Dyadique (Socrate + 1 interlocuteur), asymétrique (maître/élève apparent). & Collectif (Conseil des sages, notables, famille élargie), circulaire. \\ \hline
\textbf{Méthode} & Élenchos (réfutation), Ironie, Questionnement \emph{ti estin}. & Argumentation, Proverbes, Médiation, Compromis, Serment. \\ \hline
\textbf{Rôle de la parole} & Outil logique pour tester la cohérence. & Acte social engageant l'honneur, la parole donnée (\emph{ngondo}, \emph{mboa}). \\ \hline
\textbf{Issue} & Souvent l'Aporie (savoir qu'on ne sait pas). & Décision contraignante, réconciliation, sanction réparatrice. \\ \hline
\end{tabularx}
\end{center}

\begin{aretenir}[Similitude profonde : L'exigence du \textbf{Logos}]
Dans les deux cas, la \textbf{parole raisonnée} (\emph{logos}) prime sur la violence ou l'arbitraire. On ne tranche pas par la force, mais par l'argument. La palabre, comme le dialogue socratique, est un \textbf{dispositif d'intelligence collective} : l'intelligence du groupe (ou du couple dialectique) dépasse la somme des intelligences individuelles. \textbf{Différence cruciale :} Socrate vise l'universel abstrait (le Juste \emph{en soi}) ; la palabre vise le particulier concret (la justice \emph{ici et maintenant} pour ces personnes).
\end{aretenir}

\courssub{Atelier méthodologique (TP) : Simulation d'un dialogue socratique}

\textbf{Thème :} \og Qu'est-ce que la réussite scolaire ? \fg{}

\begin{methode}[Déroulement de la séance (45 min)]
\begin{enumerate}
    \item \textbf{Phase 1 : Doxa initiale (5 min).} Chaque élève écrit sa définition spontanée sur un papier : \og Réussir, c'est... \fg{} (ex. : avoir 18/20, être premier, avoir le Bac, apprendre un métier, faire fierté aux parents).
    \item \textbf{Phase 2 : Mise en commun \& Élenchos (15 min).} L'enseignant (rôle de Socrate) collecte 3-4 définitions au tableau. La classe propose des \textbf{contre-exemples} pour chaque :
    \begin{itemize}
        \item \og Avoir 18/20 \fg{} $\rightarrow$ Contre-exemple : Un élève qui triche et a 18/20 a-t-il réussi ?
        \item \og Avoir le Bac \fg{} $\rightarrow$ Contre-exemple : Un bachelier au chômage, déprimé, a-t-il réussi sa vie ?
        \item \og Faire fierté aux parents \fg{} $\rightarrow$ Contre-exemple : Des parents fiers d'un fils corrompu ?
    \end{itemize}
    \item \textbf{Phase 3 : Apore \& Nouvelle hypothèse (15 min).} Face aux contradictions, les définitions s'affinent. On cherche l'\textbf{invariant} : \og La réussite, c'est l'acquisition de \textbf{compétences} (savoirs, savoir-faire, savoir-être) permettant l'\textbf{autonomie} et la \textbf{contribution sociale}. \fg{}
    \item \textbf{Phase 4 : Construction de l'arbre des définitions (10 min).} Collectivement, on trace l'arbre au tableau (voir figure ci-dessous).
\end{enumerate}
\end{methode}

\begin{popfigure}
\begin{center}\tcbox[colback=popink,colframe=popink,coltext=white,arc=9pt,boxsep=4pt,left=6pt,right=6pt]{\bfseries\small Réussite Scolaire?}\end{center}
\vspace{-2pt}
\begin{tcbraster}[raster columns=2,raster equal height=rows,raster column skip=6pt,raster row skip=6pt,enhanced,blankest]
\begin{tcolorbox}[enhanced,colback=white,colframe=poppink,boxrule=0.9pt,arc=3pt,title={Avoir de bonnes notes},fonttitle=\bfseries\scriptsize,coltitle=white,colbacktitle=poppink,left=4pt,right=4pt,top=2pt,bottom=2pt,boxsep=1pt,toptitle=1.5pt,bottomtitle=1.5pt,halign=left]
\begin{itemize}[nosep,leftmargin=*,itemsep=1pt,font=\scriptsize]
\item Contre-ex: Tricherie
\item Contre-ex: Par cœur sans comprendre
\end{itemize}
\end{tcolorbox}
\begin{tcolorbox}[enhanced,colback=white,colframe=poppurple,boxrule=0.9pt,arc=3pt,title={Obtenir le Bac},fonttitle=\bfseries\scriptsize,coltitle=white,colbacktitle=poppurple,left=4pt,right=4pt,top=2pt,bottom=2pt,boxsep=1pt,toptitle=1.5pt,bottomtitle=1.5pt,halign=left]
\begin{itemize}[nosep,leftmargin=*,itemsep=1pt,font=\scriptsize]
\item Contre-ex: Chômage diplômé
\item Contre-ex: Bac « acheté »
\end{itemize}
\end{tcolorbox}
\begin{tcolorbox}[enhanced,colback=white,colframe=popgold,boxrule=0.9pt,arc=3pt,title={Faire fierté aux parents},fonttitle=\bfseries\scriptsize,coltitle=white,colbacktitle=popgold,left=4pt,right=4pt,top=2pt,bottom=2pt,boxsep=1pt,toptitle=1.5pt,bottomtitle=1.5pt,halign=left]
\begin{itemize}[nosep,leftmargin=*,itemsep=1pt,font=\scriptsize]
\item Contre-ex: Parents complices
\item Contre-ex: Conformisme
\end{itemize}
\end{tcolorbox}
\begin{tcolorbox}[enhanced,colback=white,colframe=popgreen,boxrule=0.9pt,arc=3pt,title={\textbf{Acquérir des} \textbf{ compétences} \textbf{ pour l'autonomie}},fonttitle=\bfseries\scriptsize,coltitle=white,colbacktitle=popgreen,left=4pt,right=4pt,top=2pt,bottom=2pt,boxsep=1pt,toptitle=1.5pt,bottomtitle=1.5pt,halign=left]
\begin{itemize}[nosep,leftmargin=*,itemsep=1pt,font=\scriptsize]
\item Savoirs validés
\item Savoir-faire techniques
\item Savoir-être (éthique)
\item Insertion pro / civique
\end{itemize}
\end{tcolorbox}
\end{tcbraster}

\legende{Arbre des définitions successives construit collectivement : élagage des définitions partielles (Doxa) par l'élenchos vers une définition essentielle (Epistémè approchée).}
\end{popfigure}

\courssub{Exercice d'application : Explication de texte guidée}

\begin{exercice}[Entraînement au Bac : \emph{Apologie de Socrate} (21d)]
\textbf{Texte :} \og Je me mis donc à l'examiner, au nom du dieu, et je raisonnai ainsi : « Celui-ci, me dis-je, se croit sage, sans l'être ; moi, je ne me crois pas sage, et je ne le suis pas non plus. Il est donc probable que je suis un peu plus sage que lui, en ce que je ne crois pas savoir ce que je ne sais pas. » \fg{} (Platon, \emph{Apologie de Socrate}, 21d).
\textbf{Questions :}
\begin{enumerate}
    \item Situer le texte (auteur, œuvre, contexte historique, circonstances de l'énonciation).
    \item Expliquer la thèse défendue : en quoi consiste la « sagesse » socratique ?
    \item Analyser l'argumentation : identifier la structure du raisonnement (comparaison, ironie, conclusion).
    \item Discuter : « Savoir qu'on ne sait rien » est-il une connaissance positive ou une simple attitude négative ? Illustrer par un exemple contemporain (ex. : scientifique face à l'inconnu, citoyen face aux fake news).
\end{enumerate}
\end{exercice}

\begin{corrige}[Exercice n°16 - Éléments de correction]
\begin{enumerate}
    \item \textbf{Contexte :} Platon (disciple), \emph{Apologie} (discours de défense), -399, procès de Socrate à l'Héliée. Socrate répond à l'accusation de « corrompre la jeunesse » en expliquant l'origine de sa mission (oracle de Delphes : « Nul n'est plus sage que Socrate »).
    \item \textbf{Thèse :} La sagesse humaine n'est pas un savoir positif (comme la technique du cordonnier), mais une \textbf{conscience de ses limites} (\emph{sophia} = savoir qu'on ne sait pas). C'est une \textbf{sagesse négative} (par opposition à l'ignorance présomptueuse des politiciens/poètes/artisans).
    \item \textbf{Argumentation :}
    \begin{itemize}
        \item Prémisse 1 : L'oracle dit « Socrate est le plus sage ».
        \item Prémisse 2 : Socrate sait qu'il ne sait pas (vérification empirique chez les « sages » réputés).
        \item Prémisse 3 : Les autres croient savoir sans savoir (ignorance double).
        \item Conclusion : Socrate est plus sage \emph{de ce petit rien} : il ne s'ajoute pas l'erreur de croire savoir.
    \end{itemize}
    L'ironie est l'outil : feindre de prendre l'oracle au pied de la lettre pour tester les réputés sages.
    \item \textbf{Discussion :} C'est une \textbf{connaissance méta-cognitive} positive (savoir \emph{sur} son savoir). Exemple : Un bon médecin dit « Je ne sais pas encore quel virus c'est, il faut faire des analyses » (sagesse socratique) vs un charlatan qui dit « C'est tel virus, prenez ceci » (ignorance présomptueuse). Face aux fake news : l'attitude socratique est la \emph{vérification} (vérifier la source, croiser les infos) avant de partager, admettant qu'on ne sait pas encore la vérité.
\end{enumerate}
\end{corrige}

\begin{methode}[Réussir l'explication de texte socratique (Bac)]
\begin{enumerate}
    \item \textbf{Identifier le stade du dialogue :} Est-on dans l'Ironie (début), l'Élenchos (milieu, réfutation), la Maïeutique (construction), l'Apore (fin) ?
    \item \textbf{Repérer la question « Ti estin ? » :} Quel concept est visé (Courage, Piété, Justice, Science, Vertu) ?
    \item \textbf{Analyser la structure argumentative :} Définition proposée $\rightarrow$ Contre-exemple $\rightarrow$ Réfutation $\rightarrow$ Nouvelle définition.
    \item \textbf{Ne pas confondre Socrate et Platon :} Dans les dialogues de jeunesse (\emph{Apologie}, \emph{Criton}, \emph{Lachès}, \emph{Charmide}, \emph{Euthyphron}), Socrate ne donne \textbf{jamais} la réponse finale. Il questionne. La théorie des Idées (monde suprasensible) est platonicienne (dialogues de maturité : \emph{République}, \emph{Phédon}, \emph{Banquet}).
\end{enumerate}
\end{methode}

\begin{popfigure}
\begin{tikzpicture}[
    scale=0.85,
    every node/.style={font=\small},
    city/.style={circle, fill=poporange, draw=popdark, thick, minimum size=6pt, inner sep=1pt},
    place/.style={circle, fill=popblue, draw=popdark, thick, minimum size=5pt, inner sep=1pt},
    building/.style={rectangle, fill=poppurple!60, draw=popdark, minimum size=5pt, inner sep=1pt},
    label/.style={anchor=north, yshift=-2pt, text=popdark}
]
% Coastline / simplified Attica
\draw[popblue!30, line width=3pt] (0,0) .. controls (1,1) and (3,1.5) .. (5,1) .. controls (6,0.5) and (6,-0.5) .. (5,-1) .. controls (3,-1.5) and (1,-1) .. (0,0);
% City walls (Long Walls schematic)
\draw[popgray, line width=1.5pt] (2.5,0.5) -- (5.5,0.5) (2.5,-0.5) -- (5.5,-0.5);
% Piraeus
\node[city] (piraeus) at (5.5,0) {};
\node[label] at (5.5,-0.8) {Pirée (Port)};
% Athens Center
\node[city] (athens) at (2.5,0) {};
\node[label] at (2.5,-0.8) {Athènes (Asty)};
% Acropolis
\node[building, minimum width=12pt, minimum height=8pt] (acropolis) at (2.0, 0.6) {};
\node[label] at (2.0, 1.1) {Acropole};
% Agora
\node[building, minimum width=10pt, minimum height=6pt, fill=popgreen!60] (agora) at (2.6, 0.1) {};
\node[label] at (2.6, -0.4) {Agora};
% Pnyx (Assembly)
\node[place] (pnyx) at (1.7, -0.4) {};
\node[label, anchor=east] at (1.4, -0.4) {Pnyx (Assemblée)};
% Areopagus
\node[place] (areopagus) at (1.3, 0.3) {};
\node[label, anchor=east] at (1.0, 0.3) {Aréopage};
% Heliaia (Court - approx location near Agora)
\node[building, minimum width=8pt, minimum height=8pt, fill=poppink!60] (heliaia) at (3.0, -0.2) {};
\node[label] at (3.0, -0.6) {Héliée (Procès)};
% Lyceum / Academy (outside walls)
\node[place] (lyceum) at (3.5, 1.2) {};
\node[label] at (3.5, 1.5) {Lycée (Socrate)};
\node[place] (academy) at (0.8, 1.2) {};
\node[label] at (0.8, 1.5) {Académie (Platon)};
% North Arrow
\node[anchor=south] at (5.5, 1.5) {\Large $\uparrow$ N};
% Legend
\begin{scope}[xshift=6.5cm, yshift=-1cm]
    \node[city] (l1) at (0,0) {}; \node[label, anchor=west] at (0.3,0) {Ville / Lieu clé};
    \node[building, fill=popgreen!60] (l2) at (0,-0.6) {}; \node[label, anchor=west] at (0.3,-0.6) {Agora (Marché/Débat)};
    \node[building, fill=poppink!60] (l3) at (0,-1.2) {}; \node[label, anchor=west] at (0.3,-1.2) {Héliée (Tribunal)};
    \node[building, fill=poppurple!60] (l4) at (0,-1.8) {}; \node[label, anchor=west] at (0.3,-1.8) {Acropole (Religieux)};
    \node[place] (l5) at (0,-2.4) {}; \node[label, anchor=west] at (0.3,-2.4) {Lieux de dialogue};
\end{scope}
\end{tikzpicture}
\legende{Carte schématique d'Athènes au V\textsuperscript{e} s. av. J.-C. : L'Agora (cœur civique et philosophique), l'Héliée (lieu du procès), l'Acropole (religieux), le Pnyx (politique). Les Longs Murs relient la cité au Pirée (port).}
\end{popfigure}

\begin{popfigure}
\begin{tikzpicture}
% Word Cloud style using nodes with varying sizes and colors
\pgfmathsetseed{42}
\foreach \word/\size/\color/\angle/\radius in {
    {Ironie/2.2/popblue/10/2.5},
    {Maïeutique/2.5/poporange/45/2.0},
    {Élenchos/2.0/popgreen/80/2.8},
    {Soin de soi/1.8/poppurple/120/2.2},
    {Vertu=Science/2.3/popgold/160/2.6},
    {Daimonion/1.9/poppink/200/2.4},
    {Apore/1.7/popblue/240/2.1},
    {Eidos/1.6/poporange/280/2.7},
    {Doxa/1.5/popgreen/320/1.9},
    {Epistémè/1.8/poppurple/350/2.3},
    {Criton/1.4/popgold/20/1.5},
    {Apologie/1.4/poppink/60/1.7},
    {Protagoras/1.3/popblue/100/1.6},
    {Relativisme/1.3/poporange/140/1.8},
    {Universel/1.6/popgreen/180/2.0},
    {Palabre/1.5/poppurple/220/1.9},
    {Consensus/1.4/popgold/260/1.6},
    {Vérité/1.7/poppink/300/2.1},
    {Ignorance/1.5/popblue/340/1.8},
    {Sagesse/2.0/poporange/15/2.2}} {
    \node[rotate=\angle, text=\color, font=\bfseries\fontsize{\size em}{1em}\selectfont, opacity=0.85] at (\angle:\radius) {\word};
}
% Central concept
\node[text=popdark, font=\bfseries\Large, fill=popcream, rounded corners=3pt, inner sep=4pt] at (0,0) {SOCRATE};
\end{tikzpicture}
\legende{Nuage de concepts socratiques : taille $\approx$ importance centrale. Au centre, Socrate ; autour, les piliers de sa méthode (Ironie, Maïeutique, Élenchos), sa thèse éthique (Vertu=Science, Daimonion), ses concepts épistémologiques (Doxa, Epistémè, Eidos, Apore), ses textes fondateurs (Apologie, Criton) et le dialogue avec son époque (Protagoras, Relativisme) et le nôtre (Palabre).}
\end{popfigure}

\courssub{Synthèse de la section : Socrate, fondateur de l'éthique rationnelle}

Avec Socrate, la philosophie devient \textbf{pratique spirituelle} et \textbf{exigence universelle}.
\begin{itemize}
    \item \textbf{Contre le relativisme sophistique :} Il ancre la morale dans la \textbf{connaissance objective} (Science du Bien), non dans l'opinion changeante.
    \item \textbf{Méthode :} Le dialogue rationnel (\emph{dialectique}) remplace la rhétorique persuasive. La vérité s'élabore \emph{avec} l'autre, par le conflit des arguments (Élenchos), non par l'autorité du maître.
    \item \textbf{Unité vie/pensée :} Le procès et la mort (ciguë) scellent l'accord entre le \emph{logos} (discours) et le \emph{bios} (vie). « Une vie sans examen ne vaut pas d'être vécue » (\emph{Apologie} 38a).
    \item \textbf{Héritage direct :} Platon (théorie des Idées, politique), Aristote (éthique à Nicomaque, logique), Stoïciens (soin de soi, intériorité), Christianisme (examen de conscience, martyr), Lumières (autonomie de la raison), Éducation moderne (pédagogie active, maïeutique).
\end{itemize}

La section suivante verra comment \textbf{Platon}, son disciple, systématise cet héritage en construisant une métaphysique (Théorie des Idées), une politique (République) et une pédagogie (Éducation des gardiens) pour garantir l'accès à ce Bien que Socrate a seulement cherché.

\coursec{La systématisation platonicienne : Être, Idées, Politique et Éducation}

Avec Socrate, nous avons vu la philosophie se retourner vers l'homme : le dialogue, la recherche des définitions morales (qu'est-ce que le courage ? la justice ? la piété ?), l'exigence du « connais-toi toi-même ». Mais Socrate n'a rien écrit et mourut sans avoir trouvé de définition pleinement satisfaisante. C'est son disciple Platon (427--347 av. J.-C.) qui va transformer cette quête socratique en un \cle{système} complet : une doctrine de l'être (ontologie), une doctrine de l'âme (anthropologie), une doctrine de la connaissance (épistémologie) et une doctrine de la cité (politique). Comprendre Platon, c'est comprendre comment la philosophie antique atteint sa forme achevée : non plus seulement des questions, mais une \emph{architecture} du réel.

\courssub{Genèse du platonisme : la synthèse de Socrate, Héraclite et Parménide}

Platon hérite de trois enseignements apparemment contradictoires, et son génie consiste à les réconcilier.

\begin{enumerate}
\item \textbf{De Socrate}, il retient la méthode : la recherche des \cle{définitions} universelles en morale. Quand Socrate demande « qu'est-ce que la justice ? », il suppose qu'il existe \emph{quelque chose} de stable, la justice elle-même, commune à toutes les actions justes. Mais où se trouve ce « quelque chose » ?
\item \textbf{D'Héraclite}, il retient que le monde sensible est en \cle{devenir} perpétuel : « tout coule », rien ne demeure. Or on ne peut connaître véritablement ce qui change sans cesse : le sensible ne peut donc pas être l'objet de la science.
\item \textbf{De Parménide}, il retient que l'\cle{être} véritable est un, immuable, éternel, et qu'il ne se saisit que par la pensée, non par les sens.
\end{enumerate}

La solution de Platon est d'une élégance radicale : il y a \emph{deux mondes}. Le monde sensible, celui que nous percevons, est bien héraclitéen : multiple, changeant, corruptible. Mais il existe un autre monde, le monde \cle{intelligible} (\emph{topos hyperouranios}, « lieu au-delà du ciel »), qui est parménidien : un, éternel, immuable. C'est là que résident les définitions que cherchait Socrate : la Justice en soi, la Beauté en soi, le Bien en soi. C'est la \cle{théorie des Idées} (ou des Formes, \emph{eidos} en grec).

\begin{definition}[L'Idée (Eidos)]
L'\cle{Idée} (grec \emph{eidos}, « forme », « aspect ») est l'\emph{essence immuable} d'une réalité, son modèle parfait, et la cause de l'être des choses sensibles qui lui ressemblent. L'Idée n'est pas une pensée subjective dans ma tête : elle existe réellement, objectivement, dans le monde intelligible, indépendamment de nous.
\end{definition}

\begin{exemplebox}[L'Idée de Cercle et les cercles du tableau]
Le professeur trace un cercle au tableau à Nkongsamba. Ce cercle est imparfait : la craie tremble, le trait est épais, la courbe n'est jamais exactement fermée. Pourtant, tous les élèves le reconnaissent comme un cercle. Comment est-ce possible, si aucun cercle sensible n'est parfait ? Réponse de Platon : parce que notre âme connaît l'\emph{Idée de Cercle} — le cercle parfait, dont tous les points sont rigoureusement équidistants du centre — et qu'elle juge les cercles tracés en les comparant à ce modèle. Le cercle de craie \emph{participe} de l'Idée de Cercle ; il en est une copie dégradée.
\end{exemplebox}

\courssub{Ontologie : les deux mondes et la participation}

L'ontologie platonicienne (doctrine de l'être) oppose rigoureusement deux ordres de réalité :

\begin{center}
\begin{tabularx}{\linewidth}{|l|X|X|}\hline
\rowcolor{popblueL} \textbf{Critère} & \textbf{Monde intelligible (les Idées)} & \textbf{Monde sensible (les choses)} \\ \hline
Exemples & Beauté, Justice, Bien, Égalité, Cercle & belles choses, actions justes, cercles tracés \\ \hline
Quantité & Un (chaque Idée est unique) & Multiple (une infinité de choses) \\ \hline
Temps & Éternel, hors du temps & Temporel, naît et périt \\ \hline
Stabilité & Immuable, toujours identique à soi & Changeant, en devenir \\ \hline
Accès & Saisi par la pensée, la raison & Perçu par les sens \\ \hline
Statut & Cause, modèle, pleinement réel & Copie, image, réalité dégradée \\ \hline
Connaissance & Science (\emph{épistémè}) & Opinion (\emph{doxa}) \\ \hline
\end{tabularx}
\end{center}

Le lien entre les deux mondes s'appelle la \cle{participation} (\emph{methexis}) ou l'\cle{imitation} (\emph{mimèsis}) : les choses sensibles « participent » des Idées, comme une photo participe de son modèle. Une action est juste parce qu'elle participe de la Justice en soi ; un cheval est beau parce qu'il participe, imparfaitement, de la Beauté.

\begin{propriete}[Le Bien, Idée suprême]
Au sommet du monde intelligible règne l'Idée du \cle{Bien} (\emph{Agathon}). Platon affirme qu'elle est « au-delà de l'être en dignité et en puissance » (\emph{République}, 509b). Comme le soleil rend à la fois visibles les choses et leur donne de croître, le Bien est à la fois la source de l'\emph{intelligibilité} des Idées (on ne peut les connaître qu'éclairé par lui) et la source de leur \emph{existence}. Connaître le Bien est le terme de toute l'éducation philosophique.
\end{propriete}

\begin{popfigure}
\begin{center}
\begin{tikzpicture}[scale=1, every node/.style={font=\small}]
% Ligne divisée (République VI)
\draw[thick, popdark] (0,0) -- (10,0);
\foreach \x in {0,3,6,8,10} \draw[thick, popdark] (\x,-0.15) -- (\x,0.15);
% Segments du bas : objets
\node[below, align=center] at (1.5,-0.2) {\textbf{Images}\\ ombres, reflets};
\node[below, align=center] at (4.5,-0.2) {\textbf{Choses sensibles}\\ animaux, objets};
\node[below, align=center] at (7,-0.2) {\textbf{Hypothèses}\\ mathématiques};
\node[below, align=center] at (9,-0.2) {\textbf{Idées}\\ (Bien, Justice)};
% Segments du haut : facultés
\node[above, align=center, popblue] at (1.5,0.2) {\emph{Eikasia}\\ imagination};
\node[above, align=center, popblue] at (4.5,0.2) {\emph{Pistis}\\ croyance};
\node[above, align=center, popblue] at (7,0.2) {\emph{Dianoia}\\ pensée discursive};
\node[above, align=center, popblue] at (9,0.2) {\emph{Noèsis}\\ intelligence};
% Grandes divisions : remplacement des accolades par des traits et nœuds
\draw[poporange, thick] (0,1.1) -- (6,1.1);
\node[poporange, above, align=center] at (3,1.35) {\textbf{Monde sensible}\\ Opinion (\emph{doxa})};
\draw[popgreen, thick] (6,1.1) -- (10,1.1);
\node[popgreen, above, align=center] at (8,1.35) {\textbf{Monde intelligible}\\ Science (\emph{épistémè})};
\end{tikzpicture}
\end{center}
\legende{La Ligne divisée (\emph{République}, VI) : quatre degrés de réalité et de connaissance, du plus obscur (les images) au plus lumineux (les Idées).}
\end{popfigure}

\courssub{Anthropologie : l'âme immortelle et tripartite}

L'homme, chez Platon, est un être \emph{partagé entre les deux mondes}. Le \cle{corps} appartient au monde sensible : il naît, vieillit, meurt. L'\cle{âme} appartient au monde intelligible : elle est immortelle, elle a existé avant notre naissance et survivra à notre mort. C'est le \cle{dualisme} anthropologique platonicien.

De cette préexistence de l'âme découle une théorie célèbre de la connaissance : la \cle{réminiscence} (\emph{anamnèse}). Si l'âme a contemplé les Idées avant de s'incarner, alors connaître n'est pas découvrir du neuf, mais \emph{se ressouvenir} de ce que l'âme a oublié en tombant dans un corps. Dans le \emph{Ménon}, Socrate fait démontrer à un jeune esclave ignorant une propriété géométrique, simplement en le questionnant : preuve que la vérité était déjà en lui.

L'âme elle-même n'est pas simple : elle comporte \textbf{trois parties} :

\begin{itemize}
\item la partie \cle{rationnelle} (\emph{logistikon}) : la raison, qui connaît et doit commander ; son siège est la tête ;
\item la partie \cle{courageuse} (\emph{thumos}) : le cœur, l'ardeur, l'indignation noble ; son siège est la poitrine ;
\item la partie \cle{désirante} (\emph{epithumia}) : les appétits (faim, soif, plaisirs, richesse) ; son siège est le ventre.
\end{itemize}

L'homme juste est celui dont la raison gouverne, aidée du courage, les désirs — comme un bon conducteur de pirogue sur le Wouri dirige son embarcation au lieu de se laisser porter par le courant.

\begin{attention}[Piège fréquent : « Platon méprise le corps »]
Il est tentant de dire que Platon hait le corps. Nuance obligatoire au Bac : dans le \emph{Phédon}, le corps est bien présenté comme la « tombe » de l'âme (jeu de mots grec \emph{sôma/sêma} : le corps est un tombeau) et comme un obstacle à la connaissance. Mais dans le \emph{Timée} et les \emph{Lois}, le corps est un instrument nécessaire, œuvre du démiurge, qu'il faut entretenir par la gymnastique. Platon condamne la \emph{tyrannie} du corps sur l'âme, non le corps lui-même.
\end{attention}

\courssub{L'allégorie de la Caverne : éducation, vérité et devoir politique}

Au livre VII de la \emph{République}, Platon condense toute sa philosophie dans un récit : l'\cle{allégorie de la Caverne}. Des prisonniers sont enchaînés depuis l'enfance au fond d'une caverne, le dos tourné à l'entrée. Derrière eux brûle un feu ; entre le feu et eux, des hommes font passer des objets. Les prisonniers ne voient que les \emph{ombres} projetées sur la paroi, et ils les prennent pour la réalité. L'un d'eux est délivré : il se retourne, découvre le feu et les objets, puis sort, ébloui, vers la lumière du soleil. Enfin, il redescend dans la caverne pour libérer ses compagnons — qui se moquent de lui et menacent de le tuer (allusion au procès de Socrate).

\begin{popfigure}
\begin{center}
\begin{tikzpicture}[scale=1, every node/.style={font=\small}]
% Caverne
\draw[thick, popdark, fill=popcream] (0,0) -- (0,2.6) -- (3.2,2.6) -- (3.2,0) -- cycle;
\draw[thick, popdark] (3.2,2.6) -- (5.2,3.4);
\draw[thick, popdark] (3.2,0) -- (5.2,0.6);
% Prisonniers
\foreach \x in {0.5,1.1,1.7} {
  \draw[fill=popmuted] (\x,0.35) circle (0.12);
  \draw[popmuted, thick] (\x,0.23) -- (\x,0.05);
}
\node[below, align=center] at (1.1,-0.05) {Prisonniers\\ enchaînés};
% Mur et ombres
\draw[popdark] (2.3,0) -- (2.3,1.1);
\node[above] at (2.3,1.1) {paroi};
\node[align=center, popmuted] at (2.75,2.1) {ombres\\ (images)};
% Feu
\node[poporange] at (2.9,0.75) {\textbf{feu}};
% Chemin de sortie
\draw[->, thick, popblue] (3.4,1.6) .. controls (4.2,2.6) .. (5.4,3.1);
\node[popblue, align=center] at (4.6,2.2) {ascension\\ difficile};
% Soleil
\draw[popgold, fill=popgoldL, thick] (6.6,3.6) circle (0.45);
\foreach \a in {0,45,...,315} \draw[popgold, thick] ({6.6+0.55*cos(\a)},{3.6+0.55*sin(\a)}) -- ({6.6+0.8*cos(\a)},{3.6+0.8*sin(\a)});
\node[popgold, below, align=center] at (6.6,3.0) {Soleil = Idée du Bien};
% Retour
\draw[->, thick, popgreen, dashed] (6.2,3.0) .. controls (5.0,1.2) .. (3.5,0.5);
\node[popgreen, align=center] at (5.6,1.2) {retour :\\ devoir politique};
\end{tikzpicture}
\end{center}
\legende{Schéma de la Caverne : des ombres (opinions) au Soleil (le Bien), puis la redescende du philosophe vers ses concitoyens.}
\end{popfigure}

Chaque élément du récit est symbolique : la caverne est le monde sensible ; les chaînes sont les préjugés et les passions ; les ombres sont les opinions ; la sortie est l'\cle{éducation} (\emph{paideia}), conversion progressive de l'âme ; le soleil est l'Idée du Bien ; le retour est le \emph{devoir} du philosophe de gouverner la cité, même au péril de sa vie.

\begin{methode}[Analyser l'allégorie de la Caverne (explication de texte)]
\begin{enumerate}
\item \textbf{Repérer les quatre stades de la connaissance} en les reliant à la Ligne divisée : les ombres = \emph{eikasia} (imagination) ; les objets portés = \emph{pistis} (croyance) ; la sortie progressive et l'accoutumance = \emph{dianoia} (pensée discursive) ; la contemplation du soleil = \emph{noèsis} (intelligence des Idées).
\item \textbf{Articuler les trois dimensions} : ontologique (les deux mondes), épistémologique (les degrés du savoir), politique (le retour du philosophe dans la cité).
\item \textbf{Expliquer la descente} : pourquoi le libéré revient-il ? Par devoir envers la cité qui l'a formé : la connaissance n'est pas un privilège, elle fonde une responsabilité. Le philosophe doit gouverner, non par goût du pouvoir, mais parce que nul autre ne voit clair.
\end{enumerate}
\end{methode}

\courssub{Politique : la Cité juste et les quatre vertus}

Pour Platon, la cité est l'âme « écrite en grandes lettres » : elle a la même structure tripartite. La cité juste, la \cle{Kallipolis} (« belle cité »), comprend donc \textbf{trois classes}, chacune correspondant à une partie de l'âme, à une vertu propre et — selon le « mythe des métaux » — à un métal mêlé par les dieux à sa nature :

\begin{center}
\begin{tabularx}{\linewidth}{|l|X|X|l|l|}\hline
\rowcolor{popblueL} \textbf{Classe} & \textbf{Partie de l'âme} & \textbf{Fonction} & \textbf{Vertu} & \textbf{Métal} \\ \hline
Gardiens-philosophes & Raison (\emph{logistikon}) & gouverner, délibérer & Sagesse & Or \\ \hline
Auxiliaires-guerriers & Courage (\emph{thumos}) & défendre la cité & Courage & Argent \\ \hline
Producteurs (artisans, paysans, commerçants) & Désir (\emph{epithumia}) & nourrir, produire & Tempérance & Fer et bronze \\ \hline
\end{tabularx}
\end{center}

La quatrième vertu, la \cle{Justice}, n'appartient à aucune classe en particulier : elle est l'\emph{harmonie} de l'ensemble. Sa formule : « \textbf{chacun fait son œuvre} » — le cordonnier fait des chaussures, le guerrier combat, le philosophe gouverne, et nul n'empiète sur la fonction d'autrui. L'injustice, au contraire, est le mélange des fonctions : un producteur riche qui voudrait gouverner corrompt la cité comme les désirs corrompent l'âme.

\begin{exoresolu}[Appliquer Platon à une situation camerounaise]
\textbf{Énoncé.} Dans un lycée technique de Douala, le conseil d'établissement confie la gestion du budget à l'élève le plus populaire, élu pour son charisme, bien qu'il soit faible en mathématiques. Les comptes deviennent vite désordonnés. Analysez cette situation à la lumière de la justice platonicienne.
\tcblower
\textbf{Solution détaillée.} Selon Platon, la justice consiste en ce que « chacun fait son œuvre », c'est-à-dire la fonction pour laquelle sa nature le destine. Or ici, la fonction de gestion (qui exige la raison, la compétence calculatrice) a été attribuée selon le désir de la foule (la popularité), et non selon la compétence rationnelle. C'est exactement ce que Platon reproche à la démocratie athénienne de son temps : confier le gouvernement au hasard du vote plutôt qu'au savoir. Le désordre des comptes illustre l'injustice entendue comme \emph{mélange des fonctions}. La solution platonicienne serait de confier le budget à l'élève formé et compétent (le « gardien » rationnel), la représentation festive au charismatique, chacun retrouvant sa place : l'harmonie, c'est-à-dire la justice, serait restaurée.
\end{exoresolu}

\courssub{Complément pour le Bac : la critique d'Aristote}

Aristote fut vingt ans le disciple de Platon à l'Académie, avant de s'en émanciper. Sa formule restée célèbre : « \emph{amis de Platon, mais plus amis de la vérité} ». Sa critique majeure vise la \cle{séparation} (\emph{chorismos}) des Idées : en plaçant l'essence des choses dans un autre monde, Platon \emph{double inutilement} le réel sans rien expliquer. Dire que ce cheval est beau « parce qu'il participe de la Beauté » n'explique rien : c'est ajouter une entité mystérieuse au lieu de décrire la réalité. Pour Aristote, la forme n'est pas séparée : elle est \emph{dans} les choses (la forme du cheval est dans chaque cheval). Cette discussion entre Platon et Aristote — les essences sont-elles séparées ou immanentes ? — traversera toute l'histoire de la philosophie.

\begin{exemplebox}[Un repère chiffré pour les épreuves]
La \emph{République} compte \textbf{10 livres}. Le livre VII contient l'allégorie de la Caverne. Le livre X développe la critique de l'art : le tableau d'un lit est la copie d'un lit de menuisier, lui-même copie de l'Idée de lit — l'art est donc une « mimèsis de mimèsis », au \textbf{troisième degré} d'éloignement de la vérité. Retenir ces références précises valorise une copie de dissertation.
\end{exemplebox}

\begin{savaistu}[L'Académie, neuf siècles d'école]
Platon fonde vers 387 av. J.-C. l'\cle{Académie}, école installée dans un jardin public d'Athènes. Elle fonctionnera près de \textbf{900 ans}, jusqu'à sa fermeture par l'empereur Justinien en 529 apr. J.-C. À son fronton, une inscription légendaire : « Que nul n'entre ici s'il n'est géomètre. » Signification : les mathématiques, qui habituent l'âme à raisonner sur des objets intelligibles et non sensibles, sont la \emph{propédeutique} (la préparation) indispensable à la dialectique, science suprême des Idées.
\end{savaistu}

\begin{aretenir}[L'essentiel de la section]
\begin{itemize}
\item Platon \emph{ synthétise} Socrate (définitions morales), Héraclite (sensible en devenir) et Parménide (intelligible immuable) dans la théorie des \cle{Idées}.
\item Deux mondes : l'\cle{intelligible} (un, éternel, cause) et le \cle{sensible} (multiple, changeant, copie) ; le lien est la \cle{participation}. Le \cle{Bien} est l'Idée suprême.
\item L'\cle{âme} est immortelle et tripartite (raison, courage, désir) ; connaître est se ressouvenir (\cle{anamnèse}).
\item La \cle{Caverne} figure l'éducation comme conversion et fonde le devoir politique du philosophe.
\item La cité juste (\cle{Kallipolis}) : trois classes, quatre vertus ; la justice = « chacun fait son œuvre ».
\item Aristote critique la \cle{séparation} des Idées : « amis de Platon, mais plus amis de la vérité ».
\end{itemize}
\end{aretenir}

\coursec{Synthèse : Problématique de l'Unité et du Multiple – Application aux situations camerounaises}

Dans la section précédente, nous avons vu comment Platon achève le mouvement de la philosophie antique en construisant un système complet : le monde sensible, multiple et changeant, \emph{participe} d'un monde intelligible, un et stable, celui des Idées. Mais Platon n'a pas inventé ce problème : il en est l'héritier. Depuis les temples de l'Égypte pharaonique jusqu'à l'Académie d'Athènes, toute la philosophie antique peut se lire comme une série de réponses à une même question : \textbf{comment penser ensemble l'Un et le Multiple ?} C'est cette synthèse transversale que nous allons maintenant construire, avant de la mettre à l'épreuve de situations concrètes camerounaises.

\begin{definition}[Unité et Multiple]
L'\cle{Unité} est le principe de cohérence, d'identité et de rassemblement : ce qui fait qu'une réalité tient ensemble et demeure elle-même. Le \cle{Multiple} est le principe de différence, de dispersion et de particularité : la diversité des êtres, des opinions et des changements. La philosophie ne choisit pas l'un contre l'autre : elle \emph{pense leur articulation}, c'est-à-dire la manière dont l'unité peut se retrouver dans la diversité, et la diversité s'ordonner dans l'unité.
\end{definition}

\courssub{Reprise transversale : comment les Anciens ont-ils résolu le problème ?}

\textbf{Intuition.} Observez un fleuve comme le Sanaga : l'eau n'est jamais la même, les berges changent, les crues succèdent aux étiages — c'est le Multiple. Pourtant nous disons « le Sanaga », comme s'il s'agissait d'une seule et même réalité — c'est l'Un. De même, un village réunit des familles rivales, des langues parfois différentes, des générations opposées, et pourtant il reste « un » village. Le problème de l'Un et du Multiple n'est pas abstrait : il est la structure même de notre expérience.

\textbf{Énoncé rigoureux.} Chaque penseur étudié propose une \emph{solution} différente à ce problème :

\begin{itemize}
\item \textbf{La pensée pharaonique} résout le problème par le haut : \cle{Maât}, ordre cosmique, politique et moral à la fois, est le principe d'unité qui rassemble le cosmos, la société et les conduites individuelles. Le Pharaon, garant de Maât, incarne politiquement cette unité ; le désordre (\emph{Isfet}) est le multiple chaotique à contenir.
\item \textbf{Héraclite} affirme que tout coule (\emph{panta rhei}) : le Multiple est premier, tout change. Mais ce changement obéit à une loi universelle, le \cle{Logos} : « l'unité naît de la tension des contraires », comme l'arc et la lyre. L'Un est \emph{dans} le Multiple, non au-delà.
\item \textbf{Parménide} tranche radicalement : seul l'\cle{Être} est, un, immobile, éternel. Le changement et la multiplicité ne sont que des illusions des sens, des opinions (\emph{doxa}). L'Un écrase le Multiple.
\item \textbf{Socrate} déplace le problème vers l'humain : derrière la multiplicité des opinions morales (chacun a « sa » justice, « son » courage), il cherche par le dialogue la \emph{définition universelle} : une seule et même Vertu, identique à la Science du bien.
\item \textbf{Platon} hiérarchise : le Multiple sensible \emph{participe} de l'Un intelligible. Les Idées sont les modèles uniques ; les choses et les opinions en sont les copies multiples. Dans la Cité, la diversité des fonctions est ordonnée par la justice, chacun à sa place.
\end{itemize}

\begin{center}
\begin{tabularx}{\linewidth}{|l|X|X|X|}
\hline
\rowcolor{popblueL} \textbf{Pensée} & \textbf{Où est l'Un ?} & \textbf{Où est le Multiple ?} & \textbf{Solution} \\
\hline
Pharaonique (Maât) & Ordre cosmique et politique incarné par Pharaon & Désordre, chaos (Isfet) & Unité politico-cosmique imposée et maintenue \\
\hline
Héraclite & Logos, loi commune du devenir & Flux perpétuel, contraires & Unité \emph{dans} le changement (harmonie des tensions) \\
\hline
Parménide & Être unique, immobile & Illusion sensible, doxa & Négation du Multiple au profit de l'Un \\
\hline
Socrate & Définition universelle de la Vertu & Opinions particulières & Unité de la Vertu et de la Science par le dialogue \\
\hline
Platon & Idées, monde intelligible & Choses sensibles, copies & Hiérarchie : participation du Multiple à l'Un \\
\hline
\end{tabularx}
\end{center}

\begin{attention}[Deux pièges symétriques]
Deux erreurs guettent l'élève (et le citoyen). Le \textbf{relativisme facile} : « chacun sa vérité, tout se vaut » — c'est dissoudre l'Un dans le Multiple, renoncer à toute recherche d'universalité. Le \textbf{dogmatisme rigide} : « une seule vérité, imposée à tous » — c'est écraser le Multiple sous un Un abstrait. La philosophie antique, dans sa grandeur, cherche l'\emph{universel dans le particulier} : ni dispersion ni écrasement, mais articulation.
\end{attention}

\begin{popfigure}
\begin{center}
\begin{center}\tcbox[colback=popdark,colframe=popdark,coltext=white,arc=9pt,boxsep=4pt,left=6pt,right=6pt]{\bfseries\small Problème Un / Multiple}\end{center}
\vspace{-2pt}
\begin{tcbraster}[raster columns=3,raster equal height=rows,raster column skip=6pt,raster row skip=6pt,enhanced,blankest]
\begin{tcolorbox}[enhanced,colback=white,colframe=popgold,boxrule=0.9pt,arc=3pt,title={Pharaon / Maât},fonttitle=\bfseries\scriptsize,coltitle=white,colbacktitle=popgold,left=4pt,right=4pt,top=2pt,bottom=2pt,boxsep=1pt,toptitle=1.5pt,bottomtitle=1.5pt,halign=left]
\begin{itemize}[nosep,leftmargin=*,itemsep=1pt,font=\scriptsize]
\item Ordre cosmique
\item Unité politique
\end{itemize}
\end{tcolorbox}
\begin{tcolorbox}[enhanced,colback=white,colframe=poporange,boxrule=0.9pt,arc=3pt,title={Héraclite},fonttitle=\bfseries\scriptsize,coltitle=white,colbacktitle=poporange,left=4pt,right=4pt,top=2pt,bottom=2pt,boxsep=1pt,toptitle=1.5pt,bottomtitle=1.5pt,halign=left]
\begin{itemize}[nosep,leftmargin=*,itemsep=1pt,font=\scriptsize]
\item Logos
\item Devenir, contraires
\end{itemize}
\end{tcolorbox}
\begin{tcolorbox}[enhanced,colback=white,colframe=popblue,boxrule=0.9pt,arc=3pt,title={Parménide},fonttitle=\bfseries\scriptsize,coltitle=white,colbacktitle=popblue,left=4pt,right=4pt,top=2pt,bottom=2pt,boxsep=1pt,toptitle=1.5pt,bottomtitle=1.5pt,halign=left]
\begin{itemize}[nosep,leftmargin=*,itemsep=1pt,font=\scriptsize]
\item Être un, immobile
\item Doxa = illusion
\end{itemize}
\end{tcolorbox}
\begin{tcolorbox}[enhanced,colback=white,colframe=popgreen,boxrule=0.9pt,arc=3pt,title={Socrate},fonttitle=\bfseries\scriptsize,coltitle=white,colbacktitle=popgreen,left=4pt,right=4pt,top=2pt,bottom=2pt,boxsep=1pt,toptitle=1.5pt,bottomtitle=1.5pt,halign=left]
\begin{itemize}[nosep,leftmargin=*,itemsep=1pt,font=\scriptsize]
\item Dialogue, maïeutique
\item Vertu = Science
\end{itemize}
\end{tcolorbox}
\begin{tcolorbox}[enhanced,colback=white,colframe=poppurple,boxrule=0.9pt,arc=3pt,title={Platon},fonttitle=\bfseries\scriptsize,coltitle=white,colbacktitle=poppurple,left=4pt,right=4pt,top=2pt,bottom=2pt,boxsep=1pt,toptitle=1.5pt,bottomtitle=1.5pt,halign=left]
\begin{itemize}[nosep,leftmargin=*,itemsep=1pt,font=\scriptsize]
\item Idées, participation
\item Cité hiérarchisée
\end{itemize}
\end{tcolorbox}
\end{tcbraster}

\end{center}
\legende{Carte heuristique centrale de la leçon : le problème de l'Un et du Multiple et ses cinq réponses antiques. Chaque branche colorée correspond à un auteur, chaque feuille à un concept-clé.}
\end{popfigure}

\begin{popfigure}
\begin{center}
\begin{tikzpicture}[scale=1.05]
% Grille radar 5 axes
\foreach \r in {1,2,3,4} {
  \draw[popline] (90:\r) -- (162:\r) -- (234:\r) -- (306:\r) -- (18:\r) -- cycle;
}
\foreach \a/\t in {90/{Stabilité},162/{Changement},234/{Raison},306/{Politique},18/{Éthique}} {
  \draw[popline] (0,0) -- (\a:4);
  \node[font=\scriptsize\bfseries, popdark] at (\a:4.55) {\t};
}
% Maât (or) : stabilité 4, changement 1, raison 2, politique 4, éthique 3
\draw[thick, popgold, fill=popgold!25, fill opacity=0.5] (90:4) -- (162:1) -- (234:2) -- (306:4) -- (18:3) -- cycle;
% Héraclite (orange) : 1,4,3,2,2
\draw[thick, poporange, fill=poporange!20, fill opacity=0.5] (90:1) -- (162:4) -- (234:3) -- (306:2) -- (18:2) -- cycle;
% Parménide (bleu) : 4,0.3,4,1,1
\draw[thick, popblue, fill=popblue!15, fill opacity=0.5] (90:4) -- (162:0.3) -- (234:4) -- (306:1) -- (18:1) -- cycle;
% Socrate (vert) : 2,2,4,2,4
\draw[thick, popgreen, fill=popgreen!15, fill opacity=0.5] (90:2) -- (162:2) -- (234:4) -- (306:2) -- (18:4) -- cycle;
% Platon (violet) : 4,1,4,4,4
\draw[thick, poppurple, fill=poppurple!15, fill opacity=0.5] (90:4) -- (162:1) -- (234:4) -- (306:4) -- (18:4) -- cycle;
% Légende
\node[font=\scriptsize, anchor=west] at (4.6,2.2) {\textcolor{popgold}{\rule{0.4cm}{0.15cm}} Maât};
\node[font=\scriptsize, anchor=west] at (4.6,1.7) {\textcolor{poporange}{\rule{0.4cm}{0.15cm}} Héraclite};
\node[font=\scriptsize, anchor=west] at (4.6,1.2) {\textcolor{popblue}{\rule{0.4cm}{0.15cm}} Parménide};
\node[font=\scriptsize, anchor=west] at (4.6,0.7) {\textcolor{popgreen}{\rule{0.4cm}{0.15cm}} Socrate};
\node[font=\scriptsize, anchor=west] at (4.6,0.2) {\textcolor{poppurple}{\rule{0.4cm}{0.15cm}} Platon};
\node[font=\scriptsize, popmuted] at (0,-4.6) {Échelle de 0 (centre) à 4 (bord) : intensité accordée à chaque axe.};
\end{tikzpicture}
\end{center}
\legende{Diagramme radar comparatif des cinq pensées antiques sur cinq axes : Stabilité, Changement, Raison, Politique, Éthique. Parménide et Platon privilégient la stabilité ; Héraclite est la seule pensée où le changement domine ; Socrate concentre tout sur la raison éthique.}
\end{popfigure}

\courssub{Application 1 : L'identité nationale camerounaise}

\textbf{La situation.} La devise du Cameroun proclame : « Un peuple, un but, une foi ». Pourtant le pays compte plus de 250 groupes ethniques, deux langues officielles (français et anglais), des traditions religieuses plurielles (christianisme, islam, religions traditionnelles). L'indice de fragmentation ethnolinguistique du Cameroun est d'environ \num{0,89} — l'un des plus élevés au monde : deux personnes prises au hasard ont près de neuf chances sur dix d'appartenir à des groupes linguistiques différents.

\begin{exemplebox}[Lecture philosophique du défi national]
Le défi politique camerounais est exactement le problème antique : construire l'\textbf{Unité} (la Nation) sans nier le \textbf{Multiple} (les communautés). Deux modèles s'affrontent :
\begin{itemize}
\item Le \textbf{modèle héraclitéen} : l'unité naît de la tension harmonieuse des différences. Le bilinguisme, la diversité culturelle ne sont pas des obstacles mais des « contraires accordés », comme l'arc et la lyre. L'unité est \emph{dans} la diversité.
\item Le \textbf{modèle parménidien-platonicien} : l'unité suppose un modèle unique auquel les particularités doivent participer — une langue, une loi, une identité uniformes. Risque : nier le Multiple réel au nom d'un Un abstrait.
\end{itemize}
La sagesse politique consiste à articuler les deux : des institutions communes (Un) qui reconnaissent et organisent la diversité (Multiple).
\end{exemplebox}

\begin{savaistu}[Le problème philosophique inscrit dans la loi]
Le Préambule de la Constitution camerounaise de 1996 affirme l'attachement du peuple à l'\emph{unité dans la diversité}. Cette formule est la traduction politique directe du problème philosophique de l'Un et du Multiple : le législateur camerounais, sans le savoir peut-être, rédige une réponse à Héraclite et à Platon.
\end{savaistu}

\courssub{Application 2 : Justice coutumière et justice moderne}

\textbf{La situation.} Dans un village de l'Ouest, un conflit foncier oppose deux familles. Deux voies sont possibles : la \emph{palabre} devant le chef et les notables, qui cherche la réconciliation et le consensus ; ou le tribunal d'instance, qui applique le Code pénal et le droit foncier, identiques pour tous.

\textbf{L'analyse.} La \cle{palabre} fonctionne comme une recherche d'unité \emph{relationnelle} : il ne s'agit pas d'appliquer une règle abstraite, mais de restaurer l'harmonie du groupe — proche de Maât (l'ordre à maintenir) et du dialogue socratique (la vérité naît de la parole échangée). Le \cle{droit moderne}, lui, ressemble au modèle platonicien : une loi universelle, abstraite, identique pour tous, comme une Idée de la justice à laquelle les cas particuliers doivent se conformer.

\textbf{Limites croisées.} La palabre peut enfermer dans le particulier (le groupe prime sur le droit de l'individu, par exemple la veuve déshéritée) ; le droit abstrait peut ignorer le particulier (jugement juste en droit mais qui détruit la paix sociale). L'articulation Un/Multiple exige que la loi universelle laisse une place aux médiations coutumières, et que celles-ci respectent les droits fondamentaux.

\courssub{Application 3 : Environnement et forêt du bassin du Congo}

\textbf{La situation.} La forêt camerounaise, partie du bassin du Congo, est exploitée (bois, agriculture) et protégée (aires protégées, gestion communautaire). Comment penser notre rapport à la nature ?

\textbf{L'analyse.} Deux visions s'opposent. La \textbf{vision cyclique et animiste}, proche d'Héraclite et de la pensée pharaonique : la nature est un tout vivant, ordonné, en équilibre ; l'humain est \emph{dans} le cosmos, soumis à ses cycles (saisons, crues), et les traditions (forêts sacrées, interdits de chasse) maintiennent l'harmonie — une Maât écologique. La \textbf{vision rationnelle et quantifiée}, héritière du Logos et du Noûs : la nature est un objet de science, mesurable, gérable — quotas d'exploitation, inventaires forestiers, plans d'aménagement.

\textbf{Articulation.} La gestion durable exige les deux : la rigueur du Logos (données, lois environnementales) et la sagesse de l'appartenance au tout (savoirs locaux, respect des cycles). Nier le Multiple vivant au nom d'un calcul unique mène à l'épuisement des ressources ; refuser toute rationalité mène à l'impuissance.

\begin{methode}[Appliquer une notion philosophique à une situation concrète]
Savoir-faire exigible au Probatoire et au Baccalauréat technique. Trois étapes :
\begin{enumerate}
\item \textbf{Identifier le concept philosophique} : par exemple, « la justice comme harmonie des parties » (Platon) ou « l'unité dans la tension des contraires » (Héraclite).
\item \textbf{Décrire précisément la situation réelle} : par exemple, un conflit foncier entre deux familles, avec ses acteurs, ses enjeux, ses règles en présence.
\item \textbf{Tester la pertinence} : le concept éclaire-t-il réellement la situation ? Que permet-il de comprendre ? Quelles sont ses \emph{limites} (ce qu'il ne permet pas de voir) ?
\end{enumerate}
Une bonne application n'est jamais un collage : elle montre à la fois la fécondité et les bornes du concept.
\end{methode}

\begin{exoresolu}[Dissertation courte : « Peut-on faire l'unité sans nier la diversité ? »]
Énoncé : À partir des pensées antiques étudiées et de l'exemple camerounais, rédigez l'introduction et le plan d'une dissertation sur le sujet : « Peut-on faire l'unité sans nier la diversité ? »
\tcblower
\textbf{Introduction.} Amorce : le Cameroun proclame « Un peuple, un but, une foi » tout en comptant plus de 250 groupes ethniques. Problème : l'unité semble exiger l'uniformité ; la diversité semble menacer l'unité. Problématique : l'unité véritable est-elle l'uniformité, ou peut-elle s'accomplir \emph{dans} la diversité ? Annonce du plan.

\textbf{Plan.}
\begin{enumerate}
\item \emph{Thèse} : l'unité paraît exiger l'effacement des différences. Parménide : seul l'Un est vrai, le multiple est illusion. Platon : les choses multiples ne valent que par leur participation à l'Idée unique. Application : la tentation de l'uniformisation (une seule langue, un seul modèle).
\item \emph{Antithèse} : mais une unité qui nie la diversité est morte ou violente. Héraclite : l'harmonie naît de la tension des contraires, comme l'arc et la lyre ; supprimer le multiple, c'est supprimer l'unité vivante elle-même. Application : une nation qui écraserait ses cultures perdrait sa richesse et ferait naître les conflits.
\item \emph{Synthèse} : l'unité véritable est \emph{articulation}, non écrasement. Maât : l'ordre rassemble sans abolir les fonctions diverses. Socrate : l'universel se cherche \emph{dans} le dialogue des opinions particulières. Application : « l'unité dans la diversité » constitutionnelle, le bilinguisme, la décentralisation comme médiations.
\end{enumerate}
\textbf{Conclusion attendue} : oui, on peut et on doit faire l'unité sans nier la diversité, à condition de penser l'unité comme harmonie ordonnée du multiple, et non comme uniformité.
\end{exoresolu}

\begin{popfigure}
\begin{center}
\begin{tikzpicture}[font=\scriptsize, box/.style={draw, rounded corners=2pt, align=center, inner sep=3pt}]
% Cité de Platon (gauche)
\node[box, fill=poppurpleL, text width=3.1cm] (p1) at (0,3.4) {\textbf{Philosophes-rois}\\ (savoir, Idées)};
\node[box, fill=poppurple!40, text width=3.1cm] (p2) at (0,2.1) {\textbf{Gardiens}\\ (courage, défense)};
\node[box, fill=poppurple!20, text width=3.1cm] (p3) at (0,0.8) {\textbf{Producteurs}\\ (travail, besoins)};
\node[font=\small\bfseries, poppurple] at (0,4.4) {Cité de Platon};
% État camerounais (droite)
\node[box, fill=popgreenL, text width=3.6cm] (c1) at (7,3.4) {\textbf{Pouvoir central}\\ Président, Gouvernement};
\node[box, fill=popgreen!40, text width=3.6cm] (c2) at (7,2.1) {\textbf{Administration}\\ Gouverneurs, forces de l'ordre};
\node[box, fill=popgreen!20, text width=3.6cm] (c3) at (7,0.8) {\textbf{Collectivités décentralisées}\\ Régions, communes, citoyens};
\node[font=\small\bfseries, popgreen!60!black] at (7,4.4) {État camerounais unitaire décentralisé};
% Flèches d'analogie
\draw[<->, thick, popmuted] (p1) -- (c1);
\draw[<->, thick, popmuted] (p2) -- (c2);
\draw[<->, thick, popmuted] (p3) -- (c3);
\node[text width=9.5cm, align=center, font=\scriptsize\itshape, popmuted] at (3.5,-0.2) {Analogies : hiérarchie des fonctions, recherche d'unité ordonnée. Limites : chez Platon, la place est fixée par nature ; dans l'État moderne, les citoyens participent et les fonctions ne sont pas héréditaires.};
\end{tikzpicture}
\end{center}
\legende{Comparaison schématique entre la Cité platonicienne (trois classes ordonnées par la justice) et l'État camerounais unitaire décentralisé. L'analogie éclaire l'idée d'unité hiérarchisée ; les limites rappellent que la démocratie moderne refuse la fixation naturelle des places.}
\end{popfigure}

\courssub{TP d'intégration : la carte heuristique géante}

\begin{experience}[Construction collective d'une carte mentale au tableau]
\textbf{Matériel} : tableau, craies ou feutres de cinq couleurs, fiches de synthèse des sections 1 à 4.
\textbf{Protocole} :
\begin{enumerate}
\item Tracer au centre du tableau un ovale : « Problème Un / Multiple ».
\item Cinq groupes d'élèves tracent chacun une branche colorée : Pharaon/Maât, Héraclite, Parménide, Socrate, Platon.
\item Chaque groupe ajoute deux sous-branches : les \emph{concepts-clés} de son auteur (ex. : Logos, devenir pour Héraclite).
\item Chaque groupe accroche enfin une \emph{feuille-exemple camerounais} : devise nationale, palabre, forêt communautaire, bilinguisme, décentralisation.
\item La classe observe la carte complète et formule oralement une phrase-synthèse : « Pour les Anciens, l'unité\ldots ».
\end{enumerate}
\textbf{Observations attendues} : les cinq branches répondent au même centre ; aucune ne supprime totalement l'Un ou le Multiple, sauf Parménide (cas limite). \textbf{Interprétation} : la philosophie progresse par déplacement du problème, non par réponses définitives.
\end{experience}

\begin{aretenir}[Boîte à outils pour la dissertation]
Mots-clés à mobiliser pour tout sujet touchant à l'Un et au Multiple : \cle{articulation}, \cle{dépassement}, \cle{hiérarchie}, \cle{participation}, \cle{médiation}, \cle{dialectique}. Exemple d'emploi : « Platon pense la \emph{participation} du multiple sensible à l'Idée une ; la politique camerounaise cherche une \emph{médiation} entre unité nationale et diversité communautaire ; la pensée héraclitéenne est déjà \emph{dialectique}, car elle fait naître l'harmonie de la tension des contraires. »
\end{aretenir}

\begin{exercice}[Pour s'entraîner]
1. Définissez l'Unité et le Multiple, puis illustrez chaque notion par un exemple tiré de la vie de votre établissement scolaire.
2. En dix lignes, montrez comment la palabre villageoise et le tribunal moderne incarnent deux solutions différentes au problème de l'Un et du Multiple.
3. Rédigez la problématique et l'annonce de plan d'une dissertation sur le sujet : « La diversité culturelle est-elle un obstacle à l'unité nationale ? »
\end{exercice}

\begin{corrige}[Corrigé des exercices d'entraînement]
1. L'Unité est le principe de cohérence et de rassemblement ; le Multiple, le principe de différence et de dispersion. Exemple : le lycée est \emph{un} (même règlement, même établissement) mais \emph{multiple} (classes, spécialités techniques, langues des élèves).
2. La palabre recherche l'unité \emph{dans} la relation : par le dialogue, elle restaure l'harmonie concrète du groupe (proche de Maât et de Socrate). Le tribunal applique une loi universelle et abstraite, identique pour tous les cas particuliers (proche du modèle platonicien de l'Idée). La première risque le particularisme, le second l'indifférence au concret : l'idéal est leur articulation.
3. Problématique attendue : la diversité semble fragmenter, l'unité semble exiger l'uniformité ; mais l'unité véritable n'est-elle pas plutôt l'harmonie ordonnée des différences ? Plan possible : I. La diversité comme menace (Parménide, uniformisation) ; II. La diversité comme richesse (Héraclite, tension harmonieuse) ; III. L'unité comme articulation (Maât, Platon, « unité dans la diversité » constitutionnelle).
\end{corrige}

\coursec{Atelier méthodologique : Rédiger et justifier une démarche philosophique (Dissertation \& Explication)}

Dans la section précédente, nous avons synthétisé l'héritage de la philosophie antique — de Maât à Platon — autour du problème de l'\cle{Unité} et du \cle{Multiple}, et nous l'avons appliqué à des situations camerounaises concrètes (cohésion nationale, diversité culturelle, éducation). Mais au Probatoire et au Baccalauréat technique, il ne suffit pas de \emph{savoir} : il faut savoir \emph{écrire} sa pensée. Le correcteur ne lit pas dans votre esprit ; il évalue un texte. Cette dernière section est donc un \textbf{atelier} : nous allons apprendre à construire, rédiger et surtout \emph{justifier} une démarche philosophique complète, dans les deux épreuves au programme : la \cle{dissertation} et l'\cle{explication de texte}.

\courssub{Les attendus officiels de l'épreuve écrite}

\begin{definition}[Les deux exercices au programme]
La \cle{dissertation} est un exercice de réflexion personnelle et argumentée à partir d'un sujet (une question ou une affirmation à discuter). L'\cle{explication de texte} est un commentaire méthodique d'un extrait d'un auteur étudié dans l'unité (ici : Platon, Héraclite, Parménide, ou un texte de la sagesse pharaonique), dont il faut dégager la thèse, la structure et la portée.
\end{definition}

\begin{exemplebox}[Barème type de l'épreuve]
Le barème indicatif se répartit classiquement ainsi :
\begin{itemize}
\item \textbf{Introduction et conclusion : 4 points} (problématisation, annonce du plan, bilan).
\item \textbf{Développement / analyse : 12 points} (arguments, exemples, articulation logique, justification de la démarche).
\item \textbf{Expression et orthographe : 4 points} (clarté, connecteurs logiques, correction de la langue).
\end{itemize}
La \emph{justification de la démarche} — le fait d'expliquer pourquoi vous choisissez tel argument ou tel exemple — est évaluée dans la rubrique « Argumentation ».
\end{exemplebox}

\courssub{La dissertation : de l'analyse du sujet à la rédaction}

\textbf{Étape 1 : analyser le sujet.} Tout commence par une lecture active. On souligne les mots-clés, on les définit rigoureusement, puis on cherche la \emph{tension} qu'ils entretiennent. Un sujet de philosophie n'est jamais une simple demande d'information : il cache un \cle{problème}, c'est-à-dire une difficulté rationnelle qui rend la réponse non évidente. Prenons le sujet d'entraînement de cette leçon : \emph{« L'unité est-elle l'ennemie de la diversité ? »}. Les termes à définir sont « unité » (ce qui fait qu'un tout est un) et « diversité » (la multiplicité des éléments). La tension apparaît immédiatement : unifier, est-ce nécessairement uniformiser et donc détruire la diversité ? Ou au contraire, l'unité peut-elle \emph{protéger} la diversité ? C'est exactement le problème que Parménide (l'Être est Un) et Héraclite (l'harmonie naît des contraires) ont posé il y a vingt-cinq siècles.

\textbf{Étape 2 : construire le plan.} Deux structures sont admises :
\begin{itemize}
\item Le \cle{plan dialectique} (thèse / antithèse / synthèse) : on examine d'abord une position, puis la position opposée, puis on dépasse l'opposition par une distinction conceptuelle ou un changement de point de vue.
\item Le \cle{plan analytique} : on suit l'ordre des concepts (constat / causes / conséquences, ou analyse des notions du sujet les unes après les autres).
\end{itemize}

\begin{popfigure}
\begin{center}
\begin{tikzpicture}[
node distance=1.2cm and 2cm,
etape/.style={draw=popblue, fill=popblueL, rounded corners, text width=5cm, align=center, font=\small, thick},
fleche/.style={-{Stealth[length=3mm]}, thick, popdark}]
\node[etape, fill=popgoldL, draw=popgold, align=center] (sujet){\textbf{SUJET}\\ Analyse des termes, définitions};
\node[etape, below=of sujet, align=center] (pb){\textbf{PROBLÈME}\\ Tension, paradoxe, question directrice};
\node[etape, below left=of pb, fill=popgreenL, draw=popgreen, align=center] (these){\textbf{THÈSE}\\ Arguments + exemples};
\node[etape, below right=of pb, fill=poporangeL, draw=poporange, align=center] (anti){\textbf{ANTITHÈSE}\\ Arguments + exemples};
\node[etape, below=2.8cm of pb, fill=poppurpleL, draw=poppurple, align=center] (syn){\textbf{SYNTHÈSE}\\ Dépassement : distinction conceptuelle};
\node[etape, below=of syn, fill=popcream, draw=popmuted, align=center] (conc){\textbf{CONCLUSION}\\ Bilan + ouverture};
\draw[fleche] (sujet) -- (pb);
\draw[fleche] (pb) -- (these);
\draw[fleche] (pb) -- (anti);
\draw[fleche] (these) -- (syn);
\draw[fleche] (anti) -- (syn);
\draw[fleche] (syn) -- (conc);
\end{tikzpicture}
\end{center}
\legende{Organigramme du plan dialectique : du sujet à la conclusion, chaque étape enchaîne logiquement avec la suivante.}
\end{popfigure}

\textbf{Étape 3 : rédiger.} La rédaction suit trois moments : l'introduction, le développement, la conclusion.

\begin{methode}[L'introduction « en entonnoir » : quatre temps]
\begin{enumerate}
\item \textbf{L'accroche} : une citation, un fait d'actualité ou une référence culturelle camerounaise (ex. : la devise « Paix-Travail-Patrie », la cohabitation de plus de 250 groupes ethniques au Cameroun).
\item \textbf{La définition rigoureuse des termes} du sujet : chaque mot-clé reçoit un sens précis, éventuellement appuyé sur un auteur au programme.
\item \textbf{La problématisation} : on formule la tension ou le paradoxe qui rend le sujet problématique, puis la question directrice.
\item \textbf{L'annonce du plan} avec des verbes d'action : « Nous examinerons d'abord..., nous montrerons ensuite..., avant de dépasser cette opposition en... ».
\end{enumerate}
L'entonnoir va donc du large (l'accroche) vers l'étroit (le problème précis).
\end{methode}

Le \textbf{développement} est composé de parties elles-mêmes divisées en paragraphes. Chaque paragraphe suit le schéma : \emph{idée — argument — exemple — analyse de l'exemple — transition}. La \textbf{conclusion} fait le bilan de la réponse apportée au problème et ouvre sur une question nouvelle (jamais une répétition sèche de l'introduction).

\begin{attention}[L'erreur fatale : le « catalogage d'auteurs »]
Écrire « Platon dit ceci, Héraclite dit cela, Parménide pense cela » sans articulation argumentative n'est pas philosopher : c'est réciter. Le candidat qui réussit \emph{pense avec} les auteurs : il les mobilise comme des ressources au service de \emph{sa} démonstration. Chaque référence doit être introduite (« Comme le montre Platon dans l'allégorie de la caverne... »), expliquée, puis reliée à votre propre thèse (« Cela montre que... »).
\end{attention}

\begin{propriete}[Réussir la synthèse]
La synthèse n'est pas un compromis « mi-figue mi-raisin » (« l'unité est un peu ennemie, un peu amie de la diversité »). Elle opère :
\begin{itemize}
\item soit une \textbf{distinction conceptuelle} : par exemple distinguer l'\emph{unité formelle} (cadre commun, loi, langue partagée) de l'\emph{uniformité réelle} (écrasement des différences) ;
\item soit un \textbf{changement de niveau} : par exemple passer du plan ontologique (l'Être de Parménide) au plan politique (la cité juste de Platon, où l'unité harmonise des fonctions diverses).
\end{itemize}
\end{propriete}

\courssub{L'explication de texte : méthode en quatre temps}

L'explication de texte obéit à une structure différente : il ne s'agit plus de discuter librement une question, mais de faire parler un texte.

\begin{enumerate}
\item \textbf{Présentation} : situer l'auteur, l'œuvre, le contexte de l'extrait, et annoncer la \emph{thèse} du texte en une phrase.
\item \textbf{Analyse structurée} : découper le texte en séquences logiques et les expliquer \emph{dans l'ordre} : idée directrice, arguments, concepts, exemples de l'auteur. On suit le fil du texte, on ne le quitte jamais longtemps.
\item \textbf{Discussion critique} : une fois le texte compris, on évalue sa portée, ses limites éventuelles et son actualité (par exemple : que vaut le mythe de la caverne à l'ère des réseaux sociaux et des « cavernes numériques » ?).
\item \textbf{Conclusion} : bilan de ce que le texte nous a appris et place de l'extrait dans l'œuvre ou dans l'histoire de la philosophie.
\end{enumerate}

\begin{popfigure}
\begin{center}
\begin{tikzpicture}[
node distance=0.4cm,
bloc/.style={draw=popdark, rounded corners, text width=11.5cm, align=left, font=\small, thick},
fleche/.style={-{Stealth[length=2.5mm]}, thick, popmuted}]
\node[bloc, fill=popgoldL, draw=popgold] (tete) {\textbf{En-tête} : Auteur — Œuvre — Contexte — Thèse annoncée};
\node[bloc, below=of tete, fill=popblueL, draw=popblue] (intro) {\textbf{Introduction (4 temps)} : accroche / présentation du texte / problème posé par le texte / annonce du plan d'analyse};
\node[bloc, below=of intro, fill=popgreenL, draw=popgreen] (corps) {\textbf{Corps} : analyse des séquences logiques dans l'ordre du texte (idée directrice, arguments, concepts, exemples de l'auteur)};
\node[bloc, below=of corps, fill=poporangeL, draw=poporange] (disc) {\textbf{Discussion} : portée, limites, actualité de la thèse};
\node[bloc, below=of disc, fill=popcream, draw=popmuted] (conc) {\textbf{Conclusion} : bilan + place de l'extrait dans l'œuvre};
\draw[fleche] (tete) -- (intro);
\draw[fleche] (intro) -- (corps);
\draw[fleche] (corps) -- (disc);
\draw[fleche] (disc) -- (conc);
\end{tikzpicture}
\end{center}
\legende{Structure type de l'explication de texte : l'analyse suit fidèlement l'ordre des séquences du texte avant la discussion critique.}
\end{popfigure}

\begin{savaistu}[Un atout majeur : le contexte exact]
Au Baccalauréat technique, l'explication de texte porte fréquemment sur un extrait de \emph{Platon} (l'allégorie de la caverne, la ligne divisée, les mythes politiques) ou sur des \emph{fragments} d'Héraclite et de Parménide. Connaître le contexte précis de l'extrait — par exemple savoir que la caverne ouvre le livre VII de \emph{La République} et illustre la conversion de l'âme vers le Bien — permet de présenter le texte avec autorité et de gagner des points dès l'introduction.
\end{savaistu}

\courssub{Justifier sa démarche : la méta-cognition philosophique}

Le savoir-faire « rédiger et \emph{justifier} une démarche » demande un effort de \cle{méta-cognition} : penser sa propre pensée. Concrètement, cela signifie expliciter vos choix :
\begin{itemize}
\item « \emph{Je structure mon devoir en plan dialectique car} le sujet oppose deux termes (unité / diversité) qu'il faut d'abord confronter avant de les réconcilier. »
\item « \emph{Cet exemple camerounais illustre} la thèse d'Héraclite : le bilinguisme officiel français-anglais montre que l'unité nationale peut naître de la tension des contraires, et non de leur suppression. »
\item « \emph{Je choisis l'allégorie de la caverne ici parce qu'elle} donne un fondement concret à l'idée que l'éducation unifie sans uniformiser : chaque prisonnier gravit le même chemin, mais depuis sa propre position. »
\end{itemize}

\begin{aretenir}[Boîte à outils : les connecteurs logiques obligatoires]
Un devoir philosophique est un raisonnement continu. Employez : \emph{En effet} (justification), \emph{Cependant} / \emph{Or} (objection), \emph{Par conséquent} / \emph{Dès lors} (conséquence), \emph{Cela montre que} (analyse d'exemple), \emph{Au-delà de cette opposition...} (transition vers la synthèse), \emph{En définitive} (conclusion).
\end{aretenir}

\courssub{Entraînement guidé : « L'unité est-elle l'ennemie de la diversité ? »}

\begin{exoresolu}[Sujet type de dissertation — corrigé guidé]
\textbf{Énoncé.} En travail de groupe : (a) problématisez le sujet « L'unité est-elle l'ennemie de la diversité ? » ; (b) construisez un plan dialectique détaillé ; (c) rédigez l'introduction et la conclusion.
\tcblower
\textbf{(a) Problématisation.} L'unité désigne ce qui fait qu'un ensemble tient ensemble ; la diversité, la multiplicité des éléments qui le composent. Le sens commun oppose les deux : unifier semble exiger d'uniformiser, donc de niveler les différences. Mais l'expérience camerounaise — un État, plus de 250 langues — suggère l'inverse : sans unité politique, la diversité éclate ; sans diversité, l'unité devient tyrannie. D'où le problème : \emph{l'unité détruit-elle nécessairement la diversité, ou peut-elle au contraire la fonder et la protéger ?}

\textbf{(b) Plan dialectique détaillé.}
\begin{itemize}
\item \textbf{Thèse — L'unité menace la diversité.} Argument 1 : unifier, c'est soumettre le multiple à une loi unique (Parménide : l'Être est Un, le multiple est illusion). Argument 2 : les politiques d'uniformisation (langue unique, pensée unique) étouffent les identités. Exemple : les nationalismes exclusifs.
\item \textbf{Antithèse — Sans unité, la diversité se détruit elle-même.} Argument 1 : Héraclite — l'harmonie naît de la tension des contraires, comme l'arc et la lyre ; sans cadre unifiant, les différences s'affrontent. Argument 2 : Platon — la cité juste harmonise des fonctions diverses sous l'idée du Bien. Exemple : la devise « Paix-Travail-Patrie » comme cadre commun à des cultures multiples.
\item \textbf{Synthèse — Distinguer unité formelle et uniformité réelle.} L'unité est ennemie de la diversité seulement lorsqu'elle se confond avec l'uniformité. Une unité \emph{formelle} (loi, langue partagée, valeurs communes) protège la diversité \emph{réelle} des contenus. L'unité véritable est harmonie, non identité.
\end{itemize}

\textbf{(c) Introduction rédigée (modèle).} « Le Cameroun se présente souvent comme "l'Afrique en miniature" : un seul État, mais plus de deux cent cinquante groupes ethniques et deux langues officielles. Cette situation résume une énigme que la philosophie antique a posée dès ses origines : comment l'un et le multiple peuvent-ils coexister ? L'\emph{unité} désigne ce qui fait qu'un ensemble forme un tout ; la \emph{diversité}, la multiplicité des éléments qui le composent. Or unifier ne semble possible qu'en uniformisant, donc en sacrifiant les différences : l'unité serait l'ennemie de la diversité. Mais inversement, une diversité sans cadre commun ne risque-t-elle pas l'éclatement ? L'unité est-elle donc nécessairement l'ennemie de la diversité ? Nous examinerons d'abord la thèse selon laquelle toute unité écrase le multiple ; nous montrerons ensuite que, sans unité, la diversité se détruit elle-même ; enfin, nous dépasserons cette opposition en distinguant l'unité formelle de l'uniformité réelle. »

\textbf{Conclusion rédigée (modèle).} « En définitive, l'unité n'est l'ennemie de la diversité que lorsqu'elle dégénère en uniformité. Comme Héraclite l'avait pressenti, l'accord le plus beau naît des contraires : une nation, comme une lyre, ne produit d'harmonie que de cordes différentes tendues vers un même but. Reste à demander qui doit accorder ces cordes : quelle place l'éducation, que Platon plaçait au sommet de la cité, doit-elle tenir dans la construction de cette unité harmonieuse ? »
\end{exoresolu}

\courssub{Auto-évaluation : la grille MINESEC}

Après chaque devoir, évaluez-vous honnêtement sur les six critères officiels, chacun noté sur 20 : \textbf{Compréhension} du sujet ou du texte, \textbf{Problématisation}, \textbf{Argumentation} (incluant la justification de la démarche), \textbf{Structuration}, \textbf{Expression}, \textbf{Exemples}. Reportez vos notes sur le radar ci-dessous : un profil équilibré est un grand hexagone ; un « creux » désigne votre priorité de travail.

\begin{popfigure}
\begin{center}
\begin{tikzpicture}[scale=1.15]
\foreach \r in {1,2,3,4} {\draw[popline, thin] (0,0) circle (\r);}
\foreach \a/\nom in {90/Compréhension, 30/Problématisation, -30/Argumentation, -90/Structuration, -150/Expression, 150/Exemples} {
\draw[popline, thin] (0,0) -- (\a:4.4);
\node[font=\footnotesize, popdark] at (\a:5.15) {\nom};}
\draw[popblue, very thick, fill=popblue, fill opacity=0.18]
(90:3.2) -- (30:2.6) -- (-30:3.4) -- (-90:2.8) -- (-150:3.0) -- (150:2.4) -- cycle;
\foreach \a/\r in {90/3.2, 30/2.6, -30/3.4, -90/2.8, -150/3.0, 150/2.4} {\fill[popblue] (\a:\r) circle (0.06);}
\end{tikzpicture}
\end{center}
\legende{Grille radar d'auto-évaluation (exemple de profil d'élève) : chaque cercle représente 5 points (rayon 1 = 5, rayon 4 = 20). Reproduisez ce radar sur votre copie de brouillon et reportez vos propres notes après chaque entraînement.}
\end{popfigure}

\begin{exercice}[À vous de philosopher]
\begin{enumerate}
\item Analysez le sujet : « Faut-il sortir de la caverne ? » (définitions, problématisation, plan dialectique).
\item Justifiez par écrit, en trois phrases, le choix de votre plan et d'un exemple camerounais.
\item Rédigez l'introduction complète du sujet, puis auto-évaluez-la avec la grille à six critères.
\end{enumerate}
\end{exercice}

\begin{corrige}[Exercice — éléments de correction]
\textbf{1.} « Caverne » renvoie à l'allégorie de Platon (\emph{La République}, VII) : condition d'ignorance où l'on prend les ombres pour la réalité ; « faut-il » interroge le devoir et la valeur de la connaissance. Problème : la vérité est-elle préférable au confort de l'illusion, alors que la sortie est douloureuse et que le philosophe libéré est rejeté par les prisonniers ? Plan : thèse (il faut sortir : la vérité libère et fonde la cité juste) ; antithèse (la sortie coûte : douleur, solitude, risque — certains préfèrent les chaînes) ; synthèse (sortir est un devoir, mais graduel et éducatif : la conversion de l'âme exige une pédagogie).
\textbf{2.} Exemple de justification : « Je choisis le plan dialectique car le sujet oppose le devoir de vérité au prix de la sortie ; j'utilise l'exemple des élèves camerounais qui découvrent l'esprit critique à l'école, car il illustre concrètement la conversion progressive décrite par Platon. »
\textbf{3.} L'introduction doit suivre les quatre temps de l'entonnoir ; l'auto-évaluation doit identifier au moins un critère faible et un moyen précis de progresser.
\end{corrige}

\begin{aretenir}[L'essentiel de l'atelier]
\begin{itemize}
\item Dissertation : analyse du sujet $\rightarrow$ problématique $\rightarrow$ plan (dialectique ou analytique) $\rightarrow$ introduction en entonnoir, développement argumenté, conclusion avec ouverture.
\item Explication de texte : présentation $\rightarrow$ analyse séquence par séquence $\rightarrow$ discussion critique $\rightarrow$ conclusion.
\item Toujours \emph{justifier} ses choix (méta-cognition) et relier chaque auteur à sa propre démonstration.
\item S'auto-évaluer sur les six critères MINESEC pour progresser à chaque devoir.
\end{itemize}
\end{aretenir}

\newpage

\coursec{Activité d'intégration}

\coursec{Leçon 16 : La philosophie antique}

\courssub{Objectifs de l'activité}
Cette activité d'intégration vise à mobiliser l'ensemble des savoirs (Maât, Héraclite, Parménide, Socrate, Platon) et savoir-faire (argumentation, application concrète, rédaction philosophique) dans une situation de délibération démocratique simulée. Elle permet à l'élève de :
\begin{itemize}
    \item \textbf{Problématiser} la tension entre unité et diversité à l'épreuve du contexte camerounais.
    \item \textbf{Articuler} des concepts philosophiques anciens (Logos, Idées, Maïeutique, Justice, Maât) à une question civique actuelle.
    \item \textbf{S'exercer} à la dissertation philosophique (choix du plan, thèse/antithèse/synthèse, exemples pertinents) et à l'explication de texte (justification de la démarche).
    \item \textbf{Développer} l'esprit critique, la tolérance et la rigueur argumentative attendus au Probatoire.
\end{itemize}

\courssub{Situation}
\emph{Contexte :} Le Ministère des Enseignements Secondaires (MINESEC) lance, pour la rentrée 2026-2027, une consultation nationale sur la refonte du programme d'\textbf{Éducation à la Citoyenneté} en classe de Première Technique. L'enjeu : définir le « socle commun » qui doit unir les jeunes Camerounais au-delà de leurs différences (plus de 250 groupes ethniques, deux langues officielles, diversités religieuses et coutumières).

Deux propositions officielles s'affrontent dans un \textbf{Livre Blanc} diffusé dans les 58 départements :

\begin{center}
\begin{tabularx}{\linewidth}{|l|X|X|}
\hline
\rowcolor{popblueL}
\textbf{Critère} & \textbf{Proposition A : « Socle Unique et Rigide »} & \textbf{Proposition B : « Socle Commun Minimal + Décentralisation Culturelle »} \\
\hline
\emph{Inspiration philosophique} & Parménide (L'Être est, immuable) / Platon (Idées parfaites, Gardiens) & Héraclite (Tout coule) / Maât (Équilibre dynamique) / Socrate (Dialogue, Maïeutique) \\
\hline
\emph{Contenu enseigné} & \begin{itemize}\item Hymne, drapeau, devise identiques.\item Histoire nationale unique (version officielle).\item Langue unique d'enseignement (français/anglais selon zone).\item Valeurs « non négociables » listées par le centre.\end{itemize} & \begin{itemize}\item Loi, Constitution, symboles républicains (communs).\item Place aux langues nationales (Bassa, Bamiléké, Fulfulde, etc.).\item Intégration de la justice coutumière.\item Modules « Cultures locales » obligatoires.\item Forums de dialogue intercommunautaires.\end{itemize} \\
\hline
\emph{Objectif affiché} & Unité formelle, stabilité, prévention du tribalisme, efficacité administrative. & Unité vivante, adaptabilité, reconnaissance de la diversité, légitimité par le consensus. \\
\hline
\emph{Risque identifié} & Rigidité, exclusion des minorités, ressentiment, fracture identitaire. & Fragmentation, relativisme, difficulté d'évaluation nationale, instrumentalisation politique. \\
\hline
\end{tabularx}
\end{center}
\par\medskip\noindent\textit{Tableau comparatif des deux propositions du MINESEC pour le nouveau cours d'Éducation à la Citoyenneté.}\par\medskip

Vous êtes délégués d'une \textbf{Cité École} (lycée technique pilote de 1~200 élèves, représentant 10 régions). Le Proviseur organise un \textbf{Grand Débat} sur 3 séances pour rendre un avis motivé au Ministre.

\courssub{Consignes}
\textbf{Séance 1 — Préparation en groupes (4-5 élèves) — 50 min}
\begin{enumerate}
    \item \textbf{Tirage au sort d'un « Rôle » :} Philosophes antiques (experts concepts) ; Chefs traditionnels (gardiens de la coutume) ; Jeunes urbains (réseaux sociaux, mobilité) ; Parents d'élèves (transmission intergénérationnelle) ; Administration (MINESEC, inspection).
    \item \textbf{Élaboration d'une « Fiche Argumentaire » (1 page A4) :} Chaque groupe doit produire 3 arguments \textbf{philosophiquement fondés} (citation ou concept précis du cours) pour défendre sa position (A ou B) ou proposer une synthèse. La fiche mentionne : \emph{Concept mobilisé $\rightarrow$ Auteur $\rightarrow$ Application au Cameroun $\rightarrow$ Argument}.
    \item \textbf{Désignation d'un orateur} et d'un secrétaire pour le compte-rendu.
\end{enumerate}

\textbf{Séance 2 — Débat contradictoire (Assemblée) — 50 min}
\begin{enumerate}
    \item \textbf{Présentation (5 min/groupe) :} L'orateur expose la position du groupe en s'appuyant sur la fiche. Le modérateur (Professeur) veille au respect du temps et à la pertinence philosophique.
    \item \textbf{Échanges libres (15 min) :} Questions croisées, objections, demandes de clarification. Le modérateur note au tableau les concepts-clés mobilisés (Maât, Logos, Idées, Maïeutique, Justice, Unité/Multiple).
    \item \textbf{Vote à main levée :} Pour A, Pour B, Pour une synthèse (majorité qualifiée des 2/3 requise pour la synthèse).
    \item \textbf{Synthèse orale du Professeur (5 min) :} Relevé des arguments forts, failles logiques, concepts mal utilisés.
\end{enumerate}

\textbf{Séance 3 — Rédaction individuelle (Évaluation sommative) — 50 min}
\begin{enumerate}
    \item \textbf{Sujet de dissertation :} \emph{« L'unité nationale requiert-elle l'uniformité culturelle ? »}
    \item \textbf{Consignes de rédaction :}
    \begin{itemize}
        \item Analyser les termes du sujet (unité, uniformité, culture, nation).
        \item Problématiser : opposition apparente vs tension féconde.
        \item Construire un plan dialectique (Thèse / Antithèse / Synthèse) ou analytique (Concepts / Acteurs / Enjeux) \textbf{justifié en introduction}.
        \item Mobiliser \textbf{au moins 3 auteurs} du cours (Maât, Héraclite, Parménide, Socrate, Platon) et \textbf{2 exemples camerounais concrets} (ex. : bilinguisme, justice coutumière, fête de l'unité, langues nationales, décentralisation).
        \item Conclure en rouvrant sur la question du « vivre-ensemble » aujourd'hui.
        \item \textbf{Expliciter la démarche} en bas de copie (10 lignes) : choix du plan, sélection des exemples, difficulté rencontrée, auteur le plus utile.
    \end{itemize}
\end{enumerate}

\courssub{Ressources à mobiliser}
\begin{itemize}
    \item \cle{Maât (Égypte pharaonique)} : Ordre cosmique et social, justice comme équilibre dynamique, \emph{« Maât est immuable, mais son application s'adapte aux circonstances »}. Unité par l'harmonie, non par l'uniformité.
    \item \cle{Héraclite (Éphèse, v. 500 av. J.-C.)} : \emph{« On ne se baigne pas deux fois dans le même fleuve »} (Fragment 12). \cle{Logos} : loi universelle du devenir, unité des contraires (\emph{« La voie qui monte et la voie qui descend sont une et la même »}). L'unité est tension, mouvement, feu toujours vivant.
    \item \cle{Parménide (Élée, v. 515 av. J.-C.)} : \emph{« L'Être est, le non-être n'est pas »}. L'Être est Un, immuable, indivisible, éternel. Le changement est illusion des sens. Modèle d'unité \textbf{formelle, statique, excluant le multiple}.
    \item \cle{Socrate (Athènes, 470-399 av. J.-C.)} : \cle{Maïeutique} (accouchement des esprits par le dialogue), \cle{« Connais-toi toi-même »} (Delphes), refus des sophistes (relativisme). L'unité naît de l'examen critique partagé, non de l'imposition. \emph{« Nul ne fait le mal volontairement »} (ignorance = source du désordre).
    \item \cle{Platon (Athènes, 428-348 av. J.-C.)} : \cle{Théorie des Idées} (monde intelligible parfait vs monde sensible changeant). \cle{Allégorie de la Caverne} (éducation = ascension vers la vérité). \cle{République} : Gardiens-Philosophes détenteurs de la science du Bien, éducation unique, censure, communauté de biens/femmes pour l'unité de la Cité. Unité \textbf{hiérarchique, rationnelle, imposée d'en haut}.
\end{itemize}

\courssub{Démarche et corrigé commenté}
\begin{exoresolu}[Dissertation : « L'unité nationale requiert-elle l'uniformité culturelle ? »]
\textbf{1. Analyse du sujet et Problématisation}
\begin{itemize}
    \item \textbf{Unité nationale :} lien politique et juridique unissant les citoyens d'un État (Cameroun : Constitution, bilinguisme, institutions).
    \item \textbf{Uniformité culturelle :} identité des mœurs, langues, valeurs, histoires (modèle jacobin, assimilationniste).
    \item \textbf{Requiert :} condition nécessaire ? suffisante ? simple moyen ?
    \item \textbf{Problème :} L'uniformité garantit-elle l'unité ou la tue-t-elle en niant le multiple ? L'unité peut-elle être \emph{vivante} (Héraclite/Maât) sans être \emph{figée} (Parménide/Platon) ?
\end{itemize}

\textbf{2. Choix du plan (Justification) : Plan dialectique}
\emph{Thèse : L'uniformité culturelle semble nécessaire à l'unité (Parménide/Platon / Proposition A).}
\emph{Antithèse : L'uniformité détruit l'unité vivante ; l'unité suppose la diversité gérée (Héraclite/Maât/Socrate / Proposition B).}
\emph{Synthèse : L'unité nationale est un \textbf{projet dynamique} : socle juridico-politique commun (lois, symboles) + reconnaissance institutionnalisée de la diversité (décentralisation, dialogue).}

\textbf{3. Développement détaillé}

\textbf{Thèse : L'uniformité comme condition de l'unité (Modèle Parménide/Platon).}\par
\begin{itemize}
    \item \textbf{Parménide :} L'Être est Un. Le multiple, le changement sont illusion (\emph{doxa}). Transposé : la Nation doit être Une, indivisible. Toute différence culturelle est source de division (tribalisme). \emph{Exemple :} L'histoire officielle unique enseignée partout (programme unique) forge une mémoire commune, évite les « histoires rivales ».
    \item \textbf{Platon (\emph{République}) :} Les Gardiens connaissent l'Idée du Bien. L'éducation est unique, censurée, pour formater les âmes à l'unité de la Cité. \emph{Exemple camerounais :} L'hymne national, le drapeau, la devise « Paix - Travail - Patrie » sont des \emph{Idées} intangibles, identiques de Kousséri à Campo. Le bilinguisme officiel (français/anglais) impose deux langues \emph{uniformes} pour l'unité administrative.
    \item \textbf{Limite :} Unité \emph{formelle}, \emph{statique}, \emph{imposée}. Risque de \emph{tyrannie de la majorité} (anglophones vs francophones), d'exclusion des minorités (Pygmées Baka, Mbororo), de rupture avec le \emph{Logos} vivant des sociétés.
\end{itemize}

\textbf{Antithèse : L'unité vivante naît de la diversité reconnue (Modèle Héraclite/Maât/Socrate).}\par
\begin{itemize}
    \item \textbf{Héraclite :} L'unité est \emph{logos} : tension des contraires (\emph{« La discorde est le père de toutes choses »}). Un fleuve reste le même \emph{en changeant}. \emph{Exemple :} Le Cameroun comme « Afrique en miniature » : sa force est sa diversité (250+ ethnies). Vouloir l'uniformiser (ex. interdire les langues nationales à l'école) tarit le fleuve, crée de la frustration (crise anglophone 2016-).
    \item \textbf{Maât (Pharaon) :} Pharaon garantit l'ordre \emph{en} maintenant l'équilibre entre les noms (fonctions), les régions, les cultes locaux. Justice = \emph{redresser ce qui est courbé}, pas aplatir. \emph{Exemple :} La \textbf{justice coutumière} (tribunaux traditionnels, chefferies) reconnue par la Constitution (Art. 37) : elle gère les conflits fonciers/familiaux selon les coutumes locales \emph{dans} le cadre de la loi nationale. Unité par \emph{subsidiarité}.
    \item \textbf{Socrate :} La vérité naît du \emph{dialogue} (maïeutique), non de l'imposition. \emph{Exemple :} Les \textbf{forums de dialogue intercommunautaire} (ex. Grand Dialogue National 2019, comités de vigilance villageois) : l'unité se \emph{construit} par la parole partagée, le compromis, l'écoute. L'uniformité tue le dialogue (pensée unique).
\end{itemize}

\textbf{Synthèse : Unité nationale = « Unité dans la diversité » (Projet politique éthique).}\par
\begin{itemize}
    \item \textbf{Socle dur (Parménide/Platon) :} Constitution, lois, institutions, symboles républicains, service national, éducation civique de base (droits/devoirs). C'est l'Être immuable de la République. \emph{Non-négociable.}
    \item \textbf{Espace fluide (Héraclite/Maât/Socrate) :} Langues nationales à l'école (module optionnel obligatoire), gestion décentralisée des affaires coutumières, fête de l'unité (20 mai) célébrée \emph{avec} les danses/tenues régionales, médias publics multilingues.
    \item \textbf{Institutionnalisation du dialogue (Socrate) :} Conseil national de la décentralisation, observatoire des langues, « Parlement des jeunes » : lieux où l'unité se \emph{réinvente} sans cesse par le conflit réglé par la loi et la parole.
    \item \textbf{Conclusion :} L'unité nationale ne \emph{requiert} pas l'uniformité culturelle (c'est une confusion catégorie : politique vs anthropologique). Elle \emph{exige} un \textbf{contrat partagé} (lois, valeurs républicaines) qui \textbf{protège} la diversité comme richesse (Maât) et organise le \textbf{conflit permanent} (Héraclite) par le \textbf{dialogue raisonné} (Socrate). L'uniformité est une \emph{fausse unité} (mortelle) ; la vraie unité est \emph{vivante} (féconde).
\end{itemize}

\textbf{4. Explicitation de la démarche (modèle pour l'élève)}
\emph{« J'ai choisi le plan dialectique car le sujet oppose deux modèles philosophiques clairs (Parménide/Platon vs Héraclite/Maât/Socrate) qui correspondent aux deux propositions du débat. J'ai sélectionné l'exemple de la justice coutumière (Maât) et de la crise anglophone (Héraclite) car ils sont au cœur de l'actualité camerounaise et illustrent concrètement les concepts. La difficulté a été de ne pas faire un simple catalogue d'auteurs mais de les faire \emph{dialoguer} dans chaque partie. L'auteur le plus utile pour la synthèse a été Socrate : sa maïeutique offre le \emph{procédé} (le dialogue) pour articuler le socle dur (Platon) et la diversité vivante (Héraclite). »}
\tcblower
\textbf{Commentaire du correcteur (grille Probatoire) :}
\begin{itemize}
    \item \textbf{Connaissances (6 pts) :} Concepts précis (Maât, Logos, Idées, Maïeutique), auteurs nommés, citations approximatives acceptées. \checkmark
    \item \textbf{Argumentation (8 pts) :} Thèse/Antithèse/Synthèse claires, exemples camerounais pertinents (justice coutumière, bilinguisme, Grand Dialogue), articulation concept/exemple réussie. \checkmark
    \item \textbf{Problématisation (3 pts) :} Tension unité/uniformité bien posée, termes définis. \checkmark
    \item \textbf{Rédaction/Démarche (3 pts) :} Plan justifié, explicitation de la démarche en fin de copie (méta-cognition), style correct. \checkmark
    \item \textbf{Total : 20/20} (copie type « Très Bien »).
\end{itemize}
\end{exoresolu}

\courssub{Bilan de l'activité}
Cette activité a permis de mettre en évidence :
\begin{itemize}
    \item \textbf{L'opérativité des concepts antiques :} Maât, Logos, Idées, Maïeutique ne sont pas des « morceaux de musée » mais des \textbf{outils de pensée} pour analyser les tensions actuelles (décentralisation, bilinguisme, justice coutumière, vivre-ensemble).
    \item \textbf{La nécessité du dialogue philosophique :} Le passage du « choc des opinions » (doxa) au « conflit des arguments » (logos) structure le débat démocratique. L'élève a vécu la \emph{maïeutique} : ses idées ont été « accouchées » par la contradiction.
    \item \textbf{L'ancrage camerounais comme révélateur :} L'exemple de la justice coutumière (Art. 37 Constitution) ou de la crise anglophone montre que la philosophie antique éclaire des choix juridiques et politiques réels.
    \item \textbf{L'exigence de la rédaction :} L'explicitation de la démarche (méta-cognition) transforme l'exercice scolaire en compétence transférable : \emph{savoir comment on a pensé} est aussi important que \emph{ce qu'on a pensé}.
    \item \textbf{La continuité Maât $\rightarrow$ Héraclite $\rightarrow$ Socrate $\rightarrow$ Platon :} L'élève perçoit l'histoire de la philosophie non comme une suite de noms mais comme une \emph{réponse progressive} à la question : \emph{Comment l'Un (État, Nation, Vérité) pense-t-il le Multiple (cultures, opinions, devenir) ?}
\end{itemize}
\emph{Compétence certifiée :} « Rédiger et justifier une démarche philosophique en mobilisant des concepts antiques sur une situation citoyenne camerounaise. »

\newpage

\coursec{Méthodes et exercices résolus}

\courssub{Méthodes et exercices résolus}

\begin{methode}[Analyser un concept philosophique antique : la méthode du triangle conceptuel]
\textbf{Objectif :} Maîtriser la définition, l'étymologie, les contraires et les enjeux d'une notion clé (ex. : Maât, Logos, Idée, Arete).
\begin{enumerate}
    \item \textbf{Étymologie et sens premier :} Retracer l'origine du mot (grec, égyptien) et son sens littéral.
    \item \textbf{Définition philosophique :} Préciser le sens technique chez l'auteur étudié (différent du sens courant).
    \item \textbf{Champ conceptuel :} Identifier les notions voisines (synonymes), les contraires (antonymes) et les distinctions nécessaires (ex. : \emph{Être} vs \emph{Devenir}, \emph{Opinion} vs \emph{Science}).
    \item \textbf{Problématisation :} Formuler la question philosophique que ce concept permet de poser (ex. : « La justice est-elle une convention ou un ordre cosmique ? »).
    \item \textbf{Illustration concrète (Cameroun) :} Donner un exemple actuel ancré dans le réel (coutume, droit, vie scolaire, gestion publique).
\end{enumerate}
\end{methode}

\begin{exoresolu}[Analyse du concept de Maât dans la pensée pharaonique]
\textbf{Énoncé :} « Maât n'est pas seulement une déesse, c'est l'architecture même du monde. » À partir de cette affirmation, analysez le concept de Maât en montrant en quoi il fonde à la fois l'ordre cosmique, l'ordre politique et l'éthique individuelle dans l'Égypte antique. Donnez une illustration camerounaise.
\tcblower
\textbf{Solution détaillée :}

\textbf{1. Étymologie et sens premier.}
Le terme \cle{Maât} (m3ˁt) désigne en égyptien ancien la « vérité », la « justice », l'« ordre », la « droiture », mais aussi la « mesure » et l'« harmonie ». C'est un concept \emph{onto-cosmologique} : il nomme la structure fondamentale de la réalité.

\textbf{2. Définition philosophique (triple dimension).}
\begin{itemize}
    \item \textbf{Cosmique :} Maât est l'ordre stable qui s'oppose au \cle{Isfet} (le chaos, le désordre, l'injustice). Le cycle solaire (lever/coucher), la crue du Nil, le rythme des saisons sont l'actualisation de Maât. Sans Maât, le monde retourne au Nun (eaux primordiales chaotiques).
    \item \textbf{Politique :} Pharaon est le « Seigneur de Maât ». Sa fonction royale n'est pas tyrannique mais \emph{cosmothétique} : il doit « mettre Maât à la place d'Isfet ». La légitimité du pouvoir dépend du respect de cet ordre. Les vizirs portent le titre de « Prêtres de Maât ».
    \item \textbf{Éthique individuelle :} Dans les \emph{Textes des Sarcophages} et le \emph{Livre des Morts} (chapitre 125, « Confession négative »), le défunt doit déclarer : « Je n'ai pas fait tel mal, je n'ai pas volé, je n'ai pas menti ». L'accès à l'éternité (champ de Ialou) suppose une vie conforme à Maât. L'éthique est donc participation à l'ordre du monde.
\end{itemize}

\textbf{3. Champ conceptuel.}
\begin{itemize}
    \item \textbf{Contraire :} \cle{Isfet} (chaos, mensonge, violence, désordre).
    \item \textbf{Notions voisines :} \cle{Maât} (loi/divinité), \cle{Ankh} (vie), \cle{Sa} (protection magique).
    \item \textbf{Distinction :} Contrairement au \emph{Droit positif} moderne (écrit, changeant), Maât est \emph{non-écrite}, immuable, antérieure aux lois humaines. La loi humaine n'est juste que si elle \emph{révèle} Maât.
\end{itemize}

\textbf{4. Problématisation.}
Comment une société peut-elle fonder sa justice sur un principe immuable (Maât) tout en s'adaptant aux contingences historiques ? Le Pharaon est-il au-dessus de Maât ou son serviteur ?

\textbf{5. Illustration camerounaise.}
La notion de \textbf{\cle{Ngondo}} chez les Sawa (peuple de l'eau, Littoral) présente une analogie structurelle forte. Le Ngondo est à la fois :
\begin{itemize}
    \item Un \textbf{esprit de l'eau} (dimension cosmique/sacrée),
    \item Une \textbf{assemblée traditionnelle} de régulation des conflits (dimension politique/juridique),
    \item Un \textbf{code de conduite} (honnêteté, respect de la parole donnée, solidarité) conditionnant l'appartenance à la communauté (dimension éthique).
\end{itemize}
Tout comme Maât, le Ngondo assure la cohésion du corps social en le reliant à un ordre transcendant (les ancêtres, les génies de l'eau). La violation du serment sur le Ngondo (Isfet local) entraîne l'exclusion sociale, comme la non-conformité à Maât exclut de la vie éternelle.

\textbf{Conclusion :} Maât est le \emph{principe d'unité} de la civilisation pharaonique : elle unit le ciel et la terre, le roi et le peuple, l'individu et la communauté, le temps présent et l'éternité. C'est une \emph{philosophie de l'harmonie} antérieure de deux millénaires à la philosophie grecque.
\end{exoresolu}

\begin{methode}[Comparer deux thèses antagonistes : la méthode du face-à-face structuré (Héraclite vs Parménide)]
\textbf{Objectif :} Mettre en évidence l'opposition radicale entre deux systèmes sur un même problème (l'Être et le Devenir) pour en dégager l'apport à l'histoire de la pensée.
\begin{enumerate}
    \item \textbf{Identifier le problème commun :} Quelle est la nature fondamentale de la réalité ? (Unité vs Multiplicité, Permanence vs Changement).
    \item \textbf{Reconstituer la thèse A (Héraclite) :} Citation clé $\rightarrow$ Thèse centrale $\rightarrow$ Arguments (logos, feu, unité des contraires) $\rightarrow$ Conséquences (impossibilité de la connaissance stable, flux universel).
    \item \textbf{Reconstituer la thèse B (Parménide) :} Citation clé $\rightarrow$ Thèse centrale $\rightarrow$ Arguments (raison/parménide, non-être impossible, pensée=être) $\rightarrow$ Conséquences (le changement est illusion, connaissance par la raison seule).
    \item \textbf{Tableau comparatif :} Critères : Principe premier, Rôle des sens, Rôle de la raison, Statut du changement, Unité/Multiple.
    \item \textbf{Synthèse dialectique :} En quoi cette opposition force la philosophie à distinguer \emph{apparence} et \emph{réalité}, \emph{opinion} (\emph{doxa}) et \emph{science} (\emph{epistémè}) ? Préfiguration de Platon.
\end{enumerate}
\end{methode}

\begin{exoresolu}[L'opposition Héraclite / Parménide : le tournant ontologique]
\textbf{Énoncé :} « On ne se baigne jamais deux fois dans le même fleuve » (Héraclite) vs « L'Être est, le non-être n'est pas » (Parménide). Présentez l'opposition entre ces deux conceptions de la réalité. En quoi cette querelle fondatrice oblige-t-elle la philosophie à inventer la distinction entre le sensible et l'intelligible ?
\tcblower
\textbf{Solution détaillée :}

\textbf{1. Problème commun : La permanence dans le changement.}
Les deux philosophes ioniens (Éphèse pour Héraclite, Élée pour Parménide, v. 500 av. J.-C.) cherchent l'\cle{Archè} (principe premier) unique expliquant la multiplicité des choses. Ils refusent le mythe au profit du \cle{Logos} (discours rationnel).

\textbf{2. Thèse d'Héraclite : Le Devenir éternel (Le Flux).}
\begin{itemize}
    \item \textbf{Citation :} « Tout coule » (\emph{Panta rhei}), « Le chemin qui monte et celui qui descend sont un et le même ».
    \item \textbf{Principe :} Le \cle{Feu} (symbole de transformation permanente, mesure d'échange).
    \item \textbf{Thèse centrale :} La réalité est \cle{Devenir}. L'Être fixe n'existe pas. L'identité d'une chose n'est qu'un équilibre instable de forces contraires (\cle{unité des contraires} : vie/mort, veille/sommeil, chaud/froid).
    \item \textbf{Épistémologie :} Les sens perçoivent le flux, mais l'intelligence (comprendre le \cle{Logos} universel) saisit la loi immuable \emph{du} changement. « La nature aime à se cacher ».
\end{itemize}

\textbf{3. Thèse de Parménide : L'Être immuable (Le Plein).}
\begin{itemize}
    \item \textbf{Citation :} « Car il est le même, penser et être » (Fr. 3). « L'Être est ingénéré, impérissable, tout entier, unique, immobile et sans fin » (Fr. 8).
    \item \textbf{Méthode :} La \cle{voie de la vérité} (raison pure) vs \cle{voie de l'opinion} (sens trompeurs).
    \item \textbf{Argument ontologique :} Le Non-Être ne peut \emph{pas} être (contradiction logique). Donc l'Être \emph{ne peut pas ne pas être}. L'Être est \textbf{plein} (pas de vide), \textbf{immobile} (le mouvement suppose du non-être : aller là où on n'est pas), \textbf{ingénéré} (sortir du non-être ? impossible), \textbf{indivisible} (la division suppose du non-être entre les parties).
    \item \textbf{Conséquence :} Le changement, la naissance, la mort, la multiplicité sont des \cle{illusions des sens} (\emph{doxa}). Seule la raison (\emph{noêsis}) accède à la vérité (\emph{alétheia}).
\end{itemize}

\textbf{4. Tableau comparatif.}

\begin{center}
\begin{tabularx}{\linewidth}{|l|X|X|}
\hline
\rowcolor{popblueL}
\textbf{Critère} & \textbf{Héraclite (Éphèse)} & \textbf{Parménide (Élée)} \\
\hline
\textbf{Réalité fondamentale} & Le Devenir / Le Feu / Le Logos & L'Être / Le Plein / La Pensée \\
\hline
\textbf{Statut du changement} & \textbf{Réel, unique réalité} & \textbf{Illusion, apparence} \\
\hline
\textbf{Unité / Multiplicité} & Unité \emph{dans} la multiplicité (unité des contraires) & Unité \emph{excluant} la multiplicité (Un sphérique) \\
\hline
\textbf{Rôle des sens} & Nécessaires mais insuffisants (témoignent du flux) & \textbf{Trompeurs} (source de l'erreur) \\
\hline
\textbf{Rôle de la raison} & Comprendre la loi du changement (Logos) & \textbf{Seul accès à la vérité} (Être = Pensée) \\
\hline
\textbf{Connaissance} & Difficile (les hommes « ne comprennent pas ») & Possible, certaine, par déduction logique \\
\hline
\end{tabularx}
\end{center}

\textbf{5. Synthèse dialectique : L'invention du double monde.}
L'aporie (impasse) est totale : \textbf{Soit} la réalité change (Héraclite) $\rightarrow$ pas de science possible (objet fuyant), \textbf{Soit} la réalité ne change pas (Parménide) $\rightarrow$ le monde sensible est mensonge.
\emph{Réponse platonicienne (anticipée) :} Il faut \textbf{deux niveaux de réalité} :
\begin{itemize}
    \item Le \textbf{Sensible} (monde d'Héraclite) : domaine du Devenir, de l'opinion (\emph{doxa}), des sens.
    \item L'\textbf{Intelligible} (monde de Parménide) : domaine de l'Être immuable, de la science (\emph{epistémè}), des Idées/Formes.
\end{itemize}
Platon \emph{sauve} les deux : le changement est réel \emph{ici-bas}, la permanence est réelle \emph{là-haut}.

\textbf{Conclusion :} Cette opposition structure toute la métaphysique occidentale. Elle force à penser la \cle{participation} (comment le sensible « participe » de l'intelligible ?) et fonde l'exigence de \cle{concept} comme objet stable de la science.
\end{exoresolu}

\begin{methode}[Expliquer un texte philosophique (Explication linéaire) : la méthode en 4 temps]
\textbf{Objectif :} Réussir l'exercice de l'explication de texte (épreuve du Probatoire/Bac) en suivant une progression rigoureuse.
\begin{enumerate}
    \item \textbf{Mise en contexte (Introduction) :} Auteur, œuvre, date, place dans l'œuvre, thèse générale du texte, problématique du passage, annonce du plan (2 ou 3 parties logiques suivant le fil du texte).
    \item \textbf{Explication linéaire (Développement) :} Suivre le texte \emph{paragraphe par paragraphe} (ou idée par idée).
        \begin{itemize}
            \item \textbf{Reformuler} l'idée en français clair (ne pas paraphraser mot à mot).
            \item \textbf{Analyser} les concepts clés, les arguments, la logique (déduction, analogie, réfutation).
            \item \textbf{Éclairer} par des références internes (autres passages de l'auteur) ou externes (auteurs proches/opposés).
        \end{itemize}
    \item \textbf{Synthèse philosophique (Conclusion) :} Reprendre la thèse défendue, montrer la portée/les limites de l'argumentation, ouvrir sur une question prolongement.
    \item \textbf{Relecture croisée :} Vérifier : Ai-je cité le texte ? Ai-je expliqué \emph{tous} les mots techniques ? La structure suit-elle le texte ? Le style est-il impersonnel et précis ?
\end{enumerate}
\end{methode}

\begin{exoresolu}[Explication de texte : L'allégorie de la Caverne (République, VII, 514a-517a)]
\textbf{Énoncé :} Extrait : « Imagine des hommes dans une demeure souterraine... » jusqu'à « ...et qu'il faut l'habituer peu à peu à voir les choses d'en haut. » Expliquez ce texte en montrant comment Platon y figure la condition humaine, le rôle du philosophe et la difficulté de l'éducation.
\tcblower
\textbf{Solution détaillée :}

\textbf{Introduction.}
\textbf{Auteur/Œuvre :} Platon (428-348 av. J.-C.), \emph{La République} (Politeia), Livre VII. Dialogue socratique sur la Justice et l'État idéal.
\textbf{Contexte immédiat :} Socrate répond à Glaucon qui demande de décrire le philosophe-roi. L'allégorie illustre la théorie des Idées et la pédagogie politique.
\textbf{Thèse du texte :} La condition humaine ordinaire est une ignorance radicale (ombre/illusion) ; la philosophie est une ascèse douloureuse (libération/ascension) vers la vérité (Idées/Soleil) ; le philosophe a le devoir moral de redescendre pour gouverner (redescente), malgré l'hostilité des restants.
\textbf{Problématique :} Comment l'allégorie figure-t-elle simultanément une \emph{ontologie} (deux mondes), une \emph{épistémologie} (quatre degrés de connaissance) et une \emph{politique} (devoir du sage) ?
\textbf{Plan :} I. La caverne : image de l'ignorance naturelle (514a-515c). II. L'ascension : le parcours douloureux du savoir (515c-516c). III. Le retour : la charge politique du philosophe (516c-517a).

\textbf{Développement.}

\textbf{I. La caverne : image de l'ignorance naturelle (condition pré-philosophique).}
\begin{itemize}
    \item \textbf{Situation initiale :} « Demeure souterraine », « entrée ouverte à la lumière » (ironie : la lumière est dehors, ils sont dos à elle). « Chaînes aux jambes et au cou » : aliénation physique $\rightarrow$ aliénation mentale. Ils ne peuvent « regarder que devant eux ».
    \item \textbf{Le spectacle des ombres :} « Feu » (source de lumière artificielle, copie du soleil), « Muret » (écran), « Porteurs d'objets » (manipulateurs : sophistes, politiciens, médias d'aujourd'hui). « Ombres » = \cle{Images des choses fabriquées} (deuxième degré d'éloignement de la vérité : copies de copies).
    \item \textbf{Épistémologie :} C'est le degré le plus bas : l'\cle{Imagination (Eikasia)}. Les prisonniers prennent les ombres pour la réalité, les échos pour les voix vraies. « Pour eux, la vérité ne serait littéralement que les ombres des images. » $\rightarrow$ Critique de la \emph{doxa} (opinion) non examinée.
\end{itemize}

\textbf{II. L'ascension : le parcours douloureux du savoir (l'éducation/paideia).}
\begin{itemize}
    \item \textbf{Libération :} « Délivré de ses chaînes », « contraint de se lever », « tourner le cou », « marcher », « lever les yeux vers la lumière ». Verbes d'effort, de contrainte, de douleur (« souffrirait », « éblouissement »). $\rightarrow$ L'éducation n'est pas un « remplissage » mais un \textbf{tournant de l'âme} (\emph{periagogè}), une conversion totale.
    \item \textbf{Étapes de l'adaptation (pédagogie progressive) :}
        \begin{enumerate}
            \item Ombres/Reflets dans l'eau (degré 1 : encore images).
            \item Objets réels / Astres la nuit (Lune, étoiles) (degré 2 : \cle{Foi - Pistis}).
            \item Le \textbf{Soleil} lui-même (degré 3 : \cle{Intelligence / Noêsis}). « Il pourrait contempler le soleil lui-même, non pas dans l'eau... mais en lui-même, à sa place. »
        \end{enumerate}
    \item \textbf{Fonction du Soleil (Idée du Bien) :} 1. Source de \emph{visibilité} (rend les choses connaissables). 2. Source de \emph{génération} (cause de l'être et de la vie). 3. Source de \emph{croissance}. $\rightarrow$ L'Idée du Bien est au-delà de l'Être (\emph{epekeina tès ousias}), cause de la vérité et de la science.
\end{itemize}

\textbf{III. Le retour : la charge politique du philosophe (la justice dans la Cité).}
\begin{itemize}
    \item \textbf{Le devoir de redescendre :} « Ne serait-il pas tenu d'aller partager les travaux et les honneurs... » Le philosophe ne fuit pas le monde ; sa connaissance de l'Idée du Bien l'oblige à servir la Cité. C'est la \cle{justice} : chacun sa fonction, le meilleur pour le commandement.
    \item \textbf{La réadaptation difficile :} « Aveuglement » en redescendant (yeux pleins de lumière divine). Il paraît ridicule aux prisonniers (« yeux gâtés »).
    \item \textbf{L'hostilité mortelle :} « S'ils pouvaient mettre la main sur celui qui essaie de les délier et de les faire monter, ils le tueraient. » $\rightarrow$ Allusion directe au procès de \textbf{Socrate} (399 av. J.-C.). La vérité dérange l'ordre établi de l'opinion.
    \item \textbf{Conclusion du passage :} « Il faut l'habituer peu à peu. » $\rightarrow$ Nécessité d'une \textbf{éducation mathématique/dialectique} progressive (arithmétique, géométrie, astronomie, dialecte) pour former les gardiens (521d-534e).
\end{itemize}

\textbf{Synthèse / Conclusion.}
L'allégorie condense le \textbf{système platonicien} :
\begin{enumerate}
    \item \textbf{Ontologie dualiste :} Deux mondes (Sensible/Intelligible), deux réalités (Devenir/Être).
    \item \textbf{Épistémologie graduée :} 4 degrés (Eikasia, Pistis, Dianoia, Noêsis) $\leftrightarrow$ 4 segments de la Ligne (Rép. VI).
    \item \textbf{Anthropologie :} L'âme a des yeux (l'intelligence) qu'il faut tourner vers le vrai.
    \item \textbf{Politique :} Le savoir n'est pas contemplation stérile ; il fonde la légitimité du pouvoir. Le philosophe-roi est le seul juste gouverneur car il connaît la \cle{Fin} (le Bien) de la Cité.
\end{enumerate}
\textbf{Portée/Limite :} Modèle élitiste (peu d'élus), risque de totalitarisme (le sage sait, le peuple obéit), dualisme métaphysique critiqué par Aristote (« Ami de Platon, mais plus ami de la vérité »).
\textbf{Ouverture :} Cette allégorie fonde la figure moderne de l'\cle{intellectuel} (Zola, Sartre, Foucault) : celui qui « voit » et qui a le devoir d'éclairer la Cité, au risque d'être rejeté.
\end{exoresolu}

\begin{methode}[Rédiger une dissertation philosophique : la méthode du plan dialectique (Thèse / Antithèse / Synthèse)]
\textbf{Objectif :} Produire une réponse organisée, argumentée et problématisée à une question philosophique (sujet de type « Peut-on... ? », « Est-il vrai que... ? », « Faut-il... ? »).
\begin{enumerate}
    \item \textbf{Analyse du sujet (Brouillon - 30 min) :}
        \begin{itemize}
            \item Définir \emph{tous} les termes du sujet (dictionnaire + sens philosophique).
            \item Identifier le \textbf{postulat} (ce que le sujet suppose) et le \textbf{problème} (la tension interne).
            \item Formuler la \textbf{problématique} (question unique, ouverte, qui guide le plan).
            \item Trouver la \textbf{thèse} (réponse spontanée/évidente), l'\textbf{antithèse} (objection radicale), la \textbf{synthèse} (dépassement/nuance).
        \end{itemize}
    \item \textbf{Rédaction de l'Introduction (sur copie) :}
        \begin{enumerate}
            \item \textbf{Accroche :} Citation, fait historique, définition étymologique, paradoxe.
            \item \textbf{Définition des termes} (rigoureuse).
            \item \textbf{Problématique} (formulée en question).
            \item \textbf{Annonce du plan} (Thèse / Antithèse / Synthèse - ou Plan thématique si sujet "Qu'est-ce que... ?").
        \end{enumerate}
    \item \textbf{Rédaction du Développement (3 parties équilibrées) :}
        \begin{itemize}
            \item \textbf{Structure d'un paragraphe :} \textbf{Idée directrice} (phrase thèse) $\rightarrow$ \textbf{Argumentation} (raisonnement, concepts, référence auteur) $\rightarrow$ \textbf{Exemple/Illustration} (concret, précis) $\rightarrow$ \textbf{Mini-transition} (lien vers idée suivante).
            \item \textbf{Transitions majeures :} Résumer la partie finie $\rightarrow$ Montrer sa limite $\rightarrow$ Annoncer la partie qui suit.
        \end{itemize}
    \item \textbf{Rédaction de la Conclusion :}
        \begin{enumerate}
            \item \textbf{Bilan synthétique} (réponse à la problématique).
            \item \textbf{Ouverture} (prolongement : autre auteur, autre domaine, question inverse, actualité).
        \end{enumerate}
    \item \textbf{Relecture :} Orthographe, syntaxe, cohérence, respect du sujet, présence des \cle{mots-clés} du programme.
\end{enumerate}
\end{methode}

\begin{exoresolu}[Dissertation : « La politique est-elle l'art du possible ou la recherche du juste ? »]
\textbf{Énoncé :} Sujet type Bac. Traitez la question en vous appuyant sur les auteurs du programme (Platon, Aristote, Machiavel - hors programme mais référence culturelle autorisée en ouverture -, pensée pharaonique).
\tcblower
\textbf{Solution détaillée :}

\textbf{BROUILLON (Analyse).}
\begin{itemize}
    \item \textbf{Termes :} \cle{Politique} (art d'organiser la Cité/Polis), \cle{Art} (technè, savoir-faire rationnel), \cle{Possible} (réalisable, contingent, rapports de force), \cle{Juste} (conforme à la justice, à l'idéal, au Bien).
    \item \textbf{Postulat :} Opposition entre \emph{Realpolitik} (efficacité, contrainte) et \emph{Idéalisme} (valeurs, norme).
    \item \textbf{Problématique :} La politique doit-elle se réduire à la gestion des contraintes réelles (risque de cynisme) ou viser un idéal de justice inaccessible (risque d'utopie stérile) ? Comment articuler nécessité et légitimité ?
    \item \textbf{Plan dialectique :}
        \item \textbf{Thèse :} La politique comme art du possible (Réalisme / Machiavel / Aristote praticien). La contrainte du réel prime.
        \item \textbf{Antithèse :} La politique comme recherche du juste (Idéalisme / Platon / Maât). La norme éthique guide l'action, sans elle, tyrannie.
        \item \textbf{Synthèse :} La politique comme \emph{prudence} (phronésis) : ancrage dans le réel \emph{au service} du juste. Le possible est le chemin, le juste est la boussole.
\end{itemize}

\textbf{COPIE (Rédaction).}

\textbf{Introduction.}
\textbf{[Accroche]} Otto von Bismarck déclarait : « La politique est l'art du possible, le possible étant l'art du lendemain. » Cette maxime réaliste résume une tentation permanente : réduire l'action politique à l'ajustement aux contraintes du moment.
\textbf{[Définitions]} La \cle{politique} (du grec \emph{polis}, cité) est l'activité d'organisation de la vie collective. L'\cle{art} (technè) implique une rationalité, une technique. Le \cle{possible} relève de la contingence, des rapports de force, de la faisabilité. Le \cle{juste} (dikaion) relève de la norme, de la justice, de la finalité éthique.
\textbf{[Problématique]} Dès lors, la politique est-elle condamnée au cynisme du « réalisme » ou à l'impuissance de l'utopie ? Comment concilier l'exigence de résultat (le possible) et l'exigence de sens (le juste) ?
\textbf{[Annonce]} Nous verrons d'abord que la politique, comme gestion de la Cité, ne peut ignorer le réel (I). Mais que sans référence au Juste, elle perd sa légitimité et devient domination (II). Enfin, que la grandeur politique réside dans la \cle{prudence} (phronésis) : transformer le possible pour approcher le juste (III).

\textbf{Développement.}

\textbf{I. La politique comme art du possible : la contrainte du réel (Thèse réaliste).}
\textbf{Idée directrice :} L'action politique s'exerce dans la matière résistante de l'histoire, de l'économie, des passions humaines ; l'ignorer mène à l'échec.
\textbf{Argumentation :}
\begin{itemize}
    \item \textbf{Aristote (Politique, IV) :} Le législateur ne doit pas chercher la \emph{meilleure} constitution en abstracto, mais la \emph{meilleure possible} compte tenu des mœurs, de l'histoire, de la richesse du peuple. « Le meilleur est souvent l'ennemi du bien. » La \cle{moyenne constitution} (mélange oligarchique/démocratique) assure la stabilité.
    \item \textbf{Machiavel (Le Prince, XVIe) :} Rupture radicale : « Un prince qui veut faire profession de bonté en toutes choses perdra nécessairement le pouvoir. » Il faut savoir « être renard pour connaître les pièges, et lion pour effrayer les loups ». La \cle{raison d'État} prime sur la morale privée. Le \emph{virtù} (énergie, adaptation) domine la \emph{fortuna}.
    \item \textbf{Exemple camerounais :} La gestion de la crise anglophone (depuis 2016). Le « Grand Dialogue National » (2019) a été une tentative de \textbf{réalisme politique} : proposer le statut spécial, la décentralisation, la commission bilingue — mesures \emph{possibles} juridiquement et financièrement — pour désamorcer le conflit, plutôt que l'option militaire pure (impossible à gagner durablement) ou l'indépendance (impossible juridiquement/constitutionnellement).
\end{itemize}
\textbf{Transition :} Mais si la politique se réduit au possible, elle devient pure technique de pouvoir, indifférente au sort des faibles. Le « possible » sans le « juste » n'est que la loi du plus fort.

\textbf{II. La politique comme recherche du juste : l'exigence normative (Antithèse idéaliste).}
\textbf{Idée directrice :} La politique distingue l'humain de l'animal par la \cle{loi} (nomos) et la \cle{justice} (dikaiosunè) ; sans horizon éthique, elle se dénature.
\textbf{Argumentation :}
\begin{itemize}
    \item \textbf{Platon (République, Gorgias) :} Contre les sophistes (Calliclès : « Le juste, c'est l'avantage du plus fort »), Socrate et Platon affirment : « La politique est l'art de soigner les âmes » (comme la médecine soigne les corps). Le \cle{Philosophe-Roi} connaît l'\cle{Idée du Bien} (Soleil). Son art n'est pas de flatter le peuple (rhétorique/pâtisserie) mais de le guider vers la vertu. La Cité juste (Kallipolis) est le modèle régulateur.
    \item \textbf{Pensée Pharaonique (Maât) :} Pharaon n'est pas un gestionnaire de possible, il est le \textbf{garant de Maât}. S'il agit selon Isfet (injustice/chaos) pour « sauver » son pouvoir (possible), il perd la légitimité cosmique. Les \emph{Textes des Pyramides} : « J'ai fait Maât, j'ai détruit Isfet. » La justice n'est pas négociable.
    \item \textbf{Exemple camerounais :} La lutte contre la corruption (Opération \cle{Epervier} / \cle{Sparrowhawk}). C'est une politique du \textbf{Juste} (récupérer l'argent public, sanctionner les prédateurs) qui heurte des « possibles » puissants (réseaux d'influence, immunités de fait). Si l'État ne poursuit que les « petits » et épargne les « grands » par réalisme, il perd sa \emph{raison d'être} (servir le bien commun) et alimente la défiance citoyenne.
\end{itemize}
\textbf{Transition :} Pourtant, une politique purement idéale, déconnectée des rapports de force, conduit à l'impuissance ou à la terreur (Révolution française, Khmers rouges). Il faut articuler les deux.

\textbf{III. La politique comme Prudence (Phronésis) : ancrer le juste dans le possible (Synthèse).}
\textbf{Idée directrice :} La vertu politique par excellence n'est ni la technique (possible) ni la science spéculative (juste abstrait), mais la \cle{prudence} (Aristote, Éthique à Nicomaque VI) : délibérer correctement sur ce qui est bon \emph{et} réalisable \emph{ici et maintenant}.
\textbf{Argumentation :}
\begin{itemize}
    \item \textbf{Aristote :} La \emph{phronésis} concerne les « choses variables » (le contingent). Le politique prudent sait : 1. Discerner la fin (le Juste/Bien commun). 2. Analyser la situation concrète (le Possible/Contraintes). 3. Choisir les \emph{moyens} justes et efficaces (le \emph{kairos}, le moment opportun).
    \item \textbf{Le concept de \cle{Réforme} vs \cle{Révolution} :} La réforme (Bismarck, Roosevelt, Mandela) accepte le cadre institutionnel (possible) pour y introduire de la justice (juste). La révolution pure veut faire table rase du possible au nom du juste $\rightarrow$ risque de Terreur.
    \item \textbf{Exemple camerounais :} La \textbf{Décentralisation (lois 2019/2020)}. C'est un acte de \textbf{prudence politique} :
        \begin{itemize}
            \item \textbf{Juste :} Répondre à l'aspiration légitime des populations à gérer leurs affaires (proximité, développement local, identité).
            \item \textbf{Possible :} Transfert progressif de compétences et de ressources (fiscalité locale), maintien de l'unité nationale (État unitaire décentralisé), formation des élus locaux.
            \item \textbf{Résultat :} Avancée réelle (budgets communaux, plans communaux de développement) mais défis persistants (capacités techniques, ingérence tutélaire). C'est l'art du possible \emph{au service} du juste.
        \end{itemize}
\end{itemize}

\textbf{Conclusion.}
\textbf{[Bilan]} La politique n'est ni \emph{seulement} l'art du possible (cynisme machiavélien, gestionnaire technocrate), ni \emph{seulement} la recherche du juste (utopie platonicienne, rigorisme moralisateur). Elle est l'\textbf{art d'instituer le juste dans le possible}, c'est-à-dire la \cle{prudence}. Le politique véritable est celui qui, connaissant la boussole (le Juste/Bien commun/Maât), sait naviguer dans la tempête (le Possible/Rapports de force/Isfet) sans perdre le cap.
\textbf{[Ouverture]} Cette tension structure l'actualité camerounaise : \textbf{Élections 2025}, \textbf{Décentralisation effective}, \textbf{Justice sociale}. La question n'est pas « possible OU juste » mais « COMMENT rendre le juste possible ? ». Cela renvoie à la définition aristotélicienne de l'homme comme \emph{zoon politikon} : seul l'animal politique transforme la contrainte en loi, et la force en droit.
\end{exoresolu}

\begin{methode}[Appliquer une notion antique à une situation concrète camerounaise : la méthode de l'analogie structurée]
\textbf{Objectif :} Satisfaire au savoir-faire « Appliquer les notions de l'unité à des situations concrètes de la vie courante » en évitant l'anachronisme et le plaquage superficiel.
\begin{enumerate}
    \item \textbf{Extraire la structure conceptuelle :} Isoler le \emph{schème opératoire} de la notion antique (ex. : Maât = Unité Cosmique/Politique/Éthique ; Idée du Bien = Modèle régulateur / Cause finale ; Dialectique = Art du dialogue pour accéder à la vérité).
    \item \textbf{Identifier la situation cible :} Choisir un fait social, politique, éducatif, technique camerounais précis (ex. : palabre, justice coutumière, gouvernance locale, formation professionnelle, dialogue intergénérationnel).
    \item \textbf{Établir l'isomorphisme (analogie de structure) :} Montrer point par point comment la structure antique éclaire la situation moderne (formel : A est à B ce que C est à D). \emph{Ne pas dire « c'est pareil », dire « ça fonctionne sur le même principe logique »}.
    \item \textbf{Qualifier les différences (distance historique) :} Préciser ce qui change (contexte technique, échelle, laïcité, droits de l'homme) pour éviter l'essentialisme (« les Africains sont naturellement maâtiens »).
    \item \textbf{Conclure sur l'apport heuristique :} En quoi la notion antique permet-elle de \emph{mieux comprendre}, \emph{critiquer} ou \emph{orienter} la situation actuelle ?
\end{enumerate}
\end{methode}

\begin{exoresolu}[Application : La Dialectique socratique (Maïeutique) et le « Palabre » camerounais comme méthode de résolution de conflits]
\textbf{Énoncé :} En vous appuyant sur la méthode socratique (ironie, maïeutique, dialectique), analysez le fonctionnement du \cle{Palabre} traditionnel camerounais (ex. : chez les Bamiléké, Bassa, ou Sawa) comme procédure de vérité et de réconciliation. Montrez les convergences structurelles et les différences de fond.
\tcblower
\textbf{Solution détaillée :}

\textbf{1. Structure conceptuelle de la dialectique socratique (Programme : Socrate).}
\begin{itemize}
    \item \textbf{But :} Accéder à la \cle{Vérité} (définition universelle : « Qu'est-ce que le Courage ? », « Qu'est-ce que la Justice ? ») et soigner l'âme (connaissance de soi).
    \item \textbf{Méthode (3 temps) :}
        \begin{enumerate}
            \item \cle{Ironie} : Feindre l'ignorance (« Je ne sais pas »), questionner l'interlocuteur qui croit savoir. Mettre en évidence la \cle{contradiction} interne à son opinion (\emph{elenchos} / réfutation). L'interlocuteur prend conscience de son ignorance (\emph{aporie}).
            \item \cle{Maïeutique} (accouchement des esprits) : Aider l'interlocuteur à \emph{donner naissance lui-même} à la vérité par des questions guidées. Le savoir ne se transmet pas (comme un objet), il se \emph{reconnaît} (anamnèse).
            \item \cle{Dialectique} (au sens socratique) : Dialogue \emph{intersubjectif}, exigeant, sans écrit, où le \cle{Logos} (raison/discours) est le seul juge. Exclusion de la violence, de l'argent (sophistes), de l'autorité externe.
        \end{enumerate}
    \item \textbf{Éthique :} « Nul n'est méchant volontairement » (l'erreur = ignorance). Le dialogue soigne.
\end{itemize}

\textbf{2. Situation cible : Le Palabre (ex. : \cle{Mbao} Bassa, \cle{Tsoung} Bamiléké, \cle{Esèka} Sawa).}
C'est une assemblée publique, solennelle, présidée par un chef/Notables, où un conflit (vol, sorcellerie, limite de champ, succession, meurtre) est « débattu » jusqu'à l'accord.

\textbf{3. Isomorphisme structurel (Analogie de structure).}

\begin{center}
\begin{tabularx}{\linewidth}{|l|X|X|}
\hline
\rowcolor{popblueL}
\textbf{Structure logique} & \textbf{Socrate (Athènes, -400)} & \textbf{Palabre Camerounais (Terroir, aujourd'hui)} \\
\hline
\textbf{Principe de vérité} & La vérité émerge du \textbf{Logos} (discours rationnel confronté), non de l'autorité ou de la force. & La vérité émerge de la \textbf{Parole} (\emph{Nommo}, puissance verbale) échangée en cercle, non du verdict unilatéral. \\
\hline
\textbf{Rôle du « Juge »} & Socrate : \textbf{Accoucheur} (maïeutique). Ne détient pas la réponse, guide le questionnement. Ne vote pas. & Chef/Notables : \textbf{Médiateurs/Facilitateurs}. « Le chef ne juge pas, il fait juger ». Ils reformulent, posent les questions clés, rappellent les précédents (jurisprudence orale). \\
\hline
\textbf{Procédure} & \textbf{Elenchos} : Question $\rightarrow$ Réponse $\rightarrow$ Contradiction $\rightarrow$ Aporie $\rightarrow$ Nouvelle hypothèse. & \textbf{Enquête contradictoire} : Audition parties $\rightarrow$ Témoins $\rightarrow$ Confrontation $\rightarrow$ Mise en évidence du mensonge/incohérence $\rightarrow$ Recherche solution. \\
\hline
\textbf{But final} & \textbf{Consensus rationnel} (Homologia) : Définition acceptée par tous car \emph{reconnue vraie}. Soin de l'âme. & \textbf{Consensus social / Réconciliation} : Accord (\emph{Mbao} = « c'est fini », « paix »). Restauration du lien social rompu. « La parole guérit ». \\
\hline
\textbf{Exclusion} & Violence, sophistique (payer pour apprendre), écrit (fixe la pensée). & Violence privée (vendette), corruption (acheter le verdict), écriture administrative (procès-verbal froid). \\
\hline
\end{tabularx}
\end{center}

\textbf{4. Différences de fond (Distance historique/culturelle).}
\begin{itemize}
    \item \textbf{Ontologie du sujet :} Socrate = \cle{Individu} universel (raison pure, même nature chez tous). Palabre = \cle{Personne relationnelle} (insérée dans lignage, clan, ancêtres). La vérité socratique est \emph{impersonnelle} ; la vérité du palabre est \emph{relationnelle} (restaurer l'harmonie du groupe).
    \item \textbf{Statut de la contradiction :} Chez Socrate, la contradiction est \emph{logique} (A et non-A), outil de progression. Au palabre, la contradiction est souvent \emph{sociale} (offense à l'honneur, rupture d'alliance) ; la résolution passe par la \emph{réparation symbolique} (libation, offrande, compensation) plus que par la seule définition conceptuelle.
    \item \textbf{Universel vs Particulier :} Socrate cherche l'\cle{Universel} (la Justice en soi). Le palabre règle le \cle{Singulier} (cet homme, ce champ, ce clan). Le palabre n'a pas vocation à produire de la \emph{théorie} mais de la \emph{paix sociale}.
    \item \textbf{Genre / Hiérarchie :} Socrate dialogue avec tous (esclaves, femmes, étrangers) \emph{en principe} (Menon). Le palabre traditionnel exclut souvent les femmes, les jeunes, les étrangers de la parole décisive (bien que cela évolue).
\end{itemize}

\textbf{5. Apport heuristique (Ce que Socrate éclaire sur le Palabre, et inversement).}
\begin{itemize}
    \item \textbf{Critique socratique du Palabre :} Risque de \emph{doxa} consensuelle : l'accord obtenu par fatigue, respect hiérarchique ou peur des ancêtres n'est pas une \emph{connaissance}. Le « consensus » peut masquer l'injustice (ex. : femme répudiée « pour la paix du clan »). $\rightarrow$ Besoin d'une \emph{institutionnalisation} des droits (État de droit) garantissant l'universel.
    \item \textbf{Apport du Palabre à la démocratie moderne :} Contre la justice étatique lente, coûteuse, adversariale (vainqueur/vaincu), le palabre montre l'efficacité de la \textbf{justice restaurative} : responsabilité de l'auteur, réparation pour la victime, réintégration du coupable. Inspire les \textbf{Mécanismes Alternatifs de Règlement des Différends (MARD)} et la \cle{Justice Transitionnelle} (CVR - Commission Vérité Réconciliation).
    \item \textbf{Synthèse :} Le Palabre est une \cle{dialectique concrète, incarnée, communautaire}. Socrate en est la \cle{formalisation abstraite, universaliste, individuelle}. Les deux affirment que \textbf{la vérité naît de la parole librement confrontée}, non du silence imposé par la force.
\end{itemize}

\textbf{Conclusion :} Appliquer la notion de dialectique socratique au palabre ne consiste pas à dire « les Anciens faisaient du Socrate sans le savoir », mais à révéler la \textbf{rationalité communicative} (Habermas) inhérente à la tradition africaine de gestion du conflit. Cela permet de légitimer les MARD modernes par une généalogie culturelle propre, tout en les corrigeant par l'exigence universaliste des Droits de l'Homme (héritage socratique/stoïcien/Kantien).
\end{exoresolu}

\begin{methode}[Justifier une démarche philosophique : la méthode de l'argumentation explicite (Savoir-faire « Rédiger et justifier »)]
\textbf{Objectif :} Rendre visible le cheminement de la pensée : \emph{Pourquoi} je dis cela ? \emph{Sur quoi} je m'appuie ? \emph{Comment} j'articule ?
\begin{enumerate}
    \item \textbf{Expliciter le point de départ :} Question, intuition, fait observed, citation, paradoxe.
    \item \textbf{Formuler l'hypothèse de travail :} Réponse provisoire, thèse que l'on va tester.
    \item \textbf{Construire l'argumentation (Chaîne logique) :}
        \begin{itemize}
            \item \textbf{Prémisse 1 (Factuelle/Conceptuelle) :} Définition, fait établi, citation d'auteur.
            \item \textbf{Prémisse 2 (Logique/Principielle) :} Principe philosophique (ex. : « Tout ce qui change est imparfait », « L'homme est un animal politique »).
            \item \textbf{Inférence (Raisonnement) :} Déduction, induction, analogie, \emph{reductio ad absurdum}, dilemme.
        \end{itemize}
    \item \textbf{Anticiper l'objection (Esprit critique) :} « Mais on pourrait objecter que... » $\rightarrow$ Y répondre (réfutation ou intégration).
    \item \textbf{Formuler la conclusion justifiée :} « C'est pourquoi, au terme de cette analyse, on peut affirmer que... » (Nuancé, conditionnel si besoin).
    \item \textbf{Traçabilité :} Citer explicitement les auteurs/concepts du programme utilisés comme \cle{repères} (ex. : « Au sens de Platon... », « Contrairement à Héraclite... »).
\end{enumerate}
\end{methode}

\begin{exoresolu}[Justification de démarche : « En quoi la distinction platonicienne entre Opinion (Doxa) et Science (Epistémè) est-elle un outil critique pour analyser les "Fake News" au Cameroun ? »]
\textbf{Énoncé :} Justifiez, en suivant une démarche philosophique rigoureuse, que la distinction platonicienne entre \emph{doxa} et \emph{epistémè} permet d'analyser le phénomène des « Fake News » sur les réseaux sociaux au Cameroun.
\tcblower
\textbf{Solution détaillée :}

\textbf{Démarche justifiée étape par étape.}

\textbf{1. Point de départ (Constat factuel).}
Au Cameroun, la prolifération des \cle{Fake News} (infox) sur WhatsApp, Facebook, TikTok (ex. : rumeurs de coup d'État, fausses nominations, traitements miracles COVID/Ebola, manipulation électorale) crée une confusion généralisée. Le citoyen ne sait plus « quoi croire ». C'est une crise de la \cle{croyance} vs \cle{savoir}.

\textbf{2. Hypothèse de travail.}
La distinction platonicienne (République, Livres V-VII : Ligne, Caverne, Soleil) entre \cle{Doxa} (Opinion/Croyance) et \cle{Epistémè} (Science/Connaissance) offre un \emph{cadre conceptuel opératoire} pour diagnostiquer la pathologie informationnelle actuelle et prescrire un remède éducatif.

\textbf{3. Argumentation structurée (Chaîne logique).}

\textbf{Prémisse 1 : Analyse conceptuelle (Le cadre platonicien).}
\begin{itemize}
    \item \textbf{Doxa (Opinion) :} Porte sur le \emph{Devenir} (le sensible, le changeant, les apparences). Elle est \textbf{instable}, \textbf{non justifiée} (ou mal justifiée), \textbf{particulière}. Deux degrés : \emph{Eikasia} (imagination/illusion : ombres, images, rumeurs) et \emph{Pistis} (foi/croyance : objets physiques, discours non vérifiés). $\rightarrow$ \textbf{Caractère clé :} « Croire sans savoir », adhésion affective/sociale, pas de \emph{logos} démonstratif.
    \item \textbf{Epistémè (Science) :} Porte sur l'\emph{Être} (les Idées/Formes, l'intelligible, l'universel, l'immuable). Elle est \textbf{stable}, \textbf{justifiée} (démonstration, dialectique), \textbf{universelle}. Deux degrés : \emph{Dianoia} (pensée discursive : maths, hypothèses) et \emph{Noêsis} (intelligence intuitive : saisie des principes premiers, Idée du Bien). $\rightarrow$ \textbf{Caractère clé :} « Savoir pourquoi », nécessité logique, objectivité.
\end{itemize}

\textbf{Prémisse 2 : Qualification du phénomène « Fake News » (Application).}
\begin{itemize}
    \item La Fake News est une \textbf{information fausse, fabriquée, diffusée intentionnellement (désinformation) ou par négligence (mésinformation), qui mime l'apparence du vrai} (logo officiel, ton journalistique, vidéo truquée).
    \item Elle circule dans le \textbf{flux instantané} des réseaux sociaux (Héraclite : « tout coule »), sans vérification possible immédiate.
    \item Elle s'adresse aux \textbf{affects} (peur, colère, espoir, identité ethnique/religieuse) et aux \textbf{biais cognitifs} (confirmation, autorité).
    \item \textbf{Conclusion intermédiaire 1 :} La Fake News relève \textbf{ontologiquement et épistémologiquement de la Doxa (niveau Eikasia/Pistis)}. Elle est \emph{ombre} (copie d'écran truquée) prise pour \emph{réalité}. Elle n'a pas de \emph{logos} propre (pas de source traçable, pas de méthode de vérification).
\end{itemize}

\textbf{Prémisse 3 : Mécanisme de la caverne (Analogie structurelle).}
\begin{itemize}
    \item \textbf{Les porteurs d'objets (Fabricants de Fake News) :} Acteurs politiques, agences de com, trolls, algorithmes. Ils manipulent les \emph{ombres} (pixels, sons) derrière le \emph{muret} (anonymat, cryptage, serveur étranger).
    \item \textbf{Les prisonniers (Usagers passifs) :} Enchaînés par l'illettrisme numérique, la paresse cognitive, la confiance aveugle dans le « transféré par un proche ». Ils regardent le mur (écran du smartphone).
    \item \textbf{L'éblouissement (Vérification/Fact-checking) :} Tourner la tête (chercher la source, croiser les sources, consulter le site officiel) fait mal (effort cognitif, temps, rupture du confort tribal).
    \item \textbf{Conclusion intermédiaire 2 :} L'utilisateur non éduqué à l'esprit critique est \emph{structurellement} un prisonnier de la caverne numérique. La \cle{Media and Information Literacy (MIL)} / Éducation aux Médias est l'équivalent moderne du \textbf{tournant de l'âme (periagogè)}.
\end{itemize}

\textbf{4. Anticipation de l'objection (Esprit critique).}
\textbf{Objection :} « Platon est élitiste : il réserve l'Epistémè à une minorité (Gardiens/Philosophes). La démocratie numérique vise l'inclusion. Utiliser Platon pour lutter contre les Fake News, c'est risquer une \emph{censure éclairée} (décider pour le peuple ce qui est Vrai). »
\textbf{Réponse (Intégration/Nuance) :}
\begin{itemize}
    \item Platon distingue \emph{connaissance experte} (réservée) et \emph{droit à la vérité} (universel). Le philosophe \emph{redescend} dans la caverne (devoir politique).
    \item Aujourd'hui, l'\cle{Epistémè} accessible à tous n'est pas la métaphysique des Idées, mais la \textbf{méthode scientifique et journalistique} : vérifiabilité, traçabilité, contradictoire, transparence des sources. C'est une \emph{epistémè} \textbf{procédurale} (Popper, Dewey), démocratique.
    \item L'outil critique n'est pas « Platon a raison sur tout », mais « La distinction Doxa/Epistémè permet de \textbf{nommer} la différence entre \emph{partager une émotion} et \emph{valider une information} ».
\end{itemize}

\textbf{5. Conclusion justifiée.}
\textbf{Au terme de cette démarche,} on peut affirmer que la distinction platonicienne est un \textbf{outil critique pertinent} car elle :
\begin{enumerate}
    \item \textbf{Diagnostique :} Identifie la Fake News comme \emph{pathologie de la Doxa} (illusion prise pour réalité, flux non maîtrisé).
    \item \textbf{Hiérarchise :} Replace l'exigence de \emph{vérification} (logos, source, méthode) au-dessus de l'adhésion \emph{virale} (like, partage, foi).
    \item \textbf{Prescrit :} Définit l'objectif éducatif : non pas « croire les bons médias », mais \emph{acquérir la méthode} pour passer de l'opinion (doxa) à la connaissance justifiée (epistémè procédurale) $\rightarrow$ \textbf{Éducation à l'esprit critique (Citoyenneté numérique)}.
\end{enumerate}
\textbf{Justification finale :} Cette démarche est rigoureuse car elle : (1) Définit ses concepts (Doxa/Epistémè), (2) Les applique sans anachronisme (Fake News = Eikasia numérique), (3) Utilise l'analogie structurelle (Caverne = Écosystème info), (4) Intègre l'objection démocratique (Élitisme vs Éducation pour tous), (5) Produit une conclusion opératoire (MIL = Periagogè moderne).
\end{exoresolu}

\begin{methode}[Synthétiser un problème transversal : la méthode du « Fil rouge » (Unité vs Multiple)]
\textbf{Objectif :} Répondre au savoir-faire « Appliquer les notions de l'unité à des situations concrètes » en montrant comment le fil conducteur « Unité / Multiple » structure toute la leçon (Pharaonique $\rightarrow$ Présocratiques $\rightarrow$ Socrate $\rightarrow$ Platon).
\begin{enumerate}
    \item \textbf{Poser le problème métaphysique originel :} Comment l'Un (le Même, le Permanent, le Principe) peut-il rendre compte du Multiple (l'Autre, le Changeant, les Phénomènes) sans le nier ni s'y dissoudre ?
    \item \textbf{Suivre le fil à travers les 4 étapes de la leçon :}
        \begin{itemize}
            \item \textbf{Pharaonique :} Unité \emph{cosmico-politique} (Maât). L'Un (Maât/Pharaon) \emph{contient} et \emph{ordonne} le Multiple (peuple, nature, temps). Solution : \textbf{Hiérarchie sacrée / Harmonie}.
            \item \textbf{Présocratiques :} Rupture. Héraclite : L'Un \emph{est} le Multiple (Devenir/Feu/Logos). Parménide : L'Un \emph{exclut} le Multiple (Être plein/Immobile). Solution : \textbf{Opposition radicale} (Aporie).
            \item \textbf{Socrate :} Déplacement anthropologique. L'Un (Définition universelle / Soi) se cherche \emph{dans} le Multiple (opinions particulières / Dialogue). Solution : \textbf{Recherche dialectique / Universalité dans le particulier}.
            \item \textbf{Platon :} Systématisation dualiste. L'Un (Idées/Bien) \emph{transcend} le Multiple (Sensible/Devenir) et lui donne sens (Participation/Imitation). Solution : \textbf{Dualisme ontologique / Participation}.
        \end{itemize}
    \item \textbf{Dégager la leçon philosophique :} L'histoire de la philosophie antique est l'histoire des \emph{réponses successives} au problème de l'Un et du Multiple. Chaque réponse ouvre un champ nouveau (Physique, Logique, Éthique, Politique, Métaphysique).
    \item \textbf{Appliquer à une situation unificatrice camerounaise :} Ex. : « Comment l'État camerounais (Un) gère-t-il la diversité culturelle, linguistique, régionale (Multiple) ? » $\rightarrow$ Maât (Unité organique/centralisée) vs Démocratie/Décentralisation (Unité contractuelle/fédérative) vs Dialogue national (Méthode socratique).
\end{enumerate}
\end{methode}

\begin{exoresolu}[Synthèse : L'Un et le Multiple dans l'Antiquité et au Cameroun aujourd'hui]
\textbf{Énoncé :} « La philosophie antique est la recherche de l'Un derrière le Multiple. » Illustrer cette affirmation en parcourant les quatre moments de la leçon (Pharaonique, Présocratiques, Socrate, Platon). Montrer ensuite comment ce schéma éclaire le défi de l'unité nationale au Cameroun.
\tcblower
\textbf{Solution détaillée :}

\textbf{1. Le fil rouge métaphysique : L'Un et le Multiple.}
La philosophie naît de l'étonnement devant la \cle{multiplicité} du monde (choses, changements, opinions) et de la soif de l'\cle{unité} (principe unique, loi unique, vérité unique, bien unique). Le génie antique est d'avoir inventé \emph{quatre figures majeures} de cette relation.

\textbf{2. Parcours des quatre moments (Application du fil rouge).}

\begin{center}
\begin{tabularx}{\linewidth}{|l|X|X|X|}
\hline
\rowcolor{popblueL}
\textbf{Moment} & \textbf{L'Un (Principe)} & \textbf{Le Multiple (Phénomènes)} & \textbf{Mode de relation (Solution)} \\
\hline
\textbf{1. Pharaonique (Maât)} & \cle{Maât} (Ordre cosmique, Vérité, Justice, Mesure). \textbf{Un unique, immanent, vivant.} & \cle{Isfet} (Chaos), Peuple, Nature, Temps, Étrangers. \textbf{Menace permanente.} & \textbf{Hiérarchie sacrée / Maintien.} Pharaon (Un) réalise Maât \emph{sur} le Multiple. Rituel, Loi, Justice = \emph{Acte cosmique}. L'Un \textbf{ordonne} le Multiple. \\
\hline
\textbf{2. Héraclite} & \cle{Logos / Feu}. \textbf{Un dynamique.} « Tout est Un » (Fragment 50). & \cle{Devenir} (Choses, Contraires : jour/nuit, vie/mort). \textbf{Flux perpétuel.} & \textbf{Identité / Unité des contraires.} L'Un \textbf{EST} le Multiple en mouvement. L'Un est la \emph{loi} du Multiple (mesure d'échange). Pas de substance fixe. \\
\hline
\textbf{3. Parménide} & \cle{L'Être}. \textbf{Un absolu, statique, sphérique.} « L'Être est, le Non-être n'est pas. » & \cle{Le Devenir / L'Apparence / L'Opinion (Doxa)}. \textbf{Illusion des sens.} & \textbf{Exclusion / Négation.} L'Un \textbf{NIE} le Multiple. Le Multiple = Non-être = Néant. Seule la Raison (Un) atteint l'Un. \\
\hline
\textbf{4. Socrate} & \cle{L'Universel (Définition / Essence / Soi)}. « Qu'est-ce que la Justice ? » & \cle{Le Particulier (Opinions / Exemples / Actes concrets / Sophistes)}. & \textbf{Dialectique / Maïeutique.} L'Un \textbf{ÉMERGE} du Multiple par le dialogue. L'universel n'est pas \emph{au-dessus} (Platon) ni \emph{en-dessous} (Héraclite), mais \emph{dans} le particulier bien examiné. \\
\hline
\textbf{5. Platon} & \cle{Les Idées (Formes) / Idée du Bien (Soleil)}. \textbf{Un transcendant, multiple en soi (Idées multiples).} & \cle{Le Sensible (Choses particulières, Devenir, Ombres)}. \textbf{Copies imparfaites.} & \textbf{Participation (Méthexis) / Imitation (Mimèsis).} L'Un (Modèle) \textbf{FONDE} le Multiple (Copie). Le Multiple « participe » de l'Un. Dualisme ontologique. Politique : Philosophe (Un) guide Cité (Multiple). \\
\hline
\end{tabularx}
\end{center}

\textbf{3. Leçon philosophique synthétique.}
L'histoire de l'antiquité est une \textbf{progression dans la conceptualisation de la relation Un/Multiple} :
\begin{enumerate}
    \item \textbf{Fusion mystique} (Pharaonique : Un = Ordre vital).
    \item \textbf{Opposition physique} (Présocratiques : Un = Substance vs Flux ; ou Un = Être vs Néant).
    \item \textbf{Recherche logique/éthique} (Socrate : Un = Concept universel cherché dans le particulier).
    \item \textbf{Systématisation métaphysique/politique} (Platon : Un = Modèle transcendant cause du Multiple).
\end{enumerate}
Cette progression \emph{libère la raison} : elle passe du mythe (Maât) à la physique (Feu/Être), à l'éthique (Soi), à la métaphysique (Idées).

\textbf{4. Application au défi de l'Unité Nationale au Cameroun.}
Le Cameroun est l'archétype du \textbf{Multiple} : 250+ ethnies, 2 langues officielles, 2 systèmes juridiques (Common Law / Droit Civil), 2 héritages coloniaux, religions multiples, géographie contrastée (Soudan, Forêt, Côte, Montagnes). L'enjeu : construire l'\textbf{Un} (Nation, État, Peuple) sans détruire le Multiple.

\begin{itemize}
    \item \textbf{Modèle « Maât » (Centralisme organique / « Unité dans la paix ») :} L'État (Pharaon) incarne l'Unité. La diversité est gérée par \emph{hiérarchie} et \emph{intégration} au centre (Ministères, Préfets, Chefs traditionnels nommés). \textbf{Force :} Stabilité, symbole fort. \textbf{Risque :} Étouffement du Multiple (marginalisation anglophone, minorités), confusion Parti/État, Isfet si le centre faiblit.
    \item \textbf{Modèle « Héraclite » (Flux identitaire) :} L'unité n'est pas un bloc, c'est un \emph{processus} permanent de métissage, d'échanges, de « vivre-ensemble » quotidien (mariages mixtes, commerce, langues véhiculaires : Camfranglais, Pidgin). \textbf{Force :} Réalisme sociologique, vitalité. \textbf{Risque :} Dilution de l'État, absence de repères fixes, instrumentalisation des identités.
    \item \textbf{Modèle « Parménide » (Unité rigide / Négation de la différence) :} « Un peuple, un but, une foi » (Slogan). La différence (fédéralisme, particularismes) = Non-être / Danger / Sécession. \textbf{Risque :} Violence symbolique ou réelle contre la diversité (crise NOSO).
    \item \textbf{Modèle « Socrate » (Dialogue / Palabre national) :} L'Unité se \textbf{construit} par le \emph{dialogue contradictoire} (Grand Dialogue National, Concertations). On ne décrète pas l'Un, on le \emph{cherche} ensemble en acceptant le conflit d'opinions (Doxa) pour viser un consensus supérieur (Epistémè politique = Contrat social). \textbf{Force :} Légitimité, inclusion. \textbf{Condition :} Liberté d'expression, égalité des partenaires.
    \item \textbf{Modèle « Platon » (Constitution / Idée du Bien commun) :} L'Unité est une \textbf{Idée régulative} (Constitution, Décentralisation, Bilingue effectif, Équité régionale). Les institutions (Sénat, Conseil Constitutionnel, CTD) sont les \emph{Idées} qui structurent le Multiple concret. Les gouvernants (Gardiens) doivent connaître le Bien commun (Formation, Mérite, Intégrité) pour gouverner \emph{pour} le Multiple.
\end{itemize}

\textbf{Conclusion.}
La philosophie antique nous apprend qu'il n'y a \textbf{pas de solution unique, définitive et magique} au problème Un/Multiple.
\begin{itemize}
    \item Le \textbf{Pharaon} nous rappelle la \textbf{nécessité d'un centre fort} (État, Symboles, Souveraineté).
    \item \textbf{Héraclite} nous rappelle la \textbf{réalité du mouvement} (Décentralisation, Adaptation, Jeunesse).
    \item \textbf{Parménide} nous met en garde contre le \textbf{refus dogmatique de la différence}.
    \item \textbf{Socrate} nous donne la \textbf{méthode} : le dialogue vrai, exigeant, sans exclusion.
    \item \textbf{Platon} nous donne l'\textbf{exigence normative} : des institutions justes (Idées) servies par des compétents vertueux (Gardiens).
\end{itemize}
L'unité camerounaise ne sera pas \emph{donnée} (Maât), ni \emph{figée} (Parménide), ni \emph{fluide} (Héraclite), elle sera \textbf{construite par le dialogue (Socrate) sous la règle du droit juste (Platon)}. C'est tout le sens de la \cle{Décentralisation} comme \emph{participation} du Multiple à l'Un.
\end{exoresolu}

\begin{attention}[Les 5 erreurs qui coûtent des points au Bac (Probatoire / Baccalauréat Technique)]
\begin{enumerate}
    \item \textbf{Confondre « Mythos » et « Logos » (Pharaonique / Présocratiques).} \emph{Erreur :} « Les Égyptiens n'avaient pas de philosophie, c'était de la religion. » \textbf{Correction :} La pensée pharaonique (Maât) est une \emph{philosophie de l'ordre} (onto-politique), rationnelle dans son contexte, antérieure au Logos grec. Ne pas plaquer la définition grecque (théorétique) sur l'antiquité africaine. \emph{Critère :} Distinguer \emph{mythe} (récit fondateur) et \emph{philosophie} (discours rationnel sur le principe), mais reconnaître la \emph{raison dans le mythe} (théologie politique de Maât).
    \item \textbf{Réduire Héraclite à « tout change » et Parménide à « rien ne change ».} \emph{Erreur :} Caricature scolastique. \textbf{Correction :} Héraclite : « Tout change \emph{selon une mesure (Logos)} » (loi, intelligence). Parménide : « Le changement est \emph{logiquement impossible} pour l'Être vrai », mais il explique l'\emph{apparence} (Doxa) dans la Voie de l'Opinion. \emph{Nuance exigible :} Citer le \cle{Logos} (Héraclite) et la distinction \cle{Voie de Vérité / Voie d'Opinion} (Parménide).
    \item \textbf{Confondre Socrate et Platon (Le « Socratisme » de Platon).} \emph{Erreur :} Attribuer la Théorie des Idées, l'Immortalité de l'âme (démonstration), la République (Kallipolis) à Socrate. \textbf{Correction :} Socrate = \cle{Méthode} (Ironie, Maïeutique, Elenchos, Éthique intellectuelle : « Connais-toi toi-même », « Nul n'est méchant volontairement »), \emph{pas de système métaphysique écrit}. Platon = \cle{Système} (Idées, Âme tripartite, Politique, Cosmologie). \emph{Repère :} Dans les dialogues de jeunesse (Apologie, Criton, Gorgias, Lachès) $\rightarrow$ Socrate historique. Dans les dialogues de maturité (République, Phédon, Banquet, Timée) $\rightarrow$ Platon parle par la bouche de Socrate.
    \item \textbf{Traiter l'Allégorie de la Caverne comme une « simple histoire » ou une « métaphore de l'éducation scolaire ».} \emph{Erreur :} « C'est l'histoire d'un prof qui emmène ses élèves au soleil. » \textbf{Correction :} C'est une \textbf{ontologie} (deux mondes), une \textbf{épistémologie} (4 degrés de connaissance : Eikasia, Pistis, Dianoia, Noêsis), une \textbf{psychologie} (tournant de l'âme), une \textbf{politique} (devoir de redescente, philosophe-roi). \emph{Exigence :} Citer les 4 degrés de la Ligne (Rép. VI) et l'Idée du Bien (Soleil) comme cause de l'Être et de la Connaissance.
    \item \textbf{Oublier le « Cameroun » dans l'application / L'anachronisme sauvage.} \emph{Erreur :} « Chez les Bamiléké, c'est exactement comme chez Platon. » Ou : Aucune référence au contexte camerounais dans la dissertation/explication. \textbf{Correction :} Le programme exige \textbf{« Appliquer les notions... à des situations concrètes de la vie courante »}. Toute dissertation ou explication doit comporter \textbf{au moins un exemple camerounais précis, daté, nommé} (Opération Épervier, Grand Dialogue National, Ngondo, Décentralisation 2020, Camfranglais, Justice coutumière Mbao, Bourse d'études, etc.). L'analogie doit être \emph{structurée} (isomorphisme), pas \emph{identitaire}.
\end{enumerate}
\end{attention}

\newpage

\coursec{Exercices}

\courssub{Je vérifie mes connaissances}

\begin{exercice}[QCM : repères de la philosophie antique]
Pour chaque question, choisis la bonne réponse (une seule par question).
\begin{enumerate}
\item Dans la pensée égyptienne pharaonique, \cle{Maât} désigne :
\begin{enumerate}
\item la déesse de la guerre ;
\item le principe d'ordre, de justice et d'harmonie cosmique ;
\item le fleuve Nil ;
\item l'au-delà réservé aux pharaons.
\end{enumerate}
\item Pour Héraclite, la réalité est fondamentalement :
\begin{enumerate}
\item immobile et une ;
\item un devenir perpétuel symbolisé par le feu ;
\item une copie du monde des Idées ;
\item le hasard des atomes.
\end{enumerate}
\item Parménide d'Élée affirme que :
\begin{enumerate}
\item « on ne se baigne jamais deux fois dans le même fleuve » ;
\item l'Être est, le non-être n'est pas : le changement est une illusion des sens ;
\item l'homme est la mesure de toutes choses ;
\item la vertu est un don des dieux.
\end{enumerate}
\item La \cle{maïeutique} socratique consiste à :
\begin{enumerate}
\item accoucher les femmes ;
\item faire accoucher les esprits de la vérité par le questionnement ;
\item écrire des dialogues ;
\item enseigner la rhétorique contre rémunération.
\end{enumerate}
\item Dans l'allégorie de la caverne, Platon illustre :
\begin{enumerate}
\item la vie des prisonniers de guerre ;
\item le passage de l'opinion (\cle{doxa}) à la connaissance véritable (\cle{épistémè}) ;
\item l'organisation politique d'Athènes ;
\item la vie après la mort.
\end{enumerate}
\end{enumerate}
\end{exercice}

\begin{exercice}[Vrai ou faux — justifier]
Indique si chaque affirmation est vraie ou fausse, puis justifie ta réponse en une ou deux phrases.
\begin{enumerate}
\item La philosophie pharaonique ignore toute réflexion sur la justice et l'ordre du monde.
\item Héraclite et Parménide défendent exactement la même thèse sur le changement.
\item Socrate a écrit de nombreux ouvrages pour transmettre sa doctrine.
\item Pour Platon, le monde sensible est le monde de la réalité pleine et parfaite.
\item Socrate, condamné à mort, refuse de s'évader de prison par respect des lois de la cité.
\end{enumerate}
\end{exercice}

\begin{exercice}[Définitions]
Définis en quelques lignes chacune des notions suivantes, en donnant pour chacune un exemple :
\begin{enumerate}
\item le \cle{Logos} (chez Héraclite) ;
\item l'\cle{ironie} socratique ;
\item la théorie des \cle{Idées} (ou des Formes) chez Platon ;
\item le \cle{devenir}.
\end{enumerate}
\end{exercice}

\begin{exercice}[Auteurs et thèses]
Recopie et complète le tableau suivant en associant à chaque penseur sa thèse fondamentale et son concept clé.
\begin{center}
\begin{tabularx}{\linewidth}{|l|X|X|}
\hline
\rowcolor{popblueL} \textbf{Penseur} & \textbf{Thèse fondamentale} & \textbf{Concept clé} \\
\hline
Héraclite & & \\
\hline
Parménide & & \\
\hline
Socrate & & \\
\hline
Platon & & \\
\hline
\end{tabularx}
\end{center}
\end{exercice}

\courssub{Je m'entraîne}

\begin{exercice}[Le fleuve d'Héraclite au quotidien]
Héraclite affirme : « On ne se baigne jamais deux fois dans le même fleuve. »
\begin{enumerate}
\item Explique le sens de cette formule.
\item Illustre-la par deux exemples tirés de la vie camerounaise actuelle (par exemple : l'évolution d'une ville comme Douala ou Yaoundé, les transformations d'une langue comme le camfranglais).
\item Ce constat du changement permanent signifie-t-il que rien ne demeure ? Justifie ta réponse en mobilisant la notion de \cle{Logos}.
\end{enumerate}
\end{exercice}

\begin{exercice}[Maât et l'ordre social]
Dans l'Égypte pharaonique, le pharaon a pour mission de maintenir \cle{Maât}, c'est-à-dire l'ordre cosmique, la justice et l'équilibre de la société.
\begin{enumerate}
\item Explique ce que signifie le principe de Maât.
\item Compare ce principe au rôle que joue aujourd'hui la Constitution dans un État moderne comme le Cameroun : quels points communs ? quelles différences ?
\item Peut-on dire que le respect des règles dans un lycée technique relève d'une exigence comparable à Maât ? Argumente.
\end{enumerate}
\end{exercice}

\begin{exercice}[Dialogue socratique]
Un de tes camarades affirme : « Être courageux, c'est ne jamais avoir peur. »
\begin{enumerate}
\item Rédige un court dialogue (8 à 12 répliques) dans lequel tu appliques la méthode socratique : \cle{ironie} (questions qui mettent en difficulté la certitude du camarade) puis \cle{maïeutique} (questions qui l'aident à formuler une définition plus juste du courage).
\item Dégage la définition du courage à laquelle votre dialogue aboutit.
\item Explique en quelques lignes en quoi cette méthode illustre le souci de l'âme et l'examen de soi chez Socrate.
\end{enumerate}
\end{exercice}

\begin{exercice}[Sortir de la caverne]
Dans l'allégorie de la caverne (Platon, \emph{La République}, livre VII), des prisonniers enchaînés prennent des ombres pour la réalité.
\begin{enumerate}
\item Raconte en quelques lignes les étapes de la libération du prisonnier.
\item Identifie, dans la vie d'un élève camerounais d'aujourd'hui, deux situations comparables à celle des prisonniers (par exemple : les rumeurs sur les réseaux sociaux, les préjugés sur certaines filières techniques).
\item Pourquoi, selon Platon, le philosophe libéré doit-il redescendre dans la caverne ? Quelle leçon en tires-tu pour le rôle des personnes éduquées dans la société ?
\end{enumerate}
\end{exercice}

\courssub{Je me prépare au Probatoire / Baccalauréat technique}

\begin{exercice}[Dissertation : le changement et l'identité]
\textbf{Sujet :} \emph{Le changement est-il une perte d'identité ?}

\medskip
Contexte : dans de nombreuses familles camerounaises, on s'interroge : un jeune qui quitte son village pour la ville, qui parle davantage le français ou l'anglais que sa langue maternelle, qui adopte de nouveaux modes de vie, perd-il son identité, ou la transforme-t-il simplement ?

\medskip
\textbf{Travail à faire :} rédige une dissertation complète (introduction, développement en trois parties, conclusion) en suivant les consignes :
\begin{enumerate}
\item Dans l'introduction, amène le sujet à partir de la situation camerounaise décrite, définis les termes « changement » et « identité », puis formule une problématique claire.
\item Dans une première partie, expose la thèse d'Héraclite : tout s'écoule, le devenir est la loi du réel ; montre qu'alors l'identité n'est jamais figée.
\item Dans une deuxième partie, expose la thèse opposée de Parménide : l'Être est un, immobile ; le changement n'est qu'apparence sensible ; ce qui est véritablement demeure.
\item Dans une troisième partie, dépasse l'opposition en mobilisant Platon (distinction entre le monde sensible changeant et les Idées stables) : l'identité véritable d'une personne ou d'une culture est-elle dans ses apparences changeantes ou dans un fond permanent ?
\item Illustre chaque partie par des exemples camerounais précis (urbanisation, langues, traditions, chefferies).
\end{enumerate}
\end{exercice}

\begin{exercice}[Dissertation : loi et liberté]
\textbf{Sujet :} \emph{Obéir aux lois, est-ce renoncer à sa liberté ?}

\medskip
Contexte : au Cameroun, comme dans toute société organisée, les citoyens sont tenus de respecter la Constitution et les lois. Certains jeunes estiment que « être libre, c'est faire ce qu'on veut » et voient dans toute loi une entrave.

\medskip
\textbf{Travail à faire :} rédige une dissertation complète en suivant les consignes :
\begin{enumerate}
\item Problématise le sujet dans l'introduction : la liberté est-elle l'absence de contrainte, ou suppose-t-elle au contraire des règles ?
\item Première partie : montre que la loi peut apparaître comme une limitation de la liberté individuelle (appuie-toi sur le sentiment du jeune décrit dans le contexte).
\item Deuxième partie : mobilise Socrate dans le \emph{Criton} : condamné injustement, il refuse pourtant de s'évader, car désobéir aux lois de la cité, c'est détruire la cité qui l'a fait naître et éduqué. Analyse cet argument.
\item Troisième partie : confronte cette position à Platon (la caverne : la vraie liberté est libération intérieure par la connaissance) et à Maât (la justice comme condition de l'harmonie collective) ; montre que la loi juste peut être la condition de la liberté de tous.
\item Conclus en prenant position de manière argumentée, avec un exemple tiré de la vie civique camerounaise (respect du code de la route, paiement des impôts, règlement scolaire).
\end{enumerate}
\end{exercice}

\begin{exercice}[Explication de texte : l'Apologie de Socrate]
\textbf{Texte} (Platon, \emph{Apologie de Socrate}, 29d-30c, adapté) :

\begin{quote}
« Je vais vous dire pourquoi j'ai porté ce nom [\dots]. Si vous me disiez : "Socrate, nous te laisserons aller, mais à condition que tu cesses de philosopher" [\dots], je vous répondrais : Athéniens, je vous respecte et je vous aime, mais j'obéirai au dieu plutôt qu'à vous, et tant que je respirerai, je ne cesserai de philosopher, de vous exhorter et de vous faire mes reproches, disant à chacun de vous : "Comment, ô toi le meilleur des hommes, citoyen d'Athènes, la plus grande et la plus renommée des cités, n'as-tu pas honte de t'occuper de tes richesses [\dots] et de ne pas t'occuper de ta pensée, de la vérité et de ton âme, pour la rendre meilleure ?" »
\end{quote}

\textbf{Questions :}
\begin{enumerate}
\item Dégage la thèse du texte : quelle mission Socrate s'attribue-t-il, et pourquoi refuse-t-il d'y renoncer, même au prix de sa vie ?
\item Analyse les arguments de Socrate : pourquoi dit-il obéir « au dieu plutôt qu'à vous » ? En quoi ce choix engage-t-il une hiérarchie des valeurs ?
\item Explique les concepts suivants tels qu'ils fonctionnent dans le texte : l'\cle{ironie} (« ô toi le meilleur des hommes »), le \cle{soin de l'âme}, le rapport du philosophe à la \cle{cité}.
\item Pourquoi Socrate se compare-t-il, ailleurs dans l'\emph{Apologie}, à un taon (mouche à bœuf) qui pique un cheval paresseux ? Quelle fonction le philosophe remplit-il ainsi dans la cité ?
\item Portée actuelle : existe-t-il aujourd'hui, au Cameroun, des figures ou des pratiques qui jouent un rôle comparable d'« examen des consciences » collectif (journalistes, enseignants, chefs traditionnels, artistes) ? Discute en deux paragraphes argumentés.
\end{enumerate}
\end{exercice}

\courssub{Corrigés des exercices de vérification}

\begin{aretenir}[Points de vigilance pour l'examen]
\begin{itemize}
\item \textbf{Vocabulaire précis :} Distinguer \emph{doxa} (opinion/croyance sensible) et \emph{épistémè} (science/connaissance intelligible) chez Platon.
\item \textbf{Ne pas confondre :} Le « feu » d'Héraclite est symbole/métaphore du devenir, pas un élément matériel unique comme chez les milésiens. Le « Logos » est loi, raison, discours, mesure.
\item \textbf{Maât :} Concept africain fondateur (Égypte antique) : ordre, justice, vérité, harmonie. À rapprocher des notions de \emph{loi fondamentale}, \emph{constitution}, \emph{contrat social} pour l'application camerounaise.
\item \textbf{Méthode socratique :} Deux temps indissociables : 1. Ironie (destruction de la fausse certitude) $\rightarrow$ 2. Maïeutique (construction de la vraie définition). But : \emph{soin de l'âme}.
\item \textbf{Dissertation :} Respecter la structure imposée (Intro $\rightarrow$ 3 parties dialectiques $\rightarrow$ Conclusion). \textbf{Exemples camerounais obligatoires} dans chaque partie (urbanisation, bilinguisme, chefferies, code de la route, Constitution, etc.).
\item \textbf{Explication de texte :} 1. Thèse + structure. 2. Analyse arguments/concepts. 3. Portée (actualisation camerounaise). Citer le texte (guillemets, numéros de ligne si fournis).
\end{itemize}
\end{aretenir}

\newpage

\coursec{Corrigés des exercices}

\courssub{Atelier méthodologique : Rédiger et justifier une démarche philosophique (Dissertation \& Explication)}

\begin{corrige}[Exercice 1 — QCM : repères de la philosophie antique]
\begin{enumerate}
\item \textbf{Bonne réponse : b.} \cle{Maât} est le concept central de la pensée pharaonique : ordre cosmique, justice, vérité, harmonie. C'est à la fois une déesse et le principe régulateur du monde, dont le pharaon est le garant sur terre.
\item \textbf{Bonne réponse : b.} Pour Héraclite, « ce monde a toujours été, est et sera un feu toujours vivant, s'allumant avec mesure et s'éteignant avec mesure ». Le feu symbolise le \cle{devenir} perpétuel : tout s'écoule (\emph{panta rhei}).
\item \textbf{Bonne réponse : b.} Parménide affirme dans son \emph{Poème} : « L'Être est, le non-être n'est pas. » Le changement, la naissance et la mort sont des illusions des sens (\cle{doxa}) ; seule la raison accède à l'Être un et immuable. La formule du fleuve (a) appartient à Héraclite, et « l'homme mesure de toutes choses » (c) au sophiste Protagoras.
\item \textbf{Bonne réponse : b.} La \cle{maïeutique} (de \emph{maia}, sage-femme — métier de la mère de Socrate) est l'art de faire accoucher les esprits : par le questionnement, Socrate aide son interlocuteur à tirer de lui-même la vérité. Enseigner contre rémunération (d) est le fait des sophistes, que Socrate combat.
\item \textbf{Bonne réponse : b.} L'allégorie de la caverne (\emph{La République}, livre VII) figure le passage de l'opinion (\cle{doxa}, les ombres prises pour la réalité) à la connaissance véritable (\cle{épistémè}, la contemplation des réalités intelligibles, symbolisées par le soleil).
\end{enumerate}
\textbf{Points de vigilance :} ne pas attribuer à Socrate des écrits (il n'a rien écrit) ; ne pas confondre Socrate avec les sophistes ; bien distinguer la thèse d'Héraclite (devenir) de celle de Parménide (immobilité de l'Être).
\end{corrige}

\begin{corrige}[Exercice 2 — Vrai ou faux]
\begin{enumerate}
\item \textbf{Faux.} La pensée pharaonique est précisément centrée sur la justice et l'ordre du monde : le principe de \cle{Maât} règle à la fois le cosmos, la société et la conduite individuelle. Les enseignements de sagesse (par exemple l'\emph{Enseignement de Ptahhotep}) témoignent d'une réflexion éthique sur la justice, la mesure et la parole vraie.
\item \textbf{Faux.} Héraclite et Parménide défendent des thèses radicalement opposées : pour Héraclite, le changement est la loi même du réel (« on ne se baigne jamais deux fois dans le même fleuve ») ; pour Parménide, le changement est une illusion des sens, car l'Être est un, plein et immobile.
\item \textbf{Faux.} Socrate n'a rien écrit : il philosophait oralement, sur l'agora d'Athènes, par le dialogue. Sa pensée nous est connue indirectement, principalement par les dialogues de Platon et les \emph{Mémorables} de Xénophon.
\item \textbf{Faux.} Pour Platon, le monde sensible est le monde du devenir, des copies imparfaites et de l'opinion. La réalité pleine et parfaite appartient au monde intelligible, celui des \cle{Idées}, éternelles et immuables.
\item \textbf{Vrai.} Dans le \emph{Criton} de Platon, Socrate refuse l'évasion que lui proposent ses amis : les lois de la cité l'ont fait naître, éduquer et vivre ; les violer, même pour échapper à une condamnation injuste, serait commettre une injustice et ruiner la cité. Il boit donc la ciguë par respect de la loi.
\end{enumerate}
\end{corrige}

\begin{corrige}[Exercice 3 — Définitions]
\begin{enumerate}
\item \textbf{Le \cle{Logos} (chez Héraclite) :} principe rationnel et universel qui gouverne le devenir. Le changement n'est pas un chaos : il obéit à une loi, une mesure, une raison commune à tous. \emph{Exemple :} l'alternance du jour et de la nuit ou le retour des saisons (saison des pluies, saison sèche au Cameroun) manifeste un ordre régulier au cœur du changement.
\item \textbf{L'\cle{ironie} socratique :} attitude par laquelle Socrate feint l'ignorance et interroge celui qui prétend savoir, afin de révéler les contradictions de son discours et de détruire sa fausse certitude. \emph{Exemple :} Socrate demande à un interlocuteur sûr de lui « Qu'est-ce que le courage ? », puis montre que chaque définition proposée se contredit elle-même.
\item \textbf{La théorie des \cle{Idées} (Platon) :} les Idées (ou Formes) sont des essences parfaites, éternelles et immuables, accessibles seulement à l'intelligence ; les choses sensibles ne sont que des copies imparfaites qui « participent » des Idées. \emph{Exemple :} aucun cercle tracé au tableau n'est parfaitement rond ; seule l'Idée de cercle, pensée par l'esprit, est parfaite.
\item \textbf{Le \cle{devenir} :} processus de transformation continue par lequel les choses passent d'un état à un autre (naissance, croissance, déclin). Chez Héraclite, il est l'essence même du réel ; chez Platon, il caractérise le monde sensible, par opposition à l'Être stable des Idées. \emph{Exemple :} un enfant devient adolescent puis adulte ; une ville de province devient une grande agglomération.
\end{enumerate}
\textbf{Point de vigilance :} une définition attendue à l'examen comporte toujours trois éléments : la nature de la notion, sa fonction chez l'auteur, et un exemple.
\end{corrige}

\begin{corrige}[Exercice 4 — Auteurs et thèses]
\begin{center}
\begin{tabularx}{\linewidth}{|l|X|X|}
\hline
\rowcolor{popblueL} \textbf{Penseur} & \textbf{Thèse fondamentale} & \textbf{Concept clé} \\
\hline
Héraclite & La réalité est un devenir perpétuel (« tout s'écoule »), régi par une loi rationnelle ; l'unité naît de la tension des contraires. & \cle{Devenir} ; \cle{Logos} \\
\hline
Parménide & L'Être est un, immobile et éternel ; le non-être n'est pas. Le changement et la multiplicité sont des illusions des sens. & \cle{Être} ; raison contre sens \\
\hline
Socrate & « Connais-toi toi-même » : la philosophie est examen de soi et soin de l'âme ; nul ne fait le mal volontairement, la vertu est connaissance. & \cle{Ironie} ; \cle{maïeutique} \\
\hline
Platon & Le monde sensible (copies changeantes, opinion) se distingue du monde intelligible (Idées éternelles, science) ; la cité juste doit être dirigée par des philosophes formés. & \cle{Idées} ; allégorie de la caverne \\
\hline
\end{tabularx}
\end{center}
\textbf{Point de vigilance :} au Probatoire, ce type de tableau permet de vérifier la maîtrise des repères ; chaque case doit contenir une formulation complète, pas un simple mot isolé.
\end{corrige}

\begin{corrige}[Exercice 5 — Le fleuve d'Héraclite au quotidien]
\begin{enumerate}
\item \textbf{Sens de la formule.} « On ne se baigne jamais deux fois dans le même fleuve » signifie que la réalité est un flux permanent : l'eau du fleuve s'écoule sans cesse, et le baigneur lui-même change. Rien ne demeure identique à soi : le changement est la loi fondamentale de tout ce qui existe. Héraclite ajoute que ce flux n'est pas désordre : il obéit au \cle{Logos}, raison commune qui règle les transformations.
\item \textbf{Exemples camerounais.}
\begin{itemize}
\item \emph{La ville de Douala :} en quelques décennies, quartiers, populations, activités et paysages se sont profondément transformés (port, quartiers d'affaires, nouveaux ponts sur le Wouri). Celui qui revient après vingt ans ne retrouve pas « la même » ville : c'est pourtant toujours Douala.
\item \emph{Le camfranglais :} cette langue urbaine mélange français, anglais et langues locales ; son vocabulaire se renouvelle sans cesse d'une génération de lycéens à l'autre. La langue parlée hier n'est déjà plus exactement celle d'aujourd'hui.
\end{itemize}
\item \textbf{Le changement signifie-t-il que rien ne demeure ?} Non. Chez Héraclite, le devenir est ordonné par le \cle{Logos} : une loi stable gouverne le flux. De plus, « le chemin en montant et en descendant est un et le même » : l'identité subsiste \emph{à travers} le changement, comme le fleuve reste « ce » fleuve bien que ses eaux changent. Ainsi Douala change sans cesse, mais demeure Douala : l'identité n'est pas l'immobilité, elle est continuité dans la transformation.
\end{enumerate}
\textbf{Point de vigilance :} ne pas réduire Héraclite à un philosophe du chaos ; la mobilisation du \cle{Logos} est exigée par la question 3.
\end{corrige}

\begin{corrige}[Exercice 6 — Maât et l'ordre social]
\begin{enumerate}
\item \textbf{Le principe de Maât.} \cle{Maât} désigne à la fois l'ordre du cosmos (régularité du soleil, crue du Nil), la justice dans les relations humaines et la vérité dans la parole. Le pharaon a pour mission de maintenir cet équilibre : gouverner justement, c'est empêcher le désordre (\emph{isfet}) de submerger le monde. Maât fonde donc l'idée que la société ne tient que par la justice et l'harmonie.
\item \textbf{Comparaison avec la Constitution camerounaise.}
\begin{itemize}
\item \emph{Points communs :} l'une et l'autre sont des normes supérieures qui organisent la vie collective, définissent la justice et obligent aussi bien les gouvernants que les gouvernés. Comme Maât, la Constitution vise l'équilibre et la paix sociale.
\item \emph{Différences :} Maât est d'origine religieuse et cosmique, incarnée par un pharaon divin ; la Constitution est un texte écrit, œuvre humaine, modifiable, fondée sur la volonté du peuple et assortie d'institutions de contrôle (Conseil constitutionnel, tribunaux).
\end{itemize}
\item \textbf{Le règlement du lycée, une exigence comparable ?} Oui, analogiquement. Le règlement intérieur d'un lycée technique (ponctualité, respect des enseignants et des ateliers, interdiction des fraudes) vise, comme Maât, à maintenir un ordre sans lequel la vie commune — ici l'apprentissage — devient impossible. Le respect des règles n'est pas une soumission servile : c'est la condition de l'harmonie du groupe. Mais la comparaison a ses limites : le règlement scolaire est discuté et révisable démocratiquement, ce qui le rapproche de la Constitution plutôt que de l'ordre sacré pharaonique.
\end{enumerate}
\textbf{Point de vigilance :} une comparaison réussie exige à la fois des ressemblances \emph{et} des différences ; répondre seulement par l'un des deux aspects fait perdre des points.
\end{corrige}

\begin{corrige}[Exercice 7 — Dialogue socratique]
\begin{enumerate}
\item \textbf{Exemple de dialogue corrigé.}
\begin{itemize}
\item[] \emph{Moi :} Tu dis que le courageux n'a jamais peur. Dis-moi : le soldat qui avance au combat en tremblant de peur, mais qui avance quand même, est-il lâche ?
\item[] \emph{Camarade :} Non\dots il combat malgré sa peur.
\item[] \emph{Moi :} Donc on peut avoir peur \emph{et} être courageux ?
\item[] \emph{Camarade :} Il semblerait.
\item[] \emph{Moi :} Et celui qui ne craint rien parce qu'il est ivre, ou parce qu'il ignore le danger — comme un enfant qui joue au bord d'un précipice — est-il courageux ?
\item[] \emph{Camarade :} Non, il est inconscient.
\item[] \emph{Moi :} L'absence de peur ne suffit donc pas à faire le courage. Que faut-il de plus ?
\item[] \emph{Camarade :} Peut-être\dots savoir ce qui mérite d'être affronté ?
\item[] \emph{Moi :} Bien. Et celui qui affronte le danger pour voler ou pour se vanter, est-il courageux ?
\item[] \emph{Camarade :} Non, il faut une cause juste.
\item[] \emph{Moi :} Alors, qu'est-ce que le courage ?
\item[] \emph{Camarade :} C'est savoir ce qu'il faut craindre et ce qu'il ne faut pas craindre, et affronter le danger malgré la peur, pour une cause juste.
\end{itemize}
Les quatre premières répliques relèvent de l'\cle{ironie} (destruction de la fausse définition par des contre-exemples) ; les suivantes relèvent de la \cle{maïeutique} (l'interlocuteur « accouche » lui-même de la définition).
\item \textbf{Définition aboutie :} le courage est la connaissance de ce qu'il faut craindre et de ce qu'il ne faut pas craindre, jointe à la volonté d'affronter le danger pour une fin juste. Le courage n'est donc ni l'absence de peur (témérité, inconscience) ni la fuite (lâcheté) : il suppose un savoir — la vertu est connaissance, selon Socrate.
\item \textbf{Lien avec le souci de l'âme.} Cette méthode illustre le « connais-toi toi-même » : en examinant ses propres opinions, l'interlocuteur découvre qu'il croyait savoir ce qu'il ignorait. Cette prise de conscience purifie l'âme de la fausse science et la rend meilleure : philosopher, pour Socrate, c'est « s'occuper de son âme pour la rendre meilleure », et non accumuler des savoirs extérieurs.
\end{enumerate}
\textbf{Point de vigilance :} dans la rédaction du dialogue, les répliques doivent être courtes, sous forme de questions ; le maître ne donne jamais la réponse, il la fait naître.
\end{corrige}

\begin{corrige}[Exercice 8 — Sortir de la caverne]
\begin{enumerate}
\item \textbf{Les étapes de la libération.} Des prisonniers, enchaînés depuis l'enfance au fond d'une caverne, ne voient que des ombres projetées sur la paroi et les prennent pour la réalité. L'un d'eux est détaché : d'abord ébloui et souffrant, il découvre les objets qui produisaient les ombres, puis il monte vers la sortie, s'habitue progressivement à la lumière, et contemple enfin le monde réel et le soleil, symbole du Bien et de la vérité.
\item \textbf{Situations actuelles comparables.}
\begin{itemize}
\item \emph{Les rumeurs sur les réseaux sociaux :} l'élève qui croit tout ce qui circule sur WhatsApp ou Facebook sans vérifier ressemble au prisonnier qui prend les ombres pour la réalité ; vérifier les sources, comparer, réfléchir, c'est commencer à se détacher.
\item \emph{Les préjugés sur les filières techniques :} croire que « l'enseignement technique est pour les élèves qui ont échoué » est une ombre ; découvrir la valeur réelle des métiers techniques (électricité, mécanique, informatique) et leurs débouchés au Cameroun, c'est sortir de la caverne.
\end{itemize}
\item \textbf{Le retour dans la caverne.} Selon Platon, le philosophe libéré doit redescendre auprès des prisonniers pour les délivrer, même s'ils se moquent de lui ou le menacent : le savoir n'est pas un privilège égoïste, il oblige à servir la cité. \emph{Leçon pour aujourd'hui :} les personnes éduquées — ingénieurs, enseignants, techniciens formés — ont le devoir de mettre leurs compétences au service de la communauté (développement du quartier ou du village, lutte contre les rumeurs et l'ignorance), au lieu de chercher seulement leur avantage personnel.
\end{enumerate}
\textbf{Point de vigilance :} bien distinguer les deux mouvements de l'allégorie : l'ascension (éducation, passage de la \cle{doxa} à l'\cle{épistémè}) et la descente (responsabilité politique du philosophe).
\end{corrige}

\begin{corrige}[Exercice 9 — Dissertation : le changement est-il une perte d'identité ?]
\textbf{Corrigé détaillé (plan et contenus attendus).}

\medskip
\textbf{Introduction (attendus).} \emph{Amorce :} partir de la situation camerounaise — le jeune qui quitte son village pour Douala ou Yaoundé, parle davantage français ou anglais que sa langue maternelle, adopte de nouveaux modes de vie : sa famille s'inquiète de le voir « se perdre ». \emph{Définitions :} le changement est le passage d'un état à un autre ; l'identité est ce qui fait qu'un être demeure le même à travers le temps. \emph{Problème :} si tout change, comment quelque chose peut-il rester « soi » ? \emph{Problématique :} le changement détruit-il l'identité, ou bien l'identité se construit-elle à travers le changement ? \emph{Annonce du plan.}

\medskip
\textbf{Première partie — Héraclite : le devenir est la loi du réel.} Thèse : « tout s'écoule », « on ne se baigne jamais deux fois dans le même fleuve » ; le feu symbolise la transformation permanente ; l'identité figée est une illusion. Conséquence : vouloir empêcher tout changement, c'est lutter contre la nature même des choses. \emph{Exemples :} l'urbanisation de Yaoundé et Douala ; l'évolution des langues (camfranglais) ; les modes de vie des jeunes générations. Aucune culture camerounaise n'est restée identique à elle-même depuis un siècle — et pourtant elle existe toujours.

\medskip
\textbf{Deuxième partie — Parménide : ce qui est véritablement ne change pas.} Thèse : l'Être est, le non-être n'est pas ; le changement suppose qu'une chose cesse d'être ce qu'elle est, ce qui est impensable ; le devenir n'est qu'apparence sensible, la raison saisit l'identité permanente. Conséquence : si l'identité est réelle, elle doit reposer sur un fond stable ; un changement total abolirait l'identité — un jeune qui renierait entièrement sa langue, son nom, ses valeurs aurait bien « perdu » quelque chose. \emph{Exemples :} la permanence des valeurs transmises (respect des aînés, solidarité familiale), des langues et des institutions traditionnelles (chefferies) à travers les transformations.

\medskip
\textbf{Troisième partie — Dépassement platonicien.} Platon distingue le monde sensible (changeant, imparfait) et le monde intelligible (Idées stables) : le changement et la permanence ne s'opposent pas absolument, ils n'appartiennent pas au même niveau. Appliqué à notre problème : les \emph{apparences} d'une culture ou d'une personne changent (habits, langue d'usage, lieu de vie), mais son \emph{fond} — valeurs, mémoire, appartenance — peut demeurer, à condition d'être assumé et transmis. L'identité véritable n'est donc ni dans le changement pur (dispersion) ni dans l'immobilité (musée figé) : elle est continuité d'un fond à travers la transformation des formes. \emph{Exemple :} le jeune urbain peut parler français au travail et sa langue maternelle en famille, vivre en ville et rester attaché à sa chefferie : il transforme son identité sans la perdre.

\medskip
\textbf{Conclusion.} Le changement n'est une perte d'identité que s'il atteint le fond et rompt la continuité ; vécu comme transformation assumée, il est au contraire la manière dont l'identité se maintient vivante. \emph{Ouverture :} quelle politique culturelle (enseignement des langues nationales, valorisation des traditions) permettrait aux jeunes Camerounais de changer sans se perdre ?

\medskip
\textbf{Barème indicatif (sur 20) :} introduction (amorce, définitions, problématique, plan) : 4 pts ; partie 1 (Héraclite exact et illustré) : 4 pts ; partie 2 (Parménide exact et illustré) : 4 pts ; partie 3 (dépassement argumenté, Platon correctement mobilisé) : 4 pts ; conclusion et ouverture : 2 pts ; langue, organisation, exemples camerounais précis : 2 pts.

\medskip
\textbf{Points de vigilance :} ne pas opposer Héraclite et Parménide comme deux simples « opinions » mais comme deux thèses argumentées ; la troisième partie doit être un véritable dépassement, pas un résumé ; citer au moins un exemple camerounais par partie ; interdiction de recopier le contexte comme amorce sans le transformer.
\end{corrige}

\begin{corrige}[Exercice 10 — Dissertation : obéir aux lois, est-ce renoncer à sa liberté ?]
\textbf{Corrigé détaillé (plan et contenus attendus).}

\medskip
\textbf{Introduction.} \emph{Amorce :} le sentiment de certains jeunes — « être libre, c'est faire ce qu'on veut » — qui voient dans toute loi (code de la route, règlement scolaire, impôts) une entrave. \emph{Définitions :} la loi est une règle générale et obligatoire édictée par la cité ; la liberté est le pouvoir d'agir selon sa volonté. \emph{Problème :} si la loi limite la volonté individuelle, obéir semble renoncer à être libre ; mais sans lois, la liberté de chacun est menacée par celle des autres. \emph{Problématique :} la loi est-elle l'ennemie de la liberté, ou sa condition ? \emph{Annonce du plan.}

\medskip
\textbf{Première partie — La loi comme limitation de la liberté.} La loi interdit, oblige, sanctionne : elle contraint la volonté individuelle (on ne peut pas rouler à la vitesse qu'on veut, ni s'installer où l'on veut). Pour qui définit la liberté comme absence de contrainte, toute obéissance est une aliénation. \emph{Exemple :} le jeune motard qui voit dans le port du casque et le contrôle de gendarmerie une atteinte à sa liberté. Cette conception n'est pas absurde : une loi injuste ou arbitraire peut effectivement opprimer.

\medskip
\textbf{Deuxième partie — Socrate dans le \emph{Criton} : pourquoi obéir même à une loi injustement appliquée.} Condamné à mort injustement, Socrate refuse de s'évader. Son argument : les lois de la cité sont comme des parents — elles l'ont fait naître, éduquer, protéger ; il a vécu toute sa vie sous elles, acceptant tacitement leur contrat ; les violer pour sauver sa vie serait commettre une injustice en réponse à une injustice, et contribuer à détruire la cité. Sans lois respectées, il n'y a plus de cité, donc plus de vie commune possible. La désobéissance individuelle, généralisée, détruit la condition de la liberté de tous.

\medskip
\textbf{Troisième partie — Dépassement : la loi juste comme condition de la liberté.} \emph{Platon (la caverne) :} la vraie liberté n'est pas de faire ce qu'on veut — le prisonnier qui suit ses désirs et ses opinions est un esclave qui s'ignore — mais de se libérer intérieurement par la connaissance ; la loi juste, éclairée par la raison, aide à cette libération. \emph{Maât :} la justice est la condition de l'harmonie collective ; sans ordre juste, c'est le désordre (\emph{isfet}) qui règne, et avec lui l'arbitraire du plus fort — la pire des servitudes. La loi juste ne supprime donc pas la liberté : elle la rend possible pour tous, en limitant la « liberté » nuisible de chacun. Distinction décisive : liberté (agir selon la raison dans un ordre juste) contre licence (faire ce qu'on veut, qui aboutit au règne du plus fort).

\medskip
\textbf{Conclusion.} Obéir aux lois justes n'est pas renoncer à sa liberté, c'est en garantir l'exercice pour soi et pour autrui ; seules les lois injustes posent le problème de la désobéissance. \emph{Exemple civique camerounais :} respecter le code de la route, c'est renoncer à la licence de rouler n'importe comment pour gagner, avec tous les autres usagers, la liberté réelle de circuler en sécurité. \emph{Ouverture :} que faire lorsque la loi elle-même est injuste — obéir, réformer, désobéir ?

\medskip
\textbf{Barème indicatif (sur 20) :} introduction et problématique : 4 pts ; partie 1 : 3 pts ; partie 2 (argument du \emph{Criton} correctement analysé) : 4 pts ; partie 3 (Platon + Maât, distinction liberté/licence) : 5 pts ; conclusion avec prise de position et exemple civique : 2 pts ; langue et organisation : 2 pts.

\medskip
\textbf{Points de vigilance :} le cœur du sujet est la distinction entre \emph{liberté} et \emph{licence} ; l'argument de Socrate doit être reconstruit (les lois comme parents et comme contrat), pas seulement raconté ; la troisième partie doit mobiliser \emph{à la fois} Platon et Maât, comme l'exige la consigne.
\end{corrige}

\begin{corrige}[Exercice 11 — Explication de texte : l'\emph{Apologie de Socrate}]
\begin{enumerate}
\item \textbf{Thèse du texte.} Socrate s'attribue une mission philosophique : exhorter chaque citoyen à s'occuper de son âme, de la vérité et de la pensée plutôt que des richesses et des honneurs. Il refuse d'y renoncer, même au prix de sa vie, parce que cette mission lui vient du dieu (l'oracle de Delphes) et qu'elle constitue le sens de son existence : cesser de philosopher serait trahir à la fois le dieu, la cité et lui-même. Mieux vaut mourir fidèle à sa mission que vivre en renonçant à la recherche de la vérité.
\item \textbf{Analyse des arguments.} « J'obéirai au dieu plutôt qu'à vous » : Socrate établit une hiérarchie des valeurs. Au sommet se trouvent l'âme, la vérité, la justice — valeurs absolues, d'ordre divin ou rationnel ; en dessous, les ordres des hommes, les richesses, la réputation, la vie elle-même — valeurs relatives. Argument implicite : les juges peuvent disposer de son corps, non de sa conscience ; la vie ne vaut d'être vécue que si elle est examinée (« une vie sans examen ne vaut pas la peine d'être vécue », dit-il ailleurs dans l'\emph{Apologie}). Obéir aux hommes contre sa mission serait préférer l'existence à la valeur qui la justifie.
\item \textbf{Les concepts dans le texte.}
\begin{itemize}
\item \emph{L'ironie} : « ô toi le meilleur des hommes, citoyen d'Athènes, la plus grande et la plus renommée des cités » — l'éloge apparent retourne en reproche : plus la cité est glorieuse, plus l'oubli de l'âme est honteux. L'ironie démasque l'écart entre ce que l'Athénien prétend être et ce qu'il est.
\item \emph{Le soin de l'âme} : « t'occuper de ta pensée, de la vérité et de ton âme, pour la rendre meilleure » — la philosophie n'est pas un savoir théorique mais une pratique de transformation de soi, un souci éthique permanent.
\item \emph{Le rapport à la cité} : Socrate aime Athènes (« je vous respecte et je vous aime ») précisément en la critiquant : le philosophe n'est ni un ennemi de la cité ni un flatteur, mais son examinateur. Il sert la cité en lui disant ses vérités.
\end{itemize}
\item \textbf{La comparaison du taon.} Socrate se compare à un taon (mouche à bœuf) piquant un cheval grand et noble mais paresseux : Athènes, riche et puissante, s'est alourdie dans le confort, l'argent et les certitudes ; le philosophe la pique, la réveille, l'empêche de s'endormir dans l'opinion. Fonction du philosophe : être le réveilleur et l'examinateur critique de la cité — rôle gênant, qui explique l'hostilité dont il est l'objet, mais indispensable à la santé morale de la communauté.
\item \textbf{Portée actuelle au Cameroun (deux paragraphes attendus).} \emph{Paragraphe 1 :} des figures jouent aujourd'hui un rôle comparable de « taon » : les journalistes d'investigation qui dénoncent la corruption et les détournements, les enseignants qui forment à l'esprit critique, les artistes (chanteurs, slameurs, romanciers) qui dénoncent l'injustice, parfois les chefs traditionnels ou religieux qui rappellent les gouvernants et les citoyens à leurs devoirs. Comme Socrate, ils dérangent les certitudes installées et rappellent que la richesse et le pouvoir ne valent rien sans justice ni vérité. \emph{Paragraphe 2 :} cependant, ce rôle comporte les mêmes risques que celui de Socrate : celui qui « pique » la cité s'expose à l'hostilité, aux poursuites, à la marginalisation. Encore faut-il distinguer le véritable examinateur — qui critique par amour de la cité et au nom de valeurs argumentées — du simple polémiste ou du démagogue qui flatte les passions. La leçon de Socrate est double : une société a besoin de ses « taons » pour ne pas s'endormir, et chacun de nous doit commencer par s'examiner soi-même, à l'école, en famille, dans la vie civique.
\end{enumerate}
\textbf{Barème indicatif (sur 20) :} question 1 (thèse) : 4 pts ; question 2 (arguments, hiérarchie des valeurs) : 4 pts ; question 3 (trois concepts analysés dans le texte) : 4 pts ; question 4 (fonction du taon) : 3 pts ; question 5 (actualisation argumentée en deux paragraphes) : 3 pts ; langue et rigueur : 2 pts.

\medskip
\textbf{Points de vigilance :} en explication de texte, chaque concept doit être défini \emph{à partir du texte} (citation à l'appui), non récité de mémoire ; l'actualisation (question 5) doit rester une discussion argumentée, non une simple liste de noms ; respecter la consigne des deux paragraphes.
\end{corrige}

\newpage

\coursec{Fiche bilan}

\begin{popfigure}
\centering
\begin{center}\tcbox[colback=popink,colframe=popink,coltext=white,arc=9pt,boxsep=4pt,left=6pt,right=6pt]{\bfseries\small Philosophie Antique}\end{center}
\vspace{-2pt}
\begin{tcbraster}[raster columns=3,raster equal height=rows,raster column skip=6pt,raster row skip=6pt,enhanced,blankest]
\begin{tcolorbox}[enhanced,colback=white,colframe=popblue,boxrule=0.9pt,arc=3pt,title={1. Pharaonique (Maât)},fonttitle=\bfseries\scriptsize,coltitle=white,colbacktitle=popblue,left=4pt,right=4pt,top=2pt,bottom=2pt,boxsep=1pt,toptitle=1.5pt,bottomtitle=1.5pt,halign=left]
\begin{itemize}[nosep,leftmargin=*,itemsep=1pt,font=\scriptsize]
\item Ordre cosmique
\item Justice/Équilibre
\end{itemize}
\end{tcolorbox}
\begin{tcolorbox}[enhanced,colback=white,colframe=popgreen,boxrule=0.9pt,arc=3pt,title={2. Présocratiques},fonttitle=\bfseries\scriptsize,coltitle=white,colbacktitle=popgreen,left=4pt,right=4pt,top=2pt,bottom=2pt,boxsep=1pt,toptitle=1.5pt,bottomtitle=1.5pt,halign=left]
\begin{itemize}[nosep,leftmargin=*,itemsep=1pt,font=\scriptsize]
\item Héraclite : Flux
\item Parménide : Être
\end{itemize}
\end{tcolorbox}
\begin{tcolorbox}[enhanced,colback=white,colframe=poporange,boxrule=0.9pt,arc=3pt,title={3. Socrate},fonttitle=\bfseries\scriptsize,coltitle=white,colbacktitle=poporange,left=4pt,right=4pt,top=2pt,bottom=2pt,boxsep=1pt,toptitle=1.5pt,bottomtitle=1.5pt,halign=left]
\begin{itemize}[nosep,leftmargin=*,itemsep=1pt,font=\scriptsize]
\item Maïeutique
\item Connais-toi toi-même
\item Éthique rationnelle
\end{itemize}
\end{tcolorbox}
\begin{tcolorbox}[enhanced,colback=white,colframe=poppurple,boxrule=0.9pt,arc=3pt,title={4. Platon},fonttitle=\bfseries\scriptsize,coltitle=white,colbacktitle=poppurple,left=4pt,right=4pt,top=2pt,bottom=2pt,boxsep=1pt,toptitle=1.5pt,bottomtitle=1.5pt,halign=left]
\begin{itemize}[nosep,leftmargin=*,itemsep=1pt,font=\scriptsize]
\item Monde des Idées
\item Allégorie caverne
\item Politique/Éducation
\end{itemize}
\end{tcolorbox}
\begin{tcolorbox}[enhanced,colback=white,colframe=poppink,boxrule=0.9pt,arc=3pt,title={5. Problématique Unité/Multiple},fonttitle=\bfseries\scriptsize,coltitle=white,colbacktitle=poppink,left=4pt,right=4pt,top=2pt,bottom=2pt,boxsep=1pt,toptitle=1.5pt,bottomtitle=1.5pt,halign=left]
\begin{itemize}[nosep,leftmargin=*,itemsep=1pt,font=\scriptsize]
\item Un principe ?
\item Diversité réelle ?
\end{itemize}
\end{tcolorbox}
\begin{tcolorbox}[enhanced,colback=white,colframe=popgold,boxrule=0.9pt,arc=3pt,title={6. Méthode},fonttitle=\bfseries\scriptsize,coltitle=white,colbacktitle=popgold,left=4pt,right=4pt,top=2pt,bottom=2pt,boxsep=1pt,toptitle=1.5pt,bottomtitle=1.5pt,halign=left]
\begin{itemize}[nosep,leftmargin=*,itemsep=1pt,font=\scriptsize]
\item Dissertation
\item Explication texte
\end{itemize}
\end{tcolorbox}
\end{tcbraster}

\legende{Carte mentale de la leçon 16 : La philosophie antique (6 axes majeurs).}
\end{popfigure}

\begin{bilan}
\courssub{Définitions clés}
\begin{itemize}
\item \cle{Maât} : Ordre cosmique, justice, vérité, harmonie fondatrice de l'Égypte pharaonique ; principe d'unité liant le divin, le royal et le social.
\item \cle{Arché} : Principe premier, origine unique de toutes choses (recherché par les Présocratiques : feu, eau, air, être).
\item \cle{Devenir (Héraclite)} : Réalité comme flux perpétuel (\emph{Panta rhei}) ; unité par les contraires (voie montante/descendante) et le \cle{Logos} (loi rationnelle universelle).
\item \cle{Être (Parménide)} : L'Être est, le non-être n'est pas ; l'Être est un, immuable, éternel, sphérique, pensable par la raison (\cle{noûs}) seule, non par les sens.
\item \cle{Maïeutique (Socrate)} : Art d'accoucher les esprits par le dialogue (ironie $\rightarrow$ aporie $\rightarrow$ recherche commune) ; la vertu est science, nul n'est méchant volontairement.
\item \cle{Théorie des Idées (Platon)} : Deux mondes : sensible (devenir, opinion/\emph{doxa}) et intelligible (Être, Idées/\emph{eidos}, science/\emph{epistémè}). L'Idée du Bien = soleil du intelligible.
\item \cle{Allégorie de la Caverne} : Image de la condition humaine (ombres $\rightarrow$ statues $\rightarrow$ astres $\rightarrow$ Soleil) et de la mission du philosophe-roi (redescendre pour gouverner).
\end{itemize}

\courssub{Repères chronologiques \& Auteurs (Tabularx)}
\begin{center}
\begin{tabularx}{\linewidth}{|l|X|l|}\hline
\textbf{Période / Aire} & \textbf{Concepts / Thèses centrales} & \textbf{Figures} \\ \hline
Pharaonique (Égypte) & Maât, ordre cosmico-social, royauté fonction & Ptahhotep, Mérenrê \\ \hline
Présocratiques (Grèce, VIe–Ve s.) & Recherche de l'\cle{arché}, passage mythe $\rightarrow$ logos & \textbf{Héraclite} (Éphèse), \textbf{Parménide} (Élée) \\ \hline
Tournant anthropologique & Éthique, définition, dialogue, soin de l'âme & \textbf{Socrate} (Athènes, 470–399) \\ \hline
Systématisation platonicienne & Idées, dialectique, politique, éducation, âme triple & \textbf{Platon} (Athènes, 428–348) \\ \hline
\end{tabularx}
\end{center}

\courssub{Méthodes express (Bac)}
\begin{enumerate}
\item \textbf{Dissertation} : Analyser le sujet (termes, présupposés) $\rightarrow$ Problématiser (tension Unité/Multiple, Sens/Non-sens, Apparaître/Être) $\rightarrow$ Plan dialectique (Thèse : ex. Unité principe ; Antithèse : Multiple irréductible ; Synthèse : Unité dialectique du Multiple) $\rightarrow$ Intégrer auteurs (Socrate/Platon vs Héraclite/Parménide) et contexte camerounais.
\item \textbf{Explication de texte} : 1. Présentation (auteur, œuvre, contexte, thèse). 2. Analyse structurée (découpage, arguments, concepts clés, références implicites). 3. Discussion critique (portée, limites, actualité). 4. Conclusion (synthèse + ouverture).
\item \textbf{Application camerounaise} : Identifier une situation concrète (chefferie, justice coutumière, dialogue intergénérationnel, éducation bilingue) $\rightarrow$ Mobiliser une thèse antique (ex. Maât pour la médiation, Maïeutique pour le dialogue scolaire, Caverne pour l'émancipation par l'instruction technique).
\end{enumerate}
\end{bilan}

\begin{aretenir}[Checklist avant le Bac]
\begin{itemize}
\item $\square$ Je définis \cle{Maât} et j'explique son rôle unificateur (cosmos/société) en Égypte antique.
\item $\square$ Je distingue \cle{Héraclite} (devenir, feu, logos, unité des contraires) et \cle{Parménide} (être immuable, raison vs sens, critique du devenir).
\item $\square$ Je caractérise la \cle{maïeutique socratique} (ironie, aporie, accouchement des esprits, intellectualisme moral).
\item $\square$ J'expose la \cle{Théorie des Idées} (dualisme mondes, Idée du Bien, participation, réminiscence).
\item $\square$ J'analyse l'\cle{Allégorie de la Caverne} (étapes de la connaissance, rôle politique du philosophe).
\item $\square$ Je problématise le couple \cle{Unité / Multiple} à travers les réponses antiques (arché, Être, Idées, Maât).
\item $\square$ Je sais structurer une \textbf{dissertation} (intro : accroche, définition, problème, plan ; développement : 3 parties argumentées ; conclusion : bilan, ouverture).
\item $\square$ Je sais structurer une \textbf{explication de texte} (présentation, analyse suivie, discussion, conclusion).
\item $\square$ Je donne un \textbf{exemple camerounais} pertinent pour illustrer une thèse (ex. Maât et palabre, Socrate et éducation technique, Caverne et formation professionnelle).
\item $\square$ J'évite le \textbf{hors-sujet} : je ne raconte pas la vie des auteurs, j'utilise leurs concepts pour répondre à la question.
\item $\square$ Je soigne la \textbf{langue} (philosophique, précise, connecteurs logiques : \emph{donc, or, toutefois, en effet, par conséquent}).
\end{itemize}
\end{aretenir}

\findecours

\end{document}
